% Maintenance de premier niveau — Informatique 3ème — Cours signé OnBuch+
% Programme officiel MINESEC. Compiler avec : tectonic N03-maintenance-premier-niveau.tex
\newcommand{\DOCMATIERE}{Informatique}
\newcommand{\DOCNIVEAU}{3ème}
\newcommand{\DOCMODULE}{Informatique – classe de 3ème}
\newcommand{\DOCLECON}{N03}
\newcommand{\DOCTITRE}{Maintenance de premier niveau}
% OnBuch+ — préambule « cours signés » ENRICHI (compilé avec tectonic / XeLaTeX)
% Système visuel « Pop » d'OnBuch+ : fond crème (jamais blanc), bordures encre
% épaisses, ombres « dures » décalées, titres Archivo Black en capitales.
% Version MATHÉMATIQUES : graphes (pgfplots/tikz), boîtes pédagogiques
% (définition, propriété, méthode, attention, le savais-tu, exercice résolu,
% exercice, TP, bilan), page de garde et signature OnBuch+ ancrée en bas.
%
% Macros attendues dans main.tex AVANT \input{preamble} :
%   \DOCMATIERE  \DOCNIVEAU  \DOCMODULE  \DOCLECON (numéro)  \DOCTITRE
\documentclass[11pt]{article}

\usepackage{fontspec}
\usepackage[french]{babel}
\usepackage[a4paper,margin=2cm,top=3.4cm,bottom=2.8cm,headheight=1.4cm,footskip=1.3cm]{geometry}
\usepackage{amsmath,amssymb,amsfonts}
\usepackage{mathtools} % \xmapsto, \coloneqq... souvent utilisés par les modèles
\usepackage{cancel} % simplifications barrées (bilans de forces)
% \abs{z} (module, valeur absolue) et \norm{u} (norme), utilisés par les modèles
\providecommand{\abs}{}\let\abs\relax\DeclarePairedDelimiter{\abs}{\lvert}{\rvert}
\providecommand{\norm}{}\let\norm\relax\DeclarePairedDelimiter{\norm}{\lVert}{\rVert}
\usepackage[table]{xcolor} % [table] : fournit aussi \rowcolors (alternance de lignes)
\usepackage{pagecolor}
\usepackage{tikz}
\usetikzlibrary{arrows.meta,positioning,calc,shapes.geometric,shapes.misc,shapes.arrows,shapes.symbols,decorations.pathmorphing,decorations.pathreplacing,decorations.markings,patterns,shadows,mindmap,trees,fit,backgrounds,matrix,intersections,angles,quotes}
\usepackage{pgfplots}
\usepackage[version=4]{mhchem} % équations nucléaires \ce{^{235}_{92}U}
\usepackage[european,siunitx]{circuitikz} % schémas électriques
\pgfplotsset{compat=1.18}
\usepgfplotslibrary{fillbetween,groupplots} % aires entre courbes (calcul intégral)
\usepackage{enumitem}
\usepackage{fancyhdr}
% [most] charge toutes les bibliothèques tcolorbox (listings, minted, raster,
% documentation...) et gonfle inutilement la mémoire TeX sur un document long
% (~300 boîtes) : on ne charge que ce qui est réellement utilisé.
\usepackage[skins,breakable]{tcolorbox}
\usepackage{graphicx}
\usepackage{colortbl}
\usepackage{booktabs}
\usepackage{tabularx}
\usepackage{diagbox} % case d'en-tête diagonale des tableaux à double entrée
% Colonnes extensibles centrée / gauche / droite (C, L, R), souvent utilisées par les modèles
\newcolumntype{C}{>{\centering\arraybackslash}X}
\newcolumntype{L}{>{\raggedright\arraybackslash}X}
\newcolumntype{R}{>{\raggedleft\arraybackslash}X}
\usepackage{array}
\usepackage{multirow}
\usepackage{multicol}
\usepackage{siunitx}
% parse-numbers=false : accepte aussi \SI{2,5 \times 10^{-3}}{...}
\sisetup{locale=FR,output-decimal-marker={,},group-minimum-digits=4,parse-numbers=false}
\DeclareSIUnit\ohm{\ensuremath{\symup{\Omega}}} % U+2126 absent de la police
\DeclareSIUnit\division{div} % réglages d'oscilloscope (ms/div, V/div)
\DeclareSIUnit\tr{tr} % tours (vitesse de rotation, ex. \SI{1500}{\tr\per\minute})
\DeclareSIUnit\bit{bit}
\DeclareSIUnit\byte{o} % octet
\DeclareSIUnit\inch{po} % pouce
\DeclareSIUnit\ppp{ppp} % points par pouce (résolution d'image)
\usepackage{qrcode}
% Blocs de code source (PHP, C, HTML, JavaScript, SQL) — spécifique à la
% matière Informatique : listings + habillage aux couleurs OnBuch+ (voir le
% style \lstset ci-dessous, à côté des autres boîtes pédagogiques).
\usepackage{listings}
% \degree : commande fréquemment utilisée par les modèles pour un angle ou une
% température en dehors de \SI{}{} (non fournie par les paquets chargés).
\providecommand{\degree}{^{\circ}}
% \qed (fin de démonstration), utilisé par les modèles sans amsthm
\providecommand{\qed}{\hfill$\square$}
% \barre : double barre des valeurs interdites dans un tableau de variations (tkz-tab)
\providecommand{\barre}{\ensuremath{\|}}
% environnement proof (amsthm non chargé) utilisé par les modèles
\ifdefined\proof\else\newenvironment{proof}[1][Démonstration]{\par\noindent\textit{#1.}\ }{\qed\par}\fi
\usepackage{lastpage}
\usepackage{hyperref}
\hypersetup{hidelinks}

% ============================================================
% Palette « Pop » OnBuch+ — mêmes hex que la classe P (plus_ui.dart)
% ============================================================
\definecolor{poporange}{HTML}{F59321}
\definecolor{popdark}{HTML}{E07A0C}
\definecolor{popcream}{HTML}{FFF6E8}
\definecolor{popink}{HTML}{1C1714}
\definecolor{poppurple}{HTML}{7A5AE0}
\definecolor{popgreen}{HTML}{2E9E6A}
\definecolor{poppink}{HTML}{E0457B}
\definecolor{popblue}{HTML}{2F7DE1}
\definecolor{popgold}{HTML}{C98A12}
\definecolor{popmuted}{HTML}{8A7F76}
\definecolor{popline}{HTML}{EADFD2}
\definecolor{popgray}{HTML}{8A7F76}
\colorlet{popgrey}{popgray}
% Alias tolérants (le modèle nomme parfois les couleurs différemment)
\colorlet{popwhite}{white}
\colorlet{popblack}{popink}
\colorlet{popred}{poppink}
\colorlet{popyellow}{popgold}
% Teintes claires pour fonds de tableaux / zones
\colorlet{popblueL}{popblue!10!white}
\colorlet{poporangeL}{poporange!14!white}
\colorlet{popgreenL}{popgreen!12!white}
\colorlet{poppurpleL}{poppurple!12!white}
\colorlet{poppinkL}{poppink!12!white}
\colorlet{popgoldL}{popgold!14!white}

\pagecolor{popcream}

% ============================================================
% Typographie
% ============================================================
\defaultfontfeatures{Ligatures=TeX}
\setmainfont{PlusJakartaSans-Medium.ttf}[Path=../../../fonts/,BoldFont=PlusJakartaSans-Bold.ttf]
% Maths en sans-serif (Fira Math) avec chiffres et lettres droites de Plus
% Jakarta, pour que les formules se fondent dans le texte.
\usepackage[math-style=ISO,bold-style=ISO]{unicode-math}
\setmathfont{FiraMath-Regular.otf}[Path=../../../fonts/]
\setmathfont{PlusJakartaSans-Medium.ttf}[Path=../../../fonts/,range={up/num,up/{Latin,latin},\mathup}]
\usepackage{icomma} % virgule décimale sans espace en mode math
% Caractères mathématiques Unicode absents des polices de texte (Plus Jakarta,
% Archivo Black) : rendus par leur équivalent mathématique, en texte comme en
% formule — sinon ils s'affichent en carré vide (ex. « Arithmétique dans ℤ »).
\usepackage{newunicodechar}
\newunicodechar{ℝ}{\ensuremath{\mathbb{R}}}
\newunicodechar{ℤ}{\ensuremath{\mathbb{Z}}}
\newunicodechar{ℂ}{\ensuremath{\mathbb{C}}}
\newunicodechar{ℕ}{\ensuremath{\mathbb{N}}}
\newunicodechar{ℚ}{\ensuremath{\mathbb{Q}}}
\newunicodechar{𝔻}{\ensuremath{\mathbb{D}}}
\newunicodechar{α}{\ensuremath{\alpha}}
\newunicodechar{β}{\ensuremath{\beta}}
\newunicodechar{γ}{\ensuremath{\gamma}}
\newunicodechar{δ}{\ensuremath{\delta}}
\newunicodechar{ε}{\ensuremath{\varepsilon}}
\newunicodechar{θ}{\ensuremath{\theta}}
\newunicodechar{λ}{\ensuremath{\lambda}}
\newunicodechar{μ}{\ensuremath{\mu}}
\newunicodechar{σ}{\ensuremath{\sigma}}
\newunicodechar{τ}{\ensuremath{\tau}}
\newunicodechar{φ}{\ensuremath{\varphi}}
\newunicodechar{ψ}{\ensuremath{\psi}}
\newunicodechar{ω}{\ensuremath{\omega}}
\newunicodechar{Σ}{\ensuremath{\Sigma}}
% majuscules produites par \MakeUppercase dans les titres (α-aminés -> Α)
\newunicodechar{Α}{\ensuremath{\alpha}}
\newunicodechar{Β}{\ensuremath{\beta}}
\newunicodechar{Γ}{\ensuremath{\Gamma}}
\newunicodechar{Δ}{\ensuremath{\Delta}}
\newunicodechar{Ε}{\ensuremath{\varepsilon}}
\newunicodechar{Θ}{\ensuremath{\Theta}}
\newunicodechar{Λ}{\ensuremath{\Lambda}}
\newunicodechar{Μ}{\ensuremath{\mu}}
\newunicodechar{Τ}{\ensuremath{\tau}}
\newunicodechar{Φ}{\ensuremath{\Phi}}
\newunicodechar{Ψ}{\ensuremath{\Psi}}
\newunicodechar{Ω}{\ensuremath{\Omega}}
\newunicodechar{↦}{\ensuremath{\mapsto}}
\newunicodechar{∈}{\ensuremath{\in}}
\newunicodechar{∉}{\ensuremath{\notin}}
\newunicodechar{∧}{\ensuremath{\wedge}}
\newunicodechar{⇔}{\ensuremath{\Leftrightarrow}}
\newunicodechar{⟺}{\ensuremath{\Longleftrightarrow}}
\newunicodechar{⇒}{\ensuremath{\Rightarrow}}
\newunicodechar{⟹}{\ensuremath{\Longrightarrow}}
\newunicodechar{∘}{\ensuremath{\circ}}
\newunicodechar{∪}{\ensuremath{\cup}}
\newunicodechar{∩}{\ensuremath{\cap}}
\newunicodechar{≡}{\ensuremath{\equiv}}
\newunicodechar{⊂}{\ensuremath{\subset}}
\newunicodechar{⊥}{\ensuremath{\perp}}
\newunicodechar{∃}{\ensuremath{\exists}}
\newunicodechar{∀}{\ensuremath{\forall}}
\newunicodechar{∛}{\ensuremath{\sqrt[3]{\,}}}
\newunicodechar{∜}{\ensuremath{\sqrt[4]{\,}}}
\newunicodechar{⃗}{\ensuremath{{}^{\to}}}
\newunicodechar{ⁿ}{\ifmmode^{n}\else\textsuperscript{\ensuremath{n}}\fi}
\newunicodechar{ˣ}{\ifmmode^{x}\else\textsuperscript{\ensuremath{x}}\fi}
\newunicodechar{ᵇ}{\ifmmode^{b}\else\textsuperscript{\ensuremath{b}}\fi}
\newunicodechar{ʳ}{\ifmmode^{r}\else\textsuperscript{\ensuremath{r}}\fi}
\newunicodechar{ᵘ}{\ifmmode^{u}\else\textsuperscript{\ensuremath{u}}\fi}
\newunicodechar{ʸ}{\ifmmode^{y}\else\textsuperscript{\ensuremath{y}}\fi}
\newunicodechar{ᵃ}{\ifmmode^{a}\else\textsuperscript{\ensuremath{a}}\fi}
\newunicodechar{ᵉ}{\ifmmode^{e}\else\textsuperscript{\ensuremath{e}}\fi}
\newunicodechar{ᵏ}{\ifmmode^{k}\else\textsuperscript{\ensuremath{k}}\fi}
\newunicodechar{ᵖ}{\ifmmode^{p}\else\textsuperscript{\ensuremath{p}}\fi}
\newunicodechar{ᴺ}{\ifmmode^{N}\else\textsuperscript{\ensuremath{N}}\fi}
\newunicodechar{ᶜ}{\ifmmode^{c}\else\textsuperscript{\ensuremath{c}}\fi}
\newunicodechar{ʰ}{\ifmmode^{h}\else\textsuperscript{\ensuremath{h}}\fi}
\newunicodechar{ᵠ}{\ifmmode^{q}\else\textsuperscript{\ensuremath{q}}\fi}
\newunicodechar{ₙ}{\ifmmode_{n}\else\textsubscript{\ensuremath{n}}\fi}
\newunicodechar{ₐ}{\ifmmode_{a}\else\textsubscript{\ensuremath{a}}\fi}
\newunicodechar{ᵢ}{\ifmmode_{i}\else\textsubscript{\ensuremath{i}}\fi}
\newunicodechar{ᵣ}{\ifmmode_{r}\else\textsubscript{\ensuremath{r}}\fi}
\newunicodechar{ₖ}{\ifmmode_{k}\else\textsubscript{\ensuremath{k}}\fi}
\newunicodechar{ᵦ}{\ifmmode_{\beta}\else\textsubscript{\ensuremath{\beta}}\fi}
\newunicodechar{ₓ}{\ifmmode_{x}\else\textsubscript{\ensuremath{x}}\fi}
\newfontfamily\popdisplay{ArchivoBlack-Regular.ttf}[Path=../../../fonts/]
\newfontfamily\poplogo{SpaceGrotesk-Bold.ttf}[Path=../../../fonts/]
\newfontfamily\popsemi{PlusJakartaSans-SemiBold.ttf}[Path=../../../fonts/]
\newfontfamily\popxbold{PlusJakartaSans-ExtraBold.ttf}[Path=../../../fonts/]
\newfontfamily\popxboldls{PlusJakartaSans-ExtraBold.ttf}[Path=../../../fonts/,LetterSpace=3.0]

\setlist[enumerate]{leftmargin=1.7em,itemsep=5pt,topsep=4pt}
\setlist[itemize]{leftmargin=1.7em,itemsep=4pt,topsep=4pt,label={\color{poporange}\textbullet}}
\setlength{\parindent}{0pt}
\setlength{\parskip}{5pt}
\renewcommand{\arraystretch}{1.25}

% pgfplots : style Pop par défaut
\pgfplotsset{
  every axis/.append style={axis line style={popink,line width=1pt},tick style={popink},
    grid style={popline,line width=0.5pt},label style={font=\small},tick label style={font=\footnotesize}},
  popcurve/.style={poporange,line width=1.8pt,no markers,smooth},
}
% Même style disponible dans un \draw TikZ simple (hors pgfplots)
\tikzset{popcurve/.style={poporange,line width=1.8pt}}
\tikzset{popgrid/.style={popline,very thin}}
\colorlet{popcurve}{poporange} % le nom du style est parfois utilisé comme couleur

% ============================================================
% Pill / squiggle
% ============================================================
\newcommand{\poppill}[3]{%
  \tikz[baseline=-0.55ex]\node[draw=popink,line width=1.2pt,rounded corners=2.6pt,%
    fill=#2,inner xsep=6.5pt,inner ysep=3pt]{%
    \popxboldls\fontsize{7}{7}\selectfont\color{#3}\MakeUppercase{#1}};%
}
\newcommand{\popsquiggle}[1]{%
  \tikz[baseline=-0.4ex]{
    \draw[#1,line width=2.6pt,line cap=round]
      plot [smooth] coordinates {(0,0.09) (0.34,0.01) (0.68,0.21) (1.02,0.01) (1.36,0.10)};
  }%
}
% Mot-clé surligné (marqueur orange sous le texte)
\newcommand{\cle}[1]{{\popxbold\color{popdark}#1}}

% ============================================================
% En-tête / pied
% ============================================================
\pagestyle{fancy}
\fancyhf{}
\renewcommand{\headrulewidth}{0pt}
\renewcommand{\footrulewidth}{0pt}
\lhead{\raisebox{-0.15ex}{\poplogo\fontsize{13}{13}\selectfont\color{popink}OnBuch\color{poporange}+}}
\rhead{\poppill{\DOCMATIERE\ \textbullet\ \DOCNIVEAU\ \textbullet\ Leçon \DOCLECON}{white}{popink}}
\lfoot{\popxbold\fontsize{7.6}{7.6}\selectfont\color{popmuted}\MakeUppercase{Cours signé OnBuch+}\ \textbullet\ {\popsemi\DOCTITRE}}
\rfoot{\tikz[baseline=-0.6ex]\node[rounded corners=8.5pt,fill=popink,minimum height=17pt,minimum width=17pt,inner xsep=5pt,inner sep=0]{\popxbold\fontsize{8}{8}\selectfont\color{popcream}\thepage\,/\,\pageref*{LastPage}};}

% ============================================================
% Titres
% ============================================================
\newcommand{\chaptitre}[2]{%
  \begin{tcolorbox}[enhanced,colback=poporange,colframe=popink,boxrule=2.6pt,arc=11pt,
      drop shadow=popink,left=15pt,right=15pt,top=12pt,bottom=11pt,before skip=3pt,after skip=18pt]
    \poppill{#2}{popink}{popcream}\par\vspace{9pt}
    {\raggedright\hyphenpenalty=10000\exhyphenpenalty=10000\popdisplay\fontsize{22}{24}\selectfont\color{popcream}\MakeUppercase{#1}\par}
  \end{tcolorbox}
}
% Section numérotée : pastille orange avec le numéro + titre Archivo
\newcounter{popsec}
\newcommand{\coursec}[1]{%
  \stepcounter{popsec}\setcounter{popsub}{0}%
  \par\vspace{16pt}\noindent
  \begin{minipage}{\linewidth}
  \tikz[baseline=-0.7ex]\node[draw=popink,line width=1.6pt,fill=poporange,rounded corners=4pt,minimum size=22pt,inner sep=0]{\popdisplay\fontsize{12}{12}\selectfont\color{popcream}\thepopsec};%
  \hspace{8pt}{\popdisplay\fontsize{15}{17}\selectfont\color{popink}\MakeUppercase{#1}}\par
  \vspace{4pt}\noindent{\color{poporange}\rule{36pt}{3.6pt}}
  \end{minipage}\par\nopagebreak\vspace{8pt}
}
\newcounter{popsub}
\newcommand{\courssub}[1]{%
  \stepcounter{popsub}%
  \par\vspace{9pt}\noindent{\popxbold\fontsize{11.5}{13}\selectfont\color{popink}{\color{poporange}\thepopsec.\thepopsub}\hspace{6pt}#1}\par\nopagebreak\vspace{4pt}
}

% ============================================================
% Boîtes pédagogiques — même recette : fond blanc, bordure épaisse colorée,
% ombre dure encre, badge plein. Titre optionnel [#1].
% ============================================================
\tcbset{popbase/.style={breakable,enhanced,colback=white,boxrule=2.3pt,arc=9pt,
  shadow={3pt}{-3pt}{0pt}{popink},left=14pt,right=14pt,top=11pt,bottom=11pt,before skip=11pt,after skip=15pt,
  fontupper=\popsemi\fontsize{10.6}{14.5}\selectfont\color{popink},
  fontlower=\popsemi\fontsize{10.6}{14.5}\selectfont\color{popink}}}
% Coche dessinée (le glyphe ✓ n'existe pas dans les polices du document)
\newcommand{\popcheck}[1]{\tikz[baseline=-0.5ex]\draw[#1,line width=1.6pt,line cap=round,line join=round](-0.13,0.0)--(-0.04,-0.09)--(0.13,0.1);}
\providecommand{\coche}{\popcheck{popgreen}}
\providecommand{\resultat}[1]{\par\smallskip\noindent\fcolorbox{popgreen}{popgreenL}{\parbox{\dimexpr\linewidth-2\fboxsep-2\fboxrule\relax}{\textbf{Résultat.} #1}}\par\smallskip}
\newcommand{\popboxhead}[3]{\poppill{#1}{#2}{white}\ifx\relax#3\relax\else\hspace{6pt}{\popxbold\fontsize{10.6}{12}\selectfont\color{popink}#3}\fi\par\vspace{7pt}}

\newtcolorbox{aretenir}[1][]{popbase,colframe=popgreen,before upper={\popboxhead{À retenir}{popgreen}{#1}}}
\newtcolorbox{exemplebox}[1][]{popbase,colframe=poporange,before upper={\popboxhead{Exemple}{poporange}{#1}}}
\newtcolorbox{definition}[1][]{popbase,colframe=popblue,before upper={\popboxhead{Définition}{popblue}{#1}}}
\newtcolorbox{propriete}[1][]{popbase,colframe=poppurple,before upper={\popboxhead{Propriété}{poppurple}{#1}}}
\newtcolorbox{methode}[1][]{popbase,colframe=popgold,before upper={\popboxhead{Méthode}{popgold}{#1}}}
\newtcolorbox{attention}[1][]{popbase,colframe=poppink,before upper={\popboxhead{Attention}{poppink}{#1}}}
\newtcolorbox{experience}[1][]{popbase,colframe=popblue,colback=popblueL,before upper={\popboxhead{Expérience}{popblue}{#1}}}
% « Le savais-tu ? » : carte violette pleine (contexte, culture, Cameroun)
\newtcolorbox{savaistu}[1][]{popbase,colback=poppurple,colframe=popink,
  fontupper=\popsemi\fontsize{10.4}{14.2}\selectfont\color{white},
  before upper={\poppill{Le savais-tu\,?}{white}{poppurple}\ifx\relax#1\relax\else\hspace{6pt}{\popxbold\color{white}#1}\fi\par\vspace{7pt}}}
% Exercice résolu : énoncé en haut, \tcblower puis la solution sur fond crème
\newcounter{popexo}\newcounter{popexr}
\newtcolorbox{exoresolu}[1][]{popbase,colframe=poporange,colbacklower=popcream,
  segmentation style={popink,line width=1.2pt,dashed},
  before upper={\stepcounter{popexr}\popboxhead{Exercice résolu \thepopexr}{poporange}{#1}},
  before lower={\poppill{Solution}{popink}{popcream}\par\vspace{6pt}}}
% Exercice (non résolu en place) — la correction est dans la partie Corrigés
\newtcolorbox{exercice}[1][]{popbase,colframe=popink,
  before upper={\stepcounter{popexo}\popboxhead{Exercice \thepopexo}{popink}{#1}}}
% Corrigé
\newtcolorbox{corrige}[1][]{popbase,colframe=popgreen,colback=popgreenL,
  before upper={\popboxhead{Corrigé}{popgreen}{#1}}}
% Objectifs / prérequis / bilan
\newtcolorbox{objectifs}{popbase,colframe=popink,colback=poporangeL,
  before upper={\popboxhead{Objectifs de la leçon}{popink}{}}}
\newtcolorbox{prerequis}{popbase,colframe=popink,colback=popblueL,
  before upper={\popboxhead{Prérequis}{popblue}{}}}
\newtcolorbox{bilan}{popbase,colframe=popink,colback=white,boxrule=2.8pt,
  before upper={\popboxhead{Fiche bilan}{poporange}{L'essentiel à connaître pour l'examen}}}
% Figure encadrée avec légende
\newtcolorbox{popfigure}[1][]{enhanced,breakable=false,colback=white,colframe=popline,boxrule=1.4pt,arc=8pt,
  left=10pt,right=10pt,top=10pt,bottom=8pt,before skip=10pt,after skip=12pt,halign=center,
  lower separated=false,halign lower=center,
  fontlower=\popsemi\fontsize{9}{11.5}\selectfont\color{popmuted},
  #1}
\newcommand{\legende}[1]{\par\vspace{6pt}{\popsemi\fontsize{9}{11.5}\selectfont\color{popmuted}#1}}

% ============================================================
% Bloc de code source — \begin{lstlisting}[language=Php]...\end{lstlisting}
% (ou language=C, HTML, JavaScript, SQL ; option omise = pseudo-code brut).
% Même recette visuelle que les figures (\popfigure) : cadre encre épais,
% fond clair (popcream), légende possible juste après via \legende{...}
% comme pour une figure (listings ne sait pas arrondir ses coins : on garde
% un cadre rectangulaire, cohérent avec les autres tableaux du document).
% CONTRAT : les modèles ne doivent utiliser QUE \begin{lstlisting}[...] tel
% quel (pas de \begin{tcblisting}, pas de \verb pour un bloc multi-lignes).
% ============================================================
\lstdefinestyle{poplst}{
  backgroundcolor=\color{popcream},
  basicstyle=\popsemi\fontsize{8.6}{11}\selectfont\color{popink},
  keywordstyle=\bfseries\color{popblue},
  commentstyle=\itshape\color{popmuted},
  stringstyle=\color{popgreen},
  numberstyle=\tiny\color{popmuted},
  identifierstyle=\color{popink},
  frame=l,
  rulecolor=\color{popink},
  framerule=1.6pt,
  framesep=7pt,
  framexleftmargin=7pt,
  framexrightmargin=7pt,
  xleftmargin=2pt,
  xrightmargin=2pt,
  breaklines=true,
  breakatwhitespace=false,
  showstringspaces=false,
  showspaces=false,
  showtabs=false,
  tabsize=2,
  columns=flexible,
  keepspaces=true,
  numbers=none,
  aboveskip=10pt,
  belowskip=10pt,
  captionpos=b,
}
\lstset{style=poplst, language=PHP}
% La distribution TeX embarquée ne fournit pas le fichier de langue
% JavaScript de listings (lstlang*.dat) : on la redéfinit nous-mêmes
% pour que \begin{lstlisting}[language=JavaScript] fonctionne.
\lstdefinelanguage{JavaScript}{
  keywords={break,case,catch,class,const,continue,debugger,default,delete,do,
    else,export,extends,finally,for,function,if,import,in,instanceof,let,new,
    return,static,super,switch,this,throw,try,typeof,var,void,while,with,yield,
    async,await,of,get,set},
  keywordstyle=\bfseries\color{popblue},
  ndkeywords={true,false,null,undefined,NaN,Infinity},
  ndkeywordstyle=\bfseries\color{poporange},
  identifierstyle=\color{popink},
  sensitive=true,
  comment=[l]{//},
  morecomment=[s]{/*}{*/},
  string=[b]",
  morestring=[b]',
  morestring=[b]`,
}
% Même limitation pour CSS et les fichiers de configuration (apacheconf,
% nginx, ini...) : ils ne sont pas fournis. CSS et un style générique de
% type "config" (commentaires # ou ;, pas de mots-clés) couvrent les besoins.
\lstdefinelanguage{CSS}{
  keywords={color,background,background-color,margin,padding,border,
    border-radius,font,font-size,font-weight,font-family,width,height,
    display,position,top,left,right,bottom,float,clear,text-align,
    line-height,list-style,cursor,overflow,box-shadow,transition,
    flex,grid,align-items,justify-content,z-index},
  keywordstyle=\bfseries\color{popblue},
  identifierstyle=\color{popink},
  sensitive=false,
  comment=[s]{/*}{*/},
  string=[b]",
  morestring=[b]',
}
\lstdefinelanguage{apacheconf}{
  sensitive=false,
  morecomment=[l]{\#},
}
\lstdefinelanguage{Apache}{
  sensitive=false,
  morecomment=[l]{\#},
}
\lstdefinelanguage{text}{sensitive=false}
\lstdefinelanguage{powershell}{
  keywords={New-Object,New-SmbShare,Get-Acl,Set-Acl,Get-ChildItem,Set-Content,
    Get-Content,Write-Host,ForEach-Object,Where-Object,Select-Object,
    Test-Path,Remove-Item,Copy-Item,Move-Item,Start-Process,Stop-Process,
    If,Else,ElseIf,ForEach,While,Function,Return,Try,Catch},
  keywordstyle=\bfseries\color{popblue},
  identifierstyle=\color{popink},
  sensitive=false,
  comment=[l]{\#},
  string=[b]",
  morestring=[b]',
}
\lstdefinelanguage{ini}{
  sensitive=false,
  morecomment=[l]{;},
  morecomment=[l]{\#},
}

% ============================================================
% Signature OnBuch+
% ============================================================
\newcommand{\poplogoinline}[1]{{\poplogo\fontsize{#1}{#1}\selectfont\color{popink}OnBuch\color{poporange}+}}
\newcommand{\signatureonbuch}{%
  \begin{tcolorbox}[enhanced,colback=popink,colframe=popink,boxrule=2.2pt,arc=10pt,
      drop shadow=poporange,left=14pt,right=14pt,top=10pt,bottom=10pt,before skip=0pt,after skip=0pt]
    \tikz[baseline=-0.7ex]\node[circle,fill=popgreen,draw=popcream,line width=1.3pt,minimum size=22pt,inner sep=0]{\popcheck{white}};%
    \hspace{9pt}\begin{minipage}[c]{0.62\linewidth}
      {\popxbold\fontsize{12}{13}\selectfont\color{popcream}Signé\ \textbullet\ {\poplogo OnBuch\color{poporange}+}}\par\vspace{2pt}
      {\raggedright\popsemi\fontsize{8}{10.5}\selectfont\color{popline}Ludovic A.\\\MakeUppercase{Conforme au programme officiel MINESEC}\par}
    \end{minipage}\hfill
    {\poplogo\fontsize{22}{22}\selectfont\color{popcream}OnBuch\color{poporange}+}
  \end{tcolorbox}%
}

% ============================================================
% Page de garde
% #1 titre · #2 matière · #3 niveau · #4 module · #5 n° leçon · #6 durée ·
% #7 description / objectifs (texte ou itemize)
% ============================================================
\newcommand{\pagegarde}[7]{%
  \thispagestyle{empty}
  \newgeometry{a4paper,margin=1.7cm,top=1.6cm,bottom=1.4cm}%
  \begin{tikzpicture}[remember picture,overlay]
    \fill[poporange!18] ([xshift=-3.2cm,yshift=-3.6cm]current page.north east) circle (4.6cm);
    \fill[poppurple!14] ([xshift=2.4cm,yshift=7.6cm]current page.south west) circle (3.8cm);
  \end{tikzpicture}%
  {\poplogo\fontsize{30}{30}\selectfont\color{popink}OnBuch\color{poporange}+}\hfill
  \raisebox{0.6ex}{\poppill{Cours signé \textbullet\ Premium}{popink}{popcream}}\par
  \vspace{1pt}\popsquiggle{poppurple}\par
  \vspace{8pt}
  \begin{tcolorbox}[enhanced,colback=poporange,colframe=popink,boxrule=2.8pt,arc=13pt,
      drop shadow=popink,left=18pt,right=18pt,top=12pt,bottom=13pt,after skip=10pt]
    \poppill{#2 \textbullet\ #3}{popink}{popcream}\hspace{5pt}\poppill{#4}{white}{popink}\par\vspace{8pt}
    {\popdisplay\fontsize{14}{14}\selectfont\color{popink}LEÇON #5}\par\vspace{4pt}
    {\raggedright\hyphenpenalty=10000\exhyphenpenalty=10000\popdisplay\fontsize{25}{27}\selectfont\color{popcream}\MakeUppercase{#1}\par}
    \vspace{8pt}
    {\popxbold\fontsize{9.5}{9.5}\selectfont\color{popink}\MakeUppercase{Durée conseillée : #6}}
  \end{tcolorbox}
  \begin{tcolorbox}[enhanced,colback=white,colframe=popink,boxrule=2.2pt,arc=10pt,
      drop shadow=popink,left=16pt,right=16pt,top=10pt,bottom=10pt,after skip=8pt]
    \poppill{Dans cette leçon}{poppurple}{white}\par\vspace{5pt}
    \popsemi\fontsize{9.3}{12.6}\selectfont\color{popink}#7
  \end{tcolorbox}
  \noindent\begin{minipage}[t]{0.487\linewidth}
    \begin{tcolorbox}[enhanced,colback=popgreen,colframe=popink,boxrule=2.2pt,arc=10pt,
        drop shadow=popink,left=11pt,right=11pt,top=10pt,bottom=10pt,height=3.1cm,valign=center]
      \tcbox[colback=white,colframe=popink,boxrule=1.3pt,arc=5pt,boxsep=0pt,nobeforeafter,
        left=3pt,right=3pt,top=3pt,bottom=3pt,valign=center]{\qrcode[height=1.9cm]{https://play.google.com/store/apps/details?id=com.luvvix.onbuch&hl=fr}}%
      \hspace{8pt}\begin{minipage}[c]{0.5\linewidth}
        {\popdisplay\fontsize{11}{12.5}\selectfont\color{white}\MakeUppercase{Télécharge OnBuch}}\par\vspace{4pt}
        {\raggedright\popsemi\fontsize{8.4}{11}\selectfont\color{white}Gratuit \textbullet\ Google Play\\Tuteur IA, annales, concours, résultats\par}
      \end{minipage}
    \end{tcolorbox}
  \end{minipage}\hfill
  \begin{minipage}[t]{0.487\linewidth}
    \begin{tcolorbox}[enhanced,colback=poppurple,colframe=popink,boxrule=2.2pt,arc=10pt,
        drop shadow=popink,left=11pt,right=11pt,top=10pt,bottom=10pt,height=3.1cm,valign=center]
      \tikz[baseline=-0.7ex]\node[circle,fill=white,draw=popink,line width=1.3pt,minimum size=1.6cm,inner sep=0]{\popdisplay\fontsize{24}{24}\selectfont\color{poppurple}+};%
      \hspace{8pt}\begin{minipage}[c]{0.6\linewidth}
        {\popdisplay\fontsize{11}{12.5}\selectfont\color{white}\MakeUppercase{Passe à OnBuch+}}\par\vspace{4pt}
        {\raggedright\popsemi\fontsize{8.4}{11}\selectfont\color{white}Tous les cours signés, Tuteur IA illimité, fiches PDF — onglet OnBuch+ de l'app\par}
      \end{minipage}
    \end{tcolorbox}
  \end{minipage}
  \vspace{8pt}
  \popcontenu
  \vfill
  \signatureonbuch
  \restoregeometry
  \newpage
}

% Rangée « Ce cours contient » (page de garde)
\newcommand{\poptile}[3]{% #1 couleur #2 gros texte #3 légende
  \begin{tcolorbox}[enhanced,colback=#1,colframe=popink,boxrule=2pt,arc=9pt,shadow={2.5pt}{-2.5pt}{0pt}{popink},
      left=6pt,right=6pt,top=9pt,bottom=9pt,halign=center,valign=center,height=1.75cm,width=\linewidth,nobeforeafter]
    {\popdisplay\fontsize{12.5}{13}\selectfont\color{white}\MakeUppercase{#2}}\par\vspace{3pt}
    {\centering\popsemi\fontsize{7.8}{9.5}\selectfont\color{white}#3\par}
  \end{tcolorbox}}
\newcommand{\popcontenu}{%
  \par\noindent{\popxboldls\fontsize{7.5}{7.5}\selectfont\color{popmuted}CE COURS SIGNÉ CONTIENT}\par\vspace{7pt}
  \noindent\begin{minipage}[t]{0.235\linewidth}\poptile{popblue}{Cours}{complet, illustré, pas à pas}\end{minipage}\hfill
  \begin{minipage}[t]{0.235\linewidth}\poptile{poporange}{Méthodes}{et exercices résolus}\end{minipage}\hfill
  \begin{minipage}[t]{0.235\linewidth}\poptile{poppink}{Exercices}{type examen + corrigés}\end{minipage}\hfill
  \begin{minipage}[t]{0.235\linewidth}\poptile{popgreen}{Fiche bilan}{l'essentiel à retenir}\end{minipage}\par}

% Sommaire
\newtcolorbox{sommaire}{enhanced,colback=white,colframe=popink,boxrule=2.2pt,arc=10pt,
  drop shadow=popink,left=18pt,right=18pt,top=13pt,bottom=13pt,before skip=6pt,after skip=20pt,
  before upper={\poppill{Sommaire}{poppurple}{white}\par\vspace{9pt}\popsemi\fontsize{10.6}{14.5}\selectfont\color{popink}}}

% Signature de fin de cours — poussée tout en bas de la dernière page
\newcommand{\findecours}{%
  \par\vfill\nopagebreak
  \begin{center}\popsquiggle{poporange}\end{center}\vspace{4pt}
  \signatureonbuch
}
\AtBeginDocument{\def\llbracket{\mathopen{[\mkern-3.5mu[}}\def\rrbracket{\mathclose{]\mkern-3.5mu]}}} % intervalles d'entiers (glyphe ⟦ absent de la police)
% Glyphes absents de Fira Math : redéfinis à partir de glyphes disponibles
\AtBeginDocument{%
  \def\setminus{\mathbin{\backslash}}\let\smallsetminus\setminus
  \def\vdots{\mathord{\vbox{\baselineskip3.2pt\lineskiplimit0pt\kern4pt\hbox{.}\hbox{.}\hbox{.}}}}%
  \def\ddots{\mathinner{\mkern1mu\raise6.5pt\hbox{.}\mkern2mu\raise3.6pt\hbox{.}\mkern2mu\raise0.7pt\hbox{.}\mkern1mu}}%
  \def\triangle{\mathord{\scalebox{0.9}{$\symup{\Delta}$}}}%
  \def\nabla{\mathord{\raisebox{\depth}{\scalebox{1}[-1]{$\symup{\Delta}$}}}}%
  \def\checkmark{\popcheck{popdark}}%
  \def\star{\mathbin{\ast}}\def\bigstar{\mathord{\ast}}%
  \def\diamond{\mathbin{\scalebox{0.6}{$\Diamond$}}}%
  \def\bigcup{\mathop{\mathchoice{\raisebox{-0.3em}{\scalebox{1.7}{$\cup$}}}{\scalebox{1.25}{$\cup$}}{\cup}{\cup}}\displaylimits}%
  \def\bigcap{\mathop{\mathchoice{\raisebox{-0.3em}{\scalebox{1.7}{$\cap$}}}{\scalebox{1.25}{$\cap$}}{\cap}{\cap}}\displaylimits}%
  \def\lesseqqgtr{\mathrel{\lessgtr}}\def\bigoplus{\mathop{\oplus}\displaylimits}%
}
\providecommand{\ssi}{si et seulement si} % abréviation fréquente
\AtBeginDocument{\providecommand{\wideparen}{\overparen}} % arc de cercle

%% FIGFIX-BEGIN v4 — correctifs globaux des figures (docs/onbuch-generation/tools/figfix)
\usepackage{adjustbox}\usepackage{etoolbox}
% toute figure TikZ/pgfplots plus large que la ligne est réduite à la largeur disponible
\BeforeBeginEnvironment{tikzpicture}{\begin{adjustbox}{max width=\linewidth}}
\AfterEndEnvironment{tikzpicture}{\end{adjustbox}}
% légendes pgfplots lisibles par-dessus les courbes
\pgfplotsset{every axis/.append style={legend style={fill=white,fill opacity=0.92,text opacity=1,draw=black!15,inner sep=3pt}}}
% étiquettes d'arêtes (syntaxe edge["..."]) avec fond blanc
\tikzset{every edge quotes/.append style={fill=white,inner sep=1.2pt}}
%% FIGFIX-END
\begin{document}

\pagegarde{Maintenance de premier niveau}{Informatique}{3ème}{Informatique – classe de 3ème}{N03}{2 séances (fiche de progression harmonisée)}{Cette leçon initie les élèves à la maintenance informatique de premier niveau : entretien matériel et logiciel courant, diagnostic des pannes simples, installation et mise à jour d'applications. Elle vise à rendre l'élève autonome face aux situations concrètes de la vie quotidienne et scolaire.
\vspace{4pt}
\begin{multicols}{2}\small
\begin{itemize}
\item À la fin de cette leçon, je sais définir la maintenance informatique et distinguer maintenance préventive, corrective et évolutive
\item À la fin de cette leçon, je sais identifier les opérations de maintenance de premier niveau réalisables par un simple utilisateur
\item À la fin de cette leçon, je sais appliquer les règles de sécurité avant toute intervention sur un ordinateur
\item À la fin de cette leçon, je sais nettoyer physiquement un poste de travail (écran, clavier, unité centrale) avec le matériel adapté
\item À la fin de cette leçon, je sais effectuer une maintenance logicielle simple : suppression de fichiers inutiles, défragmentation, désinstallation propre
\item À la fin de cette leçon, je sais installer une application simple à partir d'une source fiable et vérifier son bon fonctionnement
\end{itemize}\end{multicols}}

\begin{sommaire}
\begin{multicols}{2}
\begin{enumerate}
\item Notions de base sur la maintenance informatique
\item Règles de sécurité avant toute intervention
\item Maintenance matérielle de premier niveau
\item Maintenance logicielle de premier niveau
\item Installation d'une application simple
\item Mise à jour des applications et démarche de diagnostic
\item Activité d'intégration
\item Méthodes et exercices résolus
\item Exercices
\item Corrigés des exercices
\item Fiche bilan
\end{enumerate}
\end{multicols}
\end{sommaire}

\begin{objectifs}
\begin{itemize}
\item À la fin de cette leçon, je sais définir la maintenance informatique et distinguer maintenance préventive, corrective et évolutive
\item À la fin de cette leçon, je sais identifier les opérations de maintenance de premier niveau réalisables par un simple utilisateur
\item À la fin de cette leçon, je sais appliquer les règles de sécurité avant toute intervention sur un ordinateur
\item À la fin de cette leçon, je sais nettoyer physiquement un poste de travail (écran, clavier, unité centrale) avec le matériel adapté
\item À la fin de cette leçon, je sais effectuer une maintenance logicielle simple : suppression de fichiers inutiles, défragmentation, désinstallation propre
\item À la fin de cette leçon, je sais installer une application simple à partir d'une source fiable et vérifier son bon fonctionnement
\item À la fin de cette leçon, je sais mettre à jour une application et un antivirus, et justifier l'importance de ces mises à jour
\item À la fin de cette leçon, je sais rédiger et justifier une démarche de diagnostic face à une panne simple de la vie courante
\end{itemize}
\end{objectifs}

\begin{prerequis}
\begin{itemize}
\item Notions sur les constituants d'un ordinateur (unité centrale, périphériques d'entrée et de sortie) vues en 4ème
\item Distinction entre matériel (hardware) et logiciel (software)
\item Notion de système d'exploitation et de logiciel d'application
\item Manipulations de base : allumer/éteindre correctement un ordinateur, ouvrir un dossier, désinstaller via le panneau de configuration
\item Notions élémentaires sur les virus informatiques et l'antivirus
\end{itemize}
\end{prerequis}

\newpage

\coursec{Notions de base sur la maintenance informatique}

\courssub{Situation d'accroche et problématique}

Au cybercafé de quartier à Yaoundé, l'ordinateur de M. Nkoulou met dix minutes à démarrer, affiche des messages d'erreur et refuse d'ouvrir ses fichiers de comptabilité. Un « technicien » lui réclame \num{15000} FCFA pour une simple remise en état. Pourtant, la plupart de ces pannes peuvent être prévenues ou résolues par l'utilisateur lui-même, sans matériel coûteux. Un peu de poussière retirée à temps, un câble mal branché revérifié, un logiciel mis à jour : autant de gestes simples qui auraient évité à M. Nkoulou cette dépense et cette perte de temps.

\textbf{Problématique :} comment reconnaître une panne simple, la réparer en toute sécurité et maintenir son ordinateur en bon état de fonctionnement ? C'est tout l'objet de la \cle{maintenance de premier niveau}, que nous allons découvrir dans cette leçon. Mais avant d'intervenir sur une machine, il faut d'abord comprendre ce qu'est la maintenance informatique et quels sont ses différents types.

\courssub{Qu'est-ce que la maintenance informatique ?}

\textbf{Intuition.} Pense à ta bicyclette ou à la moto de ta famille : pour qu'elle fonctionne bien longtemps, il faut graisser la chaîne, vérifier les freins, gonfler les pneus. Et quand elle tombe en panne, il faut la réparer. Un ordinateur, c'est pareil : c'est une machine qui s'use, se salit, se désorganise. Sans entretien, il ralentit, surchauffe, plante. L'ensemble des gestes qui permettent de le garder en bon état — ou de le remettre en état — porte un nom : la \cle{maintenance}.

\begin{definition}[Maintenance informatique]
La \cle{maintenance informatique} est l'ensemble des opérations techniques visant à \textbf{maintenir} ou à \textbf{rétablir} un équipement informatique (ordinateur, imprimante, réseau, logiciel) en \textbf{état de fonctionnement}, c'est-à-dire apte à remplir correctement le service pour lequel il a été conçu.
\end{definition}

Deux idées sont essentielles dans cette définition :
\begin{itemize}
\item \emph{maintenir} : agir \textbf{avant} la panne, pour qu'elle n'arrive pas ou le plus tard possible ;
\item \emph{rétablir} : agir \textbf{après} la panne, pour remettre l'équipement en service.
\end{itemize}

La maintenance ne concerne pas seulement le \cle{matériel} (l'unité centrale, l'écran, le clavier, l'imprimante...) : elle concerne aussi le \cle{logiciel} (le système d'exploitation, les applications, l'antivirus). On distingue donc la \textbf{maintenance matérielle} de la \textbf{maintenance logicielle}, que nous détaillerons dans les sections suivantes de cette leçon.

\begin{aretenir}[L'essentiel]
La maintenance informatique regroupe toutes les opérations qui gardent un équipement en état de marche (avant la panne) ou qui le remettent en état de marche (après la panne). Elle porte sur le matériel \emph{et} sur les logiciels.
\end{aretenir}

\courssub{Les trois types de maintenance}

Selon le moment où l'on intervient et le but recherché, on distingue trois types de maintenance. Les deux premiers sont au cœur du programme ; le troisième est simplement mentionné.

\textbf{1. La maintenance préventive : agir avant la panne.}

\textbf{Intuition.} Il vaut mieux balayer régulièrement sa maison que d'attendre que la saleté la rende inhabitable. De même, il vaut mieux dépoussiérer son ordinateur tous les mois que d'attendre qu'il surchauffe et grille.

\begin{definition}[Maintenance préventive]
La \cle{maintenance préventive} regroupe les opérations effectuées \textbf{régulièrement}, alors que l'équipement fonctionne encore, afin d'\textbf{éviter les pannes} ou d'en retarder l'apparition.
\end{definition}

Exemples d'opérations préventives :
\begin{itemize}
\item le \textbf{nettoyage} : dépoussiérage de l'unité centrale et du clavier, nettoyage de l'écran ;
\item les \textbf{mises à jour} du système d'exploitation, des logiciels et de l'antivirus ;
\item les \textbf{sauvegardes} régulières des fichiers importants (sur clé USB, disque externe ou en ligne) ;
\item la \textbf{vérification} des câbles, des prises et de la ventilation ;
\item la suppression des fichiers inutiles et la désinstallation des logiciels qui ne servent plus.
\end{itemize}

\textbf{2. La maintenance corrective : agir après la panne.}

\textbf{Intuition.} Quand le pneu de la moto est crevé, on ne peut plus « prévenir » : il faut réparer. La maintenance corrective, c'est la réparation.

\begin{definition}[Maintenance corrective]
La \cle{maintenance corrective} regroupe les opérations effectuées \textbf{après une panne}, afin de \textbf{rétablir} le fonctionnement normal de l'équipement : diagnostic de la panne, réparation ou remplacement de l'élément défectueux.
\end{definition}

Exemples d'opérations correctives :
\begin{itemize}
\item réinstaller un logiciel qui ne démarre plus ;
\item remplacer une barrette de mémoire vive défectueuse ou un câble d'alimentation coupé ;
\item supprimer un virus qui bloque le système ;
\item réparer ou réinstaller le système d'exploitation après un plantage grave.
\end{itemize}

\textbf{3. La maintenance évolutive : améliorer l'équipement.}

Il existe enfin un troisième type, la \cle{maintenance évolutive}, qui consiste à \textbf{améliorer les performances} de l'équipement sans qu'il soit en panne : ajouter de la mémoire vive (RAM) pour accélérer l'ordinateur, remplacer un disque dur par un disque plus rapide, installer une version plus récente d'un logiciel. Nous la mentionnons simplement ici.

\begin{attention}[Piège classique au BEPC]
Confondre maintenance \textbf{préventive} et maintenance \textbf{corrective} est une erreur très fréquente au BEPC. Retiens le critère de distinction : c'est le \textbf{moment} de l'intervention par rapport à la panne. La maintenance préventive se fait \emph{avant} la panne (l'équipement fonctionne encore) ; la maintenance corrective se fait \emph{après} la panne (l'équipement est en panne). Ainsi, « mettre à jour l'antivirus » est préventif, mais « supprimer un virus qui a bloqué la machine » est correctif. Et « ajouter de la mémoire vive à un ordinateur qui fonctionne mais qui est lent » est de la maintenance évolutive, pas corrective !
\end{attention}

\begin{popfigure}
\begin{center}
\begin{tikzpicture}[font=\small]
% Colonne 1 : préventive
\draw[fill=popgreenL, draw=popgreen, rounded corners=3pt] (0,0) rectangle (4.6,5.4);
\node[fill=popgreen, text=white, rounded corners=2pt, minimum width=4.6cm, minimum height=0.7cm, anchor=north] at (2.3,5.4) {\textbf{PRÉVENTIVE}};
\node[anchor=north west, text width=4.2cm] at (0.2,4.5) {\textbf{Quand ?} Avant la panne\\[2pt]\textbf{But :} éviter les pannes\\[4pt]\textbf{Exemples :}\\ $\bullet$ dépoussiérage\\ $\bullet$ mises à jour\\ $\bullet$ sauvegardes\\ $\bullet$ vérification des câbles};
% Colonne 2 : corrective
\draw[fill=poporangeL, draw=poporange, rounded corners=3pt] (5.2,0) rectangle (9.8,5.4);
\node[fill=poporange, text=white, rounded corners=2pt, minimum width=4.6cm, minimum height=0.7cm, anchor=north] at (7.5,5.4) {\textbf{CORRECTIVE}};
\node[anchor=north west, text width=4.2cm] at (5.4,4.5) {\textbf{Quand ?} Après la panne\\[2pt]\textbf{But :} rétablir le fonctionnement\\[4pt]\textbf{Exemples :}\\ $\bullet$ réparer un logiciel\\ $\bullet$ remplacer une pièce\\ $\bullet$ supprimer un virus};
% Colonne 3 : évolutive
\draw[fill=popblueL, draw=popblue, rounded corners=3pt] (10.4,0) rectangle (15,5.4);
\node[fill=popblue, text=white, rounded corners=2pt, minimum width=4.6cm, minimum height=0.7cm, anchor=north] at (12.7,5.4) {\textbf{ÉVOLUTIVE}};
\node[anchor=north west, text width=4.2cm] at (10.6,4.5) {\textbf{Quand ?} À tout moment\\[2pt]\textbf{But :} améliorer les performances\\[4pt]\textbf{Exemples :}\\ $\bullet$ ajouter de la RAM\\ $\bullet$ changer de disque dur\\ $\bullet$ nouvelle version d'un logiciel};
\end{tikzpicture}
\end{center}
\legende{Comparaison des trois types de maintenance informatique.}
\end{popfigure}

\begin{propriete}[Intérêt de la maintenance préventive (admise)]
Une maintenance préventive \textbf{régulière} réduit la \textbf{fréquence} et le \textbf{coût} des pannes : un équipement entretenu tombe moins souvent en panne, dure plus longtemps et ses réparations, quand elles sont nécessaires, sont moins graves et moins coûteuses.
\end{propriete}

Cette propriété prend tout son sens dans le contexte camerounais. D'une part, les réparations sont chères par rapport aux revenus : faire appel à un technicien coûte souvent entre \num{5000} et \num{20000} FCFA, alors qu'un dépoussiérage ne coûte presque rien. D'autre part, notre environnement est particulièrement agressif pour les ordinateurs : la \textbf{poussière} (surtout en saison sèche et dans les villes du Nord comme Maroua ou Garoua), la \textbf{chaleur}, l'\textbf{humidité} en saison des pluies et les \textbf{coupures d'électricité} répétées qui endommagent les disques durs et les alimentations. Un ordinateur non entretenu dans ces conditions peut voir sa durée de vie divisée par deux ou trois.

\begin{exemplebox}[Le dépoussiérage avant la saison sèche à Maroua]
Chaque année, à l'approche de la saison sèche, le vent ramène beaucoup de poussière à Maroua. Le proviseur d'un collège fait ouvrir les unités centrales de la salle informatique par le surveillant formé à cet effet : ordinateurs éteints et débranchés, il retire la poussière accumulée sur les ventilateurs et les grilles d'aération à l'aide d'un pinceau doux et d'un souffleur. Résultat : les ventilateurs tournent librement, l'air circule, les machines ne surchauffent plus. Coût de l'opération : presque nul. Coût d'un bloc d'alimentation grillé par la surchauffe : environ \num{10000} FCFA par machine, sans compter l'immobilisation du poste pendant les réparations. Voilà la maintenance préventive en action.
\end{exemplebox}

\begin{savaistu}[La poussière, ennemie numéro un en zone tropicale]
La poussière est la première cause de surchauffe des ordinateurs en zone tropicale. Elle s'accumule sur les ventilateurs et les grilles d'aération, forme une couche isolante sur les composants et empêche la chaleur de s'évacuer. Le processeur, qui peut dépasser \SI{70}{\celsius} en fonctionnement, se protège alors en ralentissant (l'ordinateur devient très lent) puis en s'éteignant brutalement. À la longue, la chaleur excessive détruit les composants. Un simple dépoussiérage tous les trois à six mois suffit souvent à éviter ces pannes.
\end{savaistu}

\courssub{Les niveaux de maintenance}

Toutes les pannes ne se réparent pas de la même façon, et tout le monde ne peut pas tout réparer. On classe les interventions en \cle{niveaux de maintenance}, selon leur difficulté et les compétences qu'elles exigent.

\begin{definition}[Maintenance de premier niveau]
La \cle{maintenance de premier niveau} est l'ensemble des opérations de maintenance \textbf{simples}, réalisables par l'\textbf{utilisateur} lui-même, sans formation technique poussée et \textbf{sans ouverture approfondie} du matériel : nettoyage extérieur, vérification des branchements, redémarrage, mises à jour, désinstallation de logiciels.
\end{definition}

Au-delà du premier niveau, on trouve les \textbf{niveaux supérieurs}, réservés aux professionnels :
\begin{itemize}
\item le \textbf{deuxième niveau} : interventions du \textbf{technicien spécialisé} (ouverture de l'unité centrale, remplacement d'une barrette mémoire, d'un disque dur ou d'une alimentation, réinstallation complète du système) ;
\item le \textbf{troisième niveau} : réparations complexes effectuées par le \textbf{constructeur} ou un centre agréé (soudure sur la carte mère, réparation d'un écran, échange standard du matériel sous garantie).
\end{itemize}

\begin{popfigure}
\begin{center}
\begin{tikzpicture}[font=\small]
% Pyramide en 3 étages (du bas vers le haut)
\draw[fill=popgreenL, draw=popgreen] (0,0) -- (10,0) -- (8.4,2) -- (1.6,2) -- cycle;
\draw[fill=poporangeL, draw=poporange] (1.6,2) -- (8.4,2) -- (6.8,4) -- (3.2,4) -- cycle;
\draw[fill=popblueL, draw=popblue] (3.2,4) -- (6.8,4) -- (5,6) -- cycle;
\node at (5,1) {\textbf{Niveau 1 : l'utilisateur}};
\node[text=popmuted] at (5,0.45) {\footnotesize nettoyage, câbles, redémarrage, mises à jour, désinstallation};
\node at (5,3) {\textbf{Niveau 2 : le technicien spécialisé}};
\node[text=popmuted] at (5,2.5) {\footnotesize remplacement de pièces, réinstallation du système};
\node[text=white] at (5,4.9) {\textbf{Niveau 3 : le constructeur}};
\node[text=white] at (5,4.45) {\footnotesize \hspace{0.6cm}réparations complexes, garantie\hspace{0.6cm}};
\end{tikzpicture}
\end{center}
\legende{Pyramide des niveaux de maintenance : plus on monte, plus les interventions sont complexes et spécialisées.}
\end{popfigure}

\textbf{Exemples d'opérations de premier niveau} (à ta portée, donc !) :
\begin{itemize}
\item le \textbf{nettoyage} : dépoussiérage extérieur, nettoyage du clavier et de l'écran avec un chiffon doux ;
\item la \textbf{vérification des câbles} : s'assurer que les prises d'alimentation, d'écran, de clavier et de souris sont bien enfoncées ;
\item le \textbf{redémarrage} de l'ordinateur, qui règle à lui seul une grande partie des petits blocages ;
\item les \textbf{mises à jour} du système, des logiciels et de l'antivirus ;
\item la \textbf{désinstallation} des logiciels inutiles ou suspects.
\end{itemize}

\begin{attention}[Connais tes limites]
Le premier niveau s'arrête là où commence le risque. Ouvrir l'unité centrale alors qu'elle est branchée, toucher les composants internes sans précaution, démonter un écran ou une alimentation : ce sont des opérations de niveau supérieur, dangereuses pour toi (électricité) et pour la machine (décharges électrostatiques). Une règle d'or : en cas de doute, on s'arrête et on fait appel à un technicien.
\end{attention}

\courssub{Exercice résolu et bilan}

\begin{exoresolu}[Identifier le type et le niveau d'une maintenance]
Dans chacune des situations suivantes, indique le \textbf{type} de maintenance (préventive, corrective ou évolutive) et le \textbf{niveau} (utilisateur, technicien, constructeur) :
\begin{enumerate}
\item La secrétaire du lycée sauvegarde chaque vendredi les fichiers des bulletins sur une clé USB.
\item L'ordinateur de la salle informatique ne démarre plus ; le professeur remplace le câble d'alimentation coupé par un câble neuf.
\item Un élève désinstalle un jeu qui ralentit l'ordinateur familial.
\item Le constructeur répare la carte mère d'un ordinateur portable encore sous garantie.
\item On ajoute une barrette de mémoire vive à un ordinateur qui fonctionne mais qui est devenu lent.
\end{enumerate}
\tcblower
\textbf{Solution détaillée.}
\begin{enumerate}
\item Sauvegarde régulière alors que tout fonctionne : maintenance \textbf{préventive}, niveau \textbf{utilisateur} (premier niveau).
\item L'ordinateur est en panne (il ne démarre plus) : maintenance \textbf{corrective}. Remplacer un câble externe est une opération simple, sans ouverture du matériel : niveau \textbf{utilisateur}.
\item Désinstaller un logiciel pour éviter les ralentissements : maintenance \textbf{préventive} (l'ordinateur fonctionne encore), niveau \textbf{utilisateur}.
\item Réparation d'une carte mère : maintenance \textbf{corrective} (panne), niveau \textbf{constructeur} (troisième niveau, intervention très complexe).
\item L'ordinateur n'est pas en panne, on améliore ses performances : maintenance \textbf{évolutive}. L'ajout de mémoire exige l'ouverture de l'unité centrale : niveau \textbf{technicien} (deuxième niveau).
\end{enumerate}
\end{exoresolu}

\begin{exercice}[À toi de jouer]
\begin{enumerate}
\item Définis la maintenance informatique et cite ses deux grands objectifs.
\item Pour chaque opération, dis si elle relève de la maintenance préventive ou corrective : (a) mettre à jour l'antivirus ; (b) supprimer un virus qui bloque le démarrage ; (c) dépoussiérer le clavier ; (d) réinstaller le système après un plantage.
\item Explique en quelques lignes pourquoi la maintenance préventive est particulièrement importante dans une ville comme Maroua ou Garoua.
\end{enumerate}
\end{exercice}

\begin{corrige}[Exercice]
\begin{enumerate}
\item La maintenance informatique est l'ensemble des opérations visant à maintenir ou à rétablir un équipement informatique en état de fonctionnement. Ses deux objectifs : \emph{maintenir} (agir avant la panne) et \emph{rétablir} (agir après la panne).
\item (a) préventive ; (b) corrective ; (c) préventive ; (d) corrective.
\item À Maroua ou Garoua, la saison sèche apporte beaucoup de poussière qui bouche les ventilateurs et provoque des surchauffes ; les coupures d'électricité y sont fréquentes et endommagent le matériel. Un dépoussiérage régulier et des sauvegardes fréquentes évitent des pannes coûteuses, alors que les réparations y sont chères et les pièces de rechange parfois difficiles à trouver.
\end{enumerate}
\end{corrige}

\begin{aretenir}[Bilan de la section]
\begin{itemize}
\item La \cle{maintenance informatique} maintient ou rétablit un équipement en état de fonctionnement ; elle concerne le matériel et les logiciels.
\item Trois types : \textbf{préventive} (avant la panne : nettoyage, mises à jour, sauvegardes), \textbf{corrective} (après la panne : réparation, remplacement), \textbf{évolutive} (amélioration des performances).
\item Trois niveaux : \textbf{utilisateur} (premier niveau, opérations simples), \textbf{technicien spécialisé}, \textbf{constructeur}.
\item Une maintenance préventive régulière réduit la fréquence et le coût des pannes — un réflexe précieux dans notre contexte de poussière, de chaleur et de coupures d'électricité.
\end{itemize}
\end{aretenir}

Maintenant que tu sais ce qu'est la maintenance et comment elle se classe, la prochaine étape avant de toucher à la moindre machine est d'apprendre les \textbf{règles de sécurité} : c'est l'objet de la section suivante.

\coursec{Règles de sécurité avant toute intervention}

Après avoir découvert ce qu'est la maintenance de premier niveau et pourquoi elle est indispensable pour prolonger la vie de nos équipements, nous devons maintenant aborder une règle d'or : \textbf{la sécurité prime sur tout le reste}. Au Cameroun, où les variations de tension sur le réseau ENEO sont fréquentes et où l'humidité peut jouer des tours, respecter ces consignes n'est pas une option, c'est une condition de survie pour ton matériel... et pour toi ! Une intervention mal préparée peut transformer un simple nettoyage en court-circuit destructeur, voire en accident corporel. Voyons ensemble comment te préparer comme un professionnel.

\courssub{Les dangers invisibles : électricité statique et tensions résiduelles}

Avant même de toucher à un tournevis, il faut comprendre deux phénomènes physiques sournois qui menacent les composants électroniques.

\begin{definition}[Électricité statique (ESD - ElectroStatic Discharge)]
L'\cle{électricité statique} est une accumulation de charges électriques à la surface d'un corps (ton corps, tes vêtements, un sac en plastique). Lorsque tu touches un composant sensible (barrette de RAM, carte mère, processeur), cette charge se décharge brutalement : c'est la \cle{décharge électrostatique (DES)}. Elle peut détruire instantanément ou affaiblir irrémédiablement les circuits intégrés, sans laisser de trace visible.
\end{definition}

\begin{propriete}[Seuil de sensibilité des composants]
Le corps humain peut accumuler jusqu'à \textbf{\SIrange{10000}{25000}{\volt}} en marchant sur un tapis synthétique. Or, les composants CMOS modernes (RAM, CPU, puces de carte mère) peuvent être endommagés par une décharge de seulement \textbf{\SIrange{30}{100}{\volt}}, bien en dessous du seuil de perception humaine (environ \SI{3000}{\volt}). \textbf{Tu ne sentiras rien, mais le composant, lui, aura « pris le courant ».}
\end{propriete}

\begin{savaistu}[Le condensateur « fantôme »]
Savais-tu que le \cle{bloc d'alimentation (PSU)} contient de gros condensateurs électrolytiques qui stockent l'énergie ? Même \textbf{après avoir débranché le cordon secteur}, ces condensateurs peuvent conserver une charge dangereuse (souvent $> \SI{200}{\volt}$) pendant \textbf{plusieurs minutes, voire une heure}. C'est pour cette raison qu'on \textbf{ne ouvre JAMAIS un bloc d'alimentation} ni un écran CRT (tube cathodique) en maintenance de premier niveau. Seul un technicien qualifié, équipé d'outils de décharge adaptés, peut intervenir dessus.
\end{savaistu}

\courssub{La check-list de survie : 5 étapes avant d'ouvrir le boîtier}

Voici la procédure standard, à suivre \emph{à la lettre} avant toute intervention matérielle (nettoyage, ajout de RAM, changement de disque dur, remplacement de la pile CMOS...).

\begin{methode}[Les 5 étapes de préparation avant une intervention matérielle]
\begin{enumerate}
    \item \textbf{Éteindre proprement l'ordinateur} via le menu du système d'exploitation (Arrêter), pas en maintenant le bouton Power enfoncé (sauf blocage total). Cela parque les têtes de disque dur et ferme les fichiers ouverts.
    \item \textbf{Débrancher le cordon d'alimentation secteur} à l'arrière du boîtier (prise IEC C13/C14). Attends \textbf{au moins 30 secondes} (idéalement 1 à 2 minutes) pour laisser les condensateurs se décharger naturellement.
    \item \textbf{Sauvegarder les données critiques} sur un support externe (clé USB, disque dur externe, cloud) \textbf{avant} toute opération logicielle risquée (réinstallation, formatage, mise à jour majeure du BIOS/UEFI).
    \item \textbf{Se décharger de l'électricité statique} : touche une partie métallique \emph{non peinte} du châssis du boîtier (ou un radiateur métallique relié à la terre) pendant quelques secondes. L'idéal est de porter un \cle{bracelet antistatique} relié à la terre.
    \item \textbf{Documenter l'intervention} : note la date, l'opération prévue, l'état initial (symptômes), les pièces remplacées dans ton \cle{cahier de suivi de maintenance} (papier ou fichier numérique).
\end{enumerate}
\end{methode}

\begin{attention}[Risque mortel : ne jamais intervenir sur un poste sous tension]
\textbf{Ne JAMAIS ouvrir un ordinateur encore branché sur le secteur.} Même si l'interrupteur arrière du bloc d'alimentation (I/O) est sur « O », le \SI{220}{\volt} alternatif arrive encore jusqu'à l'interrupteur et aux fils internes. Risque 1 : \textbf{électrocution} (courant traversant le corps). Risque 2 : \textbf{court-circuit} si un tournevis tombe sur la carte mère sous tension $\rightarrow$ destruction immédiate de la carte mère, du CPU, de la RAM. Risque 3 : \textbf{incendie} par étincelle. \textbf{La règle est simple : Câble débranché = Sécurité. Câble branché = Danger.}
\end{attention}

\courssub{Protéger le poste contre le réseau électrique (Contexte ENEO)}

Au Cameroun, le réseau électrique (ENEO) subit des \cle{coupures intempestives}, des \cle{sous-tensions} (baisse de tension) et des \cle{surtensions} (pics brutaux, souvent au retour du courant après une coupure). Ces événements sont la première cause de mortalité des blocs d'alimentation et des cartes mères.

\begin{propriete}[Équipements de protection recommandés]
\begin{itemize}
    \item \textbf{L'onduleur (UPS - Uninterruptible Power Supply) :} C'est la protection \textbf{idéale}. Il contient une batterie qui prend le relais instantanément en cas de coupure (autonomie \SIrange{10}{30}{\minute} typiques), permet d'éteindre proprement, et filtre les surtensions/sous-tensions. Indispensable pour un poste de travail critique (serveur, caisse, poste de l'admin).
    \item \textbf{Le régulateur de tension (AVR - Automatic Voltage Regulator) :} Il corrige les variations de tension (ex: maintient \SI{220}{\volt} en entrée \SIrange{160}{260}{\volt}) mais \textbf{ne protège pas contre les coupures} (pas de batterie). Moins cher qu'un onduleur, c'est le minimum vital pour un PC fixe au Cameroun.
    \item \textbf{La multiprise parafoudre (surge protector) :} Protection de base contre les pics de foudre ou de commutation. Elle \textbf{ne régule pas la tension} et ne sauve pas en cas de coupure. À réserver pour les périphériques non critiques (imprimante, écran, chargeur téléphone).
\end{itemize}
\end{propriete}

\begin{exemplebox}[Choisir sa protection pour un cybercafé à Yaoundé]
Pour une salle de 10 PC fixes + 1 serveur de facturation :
\begin{itemize}
    \item \textbf{Serveur :} Onduleur \textit{Online} (double conversion) \SI{1500}{\volt\ampere} / \SI{900}{\watt} $\rightarrow$ extinction propre garantie, données comptables sauvées.
    \item \textbf{10 PC clients :} Régulateurs AVR \SI{1000}{\volt\ampere} (ou onduleurs \textit{Line-Interactive} \SI{650}{\volt\ampere} si budget permet) $\rightarrow$ protection contre les variations ENEO quotidiennes.
    \item \textbf{Imprimantes/Scanners :} Multiprises parafoudre certifiées NF/CE.
\end{itemize}
\textit{Astuce :} Vérifie régulièrement l'état des batteries des onduleurs (test d'autonomie tous les 6 mois). Une batterie morte = onduleur inutile.
\end{exemplebox}

\courssub{Les interdits absolus : eau, humidité et limites physiques}

\begin{propriete}[Règles d'hygiène et de sécurité physiques]
\begin{enumerate}
    \item \textbf{Mains sèches impératives :} Ne touche jamais un composant, un connecteur ou l'intérieur du boîtier avec les mains mouillées, humides ou grasses. L'eau + électricité = court-circuit + corrosion (oxydation des contacts dorés).
    \item \textbf{Environnement sec :} N'interviens pas dans une pièce humide (salle de bain, buanderie, extérieur sous la pluie, local inondé). Attends que le matériel soit à température ambiante et sec (risque de condensation si on sort un PC du froid vers le chaud).
    \item \textbf{Boissons interdites :} Aucun verre, bouteille ou tasse près du poste ouvert. Un renversement de jus de mangue sur une carte mère allumée = perte totale quasi certaine.
    \item \textbf{Ne pas forcer :} Un connecteur qui ne rentre pas n'est pas le bon ou n'est pas dans le bon sens. Forcer casse les broches (ex: port SATA, connecteur alimentation CPU 4/8 broches, slot RAM).
\end{enumerate}
\end{propriete}

\begin{aretenir}[Les « Zones Rouges » : Interdiction formelle d'ouverture en maintenance de premier niveau]
\begin{itemize}
    \item \textbf{Bloc d'alimentation (PSU) :} Condensateurs haute tension (\SIrange{200}{400}{\volt} DC) $\rightarrow$ danger de mort. Pas de pièce réparable par l'utilisateur (pas de fusible accessible en général).
    \item \textbf{Écran CRT (tube cathodique) :} Anode du tube $\rightarrow$ \SIrange{15000}{30000}{\volt} (résiduels très longtemps). Risque d'implosion du tube (éclats de verre projetés).
    \item \textbf{Écran LCD/LED (panneau) :} Inverseur/ballast (haute tension pour le rétroéclairage CCFL sur anciens modèles) + dalle fragile. Pas de réparation utilisateur.
    \item \textbf{Batterie d'onduleur :} Plomb-acide (acide sulfurique, hydrogène explosif) ou Lithium (risque incendie). Remplacement complet du bloc batterie seulement.
\end{itemize}
\textbf{Si l'un de ces éléments est en cause :} Appelle un technicien qualifié ou utilise la garantie / le SAV constructeur.
\end{aretenir}

\courssub{La sauvegarde : ton assurance-vie numérique}

Toute intervention logicielle (désinstallation, nettoyage de registre, mise à jour Windows, réinstallation OS, flashage BIOS) comporte un risque de perte de données. La règle est simple : \textbf{pas de sauvegarde = pas d'intervention}.

\begin{exemplebox}[Sauvegarder ses compositions avant de désinstaller une suite bureautique]
\textbf{Situation :} Ton PC rame, tu veux désinstaller \texttt{LibreOffice} pour le réinstaller propre. Tu as 3 compositions de français (fichiers \texttt{.odt}) et ton exposé d'histoire (\texttt{.pptx}) sur le Bureau.
\textbf{Démarche :}
\begin{enumerate}
    \item Branche ta clé USB (vérifie qu'elle est saine : \texttt{chkdsk F: /f} en invite de commande admin).
    \item Crée un dossier \texttt{Sauvegarde\_JJMMAA} (ex: \texttt{Sauvegarde\_240115}).
    \item Copie-colle (ou \texttt{Ctrl+C / Ctrl+V}) tes 4 fichiers importants dedans.
    \item \textbf{Vérifie} sur la clé que les fichiers s'ouvrent correctement (double-clic).
    \item \textbf{Seulement alors}, lance la désinstallation via \texttt{Paramètres > Applications > Applications installées}.
\end{enumerate}
\textit{Gagne-temps :} Utilise la fonction « Historique des fichiers » (Windows) ou « Time Machine » (Mac) pour une sauvegarde automatique continue sur disque externe.
\end{exemplebox}

\courssub{Le cahier de suivi : la mémoire de la machine}

Un bon technicien (même débutant) laisse une trace. Cela permet de savoir \emph{quoi}, \emph{quand}, \emph{pourquoi} et \emph{comment} a été faite une intervention. C'est indispensable pour la garantie, le dépannage futur ou la revente.

\begin{methode}[Tenir un cahier de suivi de maintenance (papier ou tableur)]
Colonnes recommandées (à créer dans un classeur \texttt{Maintenance\_PC.xlsx} ou un cahier 96 pages) :
\begin{enumerate}
    \item \textbf{Date} (JJ/MM/AAAA)
    \item \textbf{Poste concerné} (Nom PC, N° inventaire, Emplacement : ex: \texttt{PC-ADMIN-01 / Bureau Directeur})
    \item \textbf{Type d'intervention} (Matérielle / Logicielle / Nettoyage / Mise à jour)
    \item \textbf{Symptôme / Motif} (ex: « Bip 1 long 2 courts au démarrage », « Écran bleu 0x0000007E », « Maintenance préventive trimestrielle »)
    \item \textbf{Action réalisée} (ex: « Nettoyage ventilateurs + pâte thermique CPU », « Réinstallation Windows 10 22H2 + drivers », « Remplacement pile CR2032 »)
    \item \textbf{Pièces remplacées / Références} (ex: « RAM DDR4 8Go 3200MHz Kingston KVR32N22S8/8 », « Pile CR2032 Sony »)
    \item \textbf{Résultat / Observations} (ex: « OK - Plus de bip, T° CPU \SI{38}{\celsius} au repos », « Échec - Disque dur HS, à changer »)
    \item \textbf{Technicien / Élève} (Ton nom / signature)
    \item \textbf{Prochaine échéance} (ex: « Nettoyage poussière dans 6 mois », « Vérif MAJ Windows dans 1 mois »)
\end{enumerate}
\end{methode}

\begin{experience}[TP Guidés : Préparer son poste pour une intervention de nettoyage interne]
\textbf{Objectif :} Appliquer la procédure complète de mise en sécurité avant d'ouvrir un boîtier pour dépoussiérage.
\textbf{Matériel :} Un PC fixe (éteint), un tournevis cruciforme (PH2), une bombe à air comprimé, un pinceau antistatique (ou pinceau doux propre), un bracelet antistatique (si dispo), une clé USB, le cahier de suivi.
\textbf{Déroulement :}
\begin{enumerate}
    \item \textbf{Sauvegarde :} Copie tes documents récents (Bureau, Documents) sur la clé USB. Vérifie l'ouverture d'un fichier sur la clé.
    \item \textbf{Extinction :} Menu \texttt{Démarrer > Arrêter}. Attends extinction complète (ventilateurs arrêtés, LEDs éteintes).
    \item \textbf{Coupure secteur :} Débranche le cordon d'alimentation \textbf{à l'arrière du boîtier}. Débranche aussi l'écran, la souris, le clavier, le réseau (RJ45) pour avoir de l'espace.
    \item \textbf{Décharge statique :} Met le bracelet antistatique (pince crocodile sur partie métallique nue du châssis) OU touche le châssis métallique nu pendant 5 secondes avec les deux mains.
    \item \textbf{Ouverture :} Repère les vis du panneau latéral (souvent 2 vis pouce à l'arrière). Dévisse, fais glisser le panneau vers l'arrière et retire-le. Pose-le loin du plan de travail.
    \item \textbf{Observation :} Identifie : ventilateur CPU, ventilateur boîtier, alimentation, carte graphique, barrettes RAM, disques. Note la poussière accumulée.
    \item \textbf{Nettoyage (simulation) :} Tiens les pales des ventilateurs avec un doigt (pour ne pas les faire tourner trop vite = générateur). Souffle la poussière \textbf{vers l'extérieur} du boîtier à l'air comprimé (bombe verticale, bursts courts). Brosses doucement les ailettes du radiateur CPU.
    \item \textbf{Fermeture :} Remets le panneau, revisse. Rebranche alimentation, écran, périphériques.
    \item \textbf{Test :} Allume. Écoute les bruits anormaux. Vérifie T° au BIOS ou via HWMonitor au bout de 5 min.
    \item \textbf{Documentation :} Remplis une ligne dans le cahier de suivi (Date, Poste, Type: Matérielle/Nettoyage, Action: Dépoussiérage complet + pâte thermique si faite, Résultat: OK).
\end{enumerate}
\textbf{Interprétation :} Si le PC démarre normalement et que les températures ont baissé (ex: -5 à -\SI{10}{\celsius}), l'intervention est un succès. Si un ventilateur gratte ou ne tourne pas $\rightarrow$ le remplacer (maintenance curative).
\end{experience}

\begin{exoresolu}[Analyse de risque : Intervention sur un PC portable]
\textbf{Énoncé :} Tu dois remplacer la pâte thermique du processeur d'un portable \texttt{HP ProBook 450 G8}. Le portable est allumé, sur secteur, batterie interne non amovible. Liste les étapes de sécurité dans l'ordre, en justifiant.
\textbf{Solution détaillée :}
\begin{enumerate}
    \item \textbf{Sauvegarder les données} sur clé USB/Disque externe/Cloud. \textit{Justification :} Démontage complet carte mère $\rightarrow$ risque de court-circuit, perte données.
    \item \textbf{Éteindre complètement} (Arrêter, pas Veille/Hibernation). \textit{Justification :} Coupe l'alimentation logicielle des composants.
    \item \textbf{Débrancher l'adaptateur secteur} (prise murale + prise portable). \textit{Justification :} Supprime la source \SI{19}{\volt} continu.
    \item \textbf{Débrancher la batterie interne} (étape critique portable). Ouvrir le capot arrière (vis Torx souvent), déconnecter le connecteur batterie $\rightarrow$ carte mère. \textit{Justification :} Même éteint, la batterie alimente le circuit de charge et le RTC. Risque court-circuit si tournevis touche zone alimentation.
    \item \textbf{Appuyer 10 sec sur bouton Power} (après débranchement batterie). \textit{Justification :} Vide les condensateurs de la carte mère (fleeting power).
    \item \textbf{Se décharger (bracelet ou touche châssis)}. \textit{Justification :} Protection ESD sur composants nus (CPU, GPU, PCH).
    \item \textbf{Procéder au démontage} (ventirad, nettoyage, pâte, remontage).
    \item \textbf{Remonter dans l'ordre inverse} : Batterie $\rightarrow$ Capot $\rightarrow$ Secteur $\rightarrow$ Allumage $\rightarrow$ Test.
\end{enumerate}
\textbf{Piège évité :} Ne JAMAIS commencer par ouvrir le capot tant que la batterie n'est pas déconnectée électriquement.
\end{exoresolu}

\begin{exercice}[Mise en situation : Le cybercafé de quartier]
Tu gères un cybercafé de 8 postes à Douala. L'onduleur du serveur de facturation vient de rendre l'âme (batteries gonflées). Le patron veut que tu installes un nouveau régulateur AVR \SI{2000}{\volt\ampere} en attendant l'onduleur de remplacement.
\begin{enumerate}
    \item Liste les 5 étapes de sécurité à respecter avant de toucher au câblage arrière du serveur.
    \item Pourquoi est-il dangereux de brancher le serveur directement sur la prise murale « en attendant » ?
    \item Quelle information dois-tu impérativement noter dans le cahier de suivi pour cette intervention ?
\end{enumerate}
\end{exercice}

\begin{corrige}[Exercice n°1 : Mise en situation : Le cybercafé de quartier]
\begin{enumerate}
    \item \textbf{1. Sauvegarde :} Copie la base de données de facturation (fichier \texttt{.fdb} ou \texttt{.mdf} + logs) sur disque externe. \textbf{2. Arrêt propre :} Fermeture logiciel caisse $\rightarrow$ Arrêt Windows Server. \textbf{3. Débranchement secteur :} Retire prise onduleur HS du mur ET du serveur. Attends 1 min. \textbf{4. Décharge statique :} Touche châssis serveur nu. \textbf{5. Documentation :} Note l'intervention dans le cahier avant de câbler le régulateur.
    \item \textbf{Risque majeur :} Sans onduleur ni régulateur, la prochaine coupure ENEO (très fréquente à Douala) coupe net le serveur $\rightarrow$ corruption base de données (tables InnoDB/MyISAM cassées, index brisés) $\rightarrow$ perte chiffre d'affaires, impossibilité d'éditer les reçus clients. La surtension au retour du courant peut griller l'alimentation du serveur (pas de varistance de protection).
    \item \textbf{Informations obligatoires :} Date/Heure, Poste (Serveur-Facturation), Type (Matérielle - Remplacement protection électrique), Motif (Onduleur HS batteries gonflées), Action (Installation régulateur AVR \SI{2000}{\volt\ampere} marque X modèle Y, câblage direct secteur $\rightarrow$ AVR $\rightarrow$ Serveur), Pièce (Régulateur AVR \SI{2000}{\volt\ampere} N° série 12345), Résultat (Serveur redémarré OK, tension stable \SI{228}{\volt} affichée), Technicien (Ton nom), Prochaine échéance (Commande nouvel onduleur \SI{1500}{\volt\ampere} Online).
\end{enumerate}
\end{corrige}

\begin{popfigure}
\begin{tikzpicture}[
    node distance=1.2cm and 1.5cm,
    box/.style={draw=popdark, thick, rounded corners=6pt, fill=popcream, text width=3.2cm, align=center, font=\small, minimum height=1.8cm},
    icon/.style={font=\Large, text=popblue},
    arrow/.style={-Stealth, thick, popblue},
    title/.style={font=\bfseries\large, text=popdark, align=center}
]
% Title
\node[title] (titre) at (0, 4.5) {\textbf{Check-list Sécurité : Avant d'ouvrir le boîtier}};

% Column 1: Electrical Safety
\node[box, anchor=north west, align=center] (elec) at (-7, 3.5){
    \\[4pt]
    \textbf{Électricité}\\
    \begin{itemize}
        \item[] \textcolor{popgreen}{✓} Éteindre (OS)
        \item[] \textcolor{popgreen}{✓} Débrancher cordon secteur
        \item[] \textcolor{popgreen}{✓} Attendre 1--2 min (condensateurs)
        \item[] \textcolor{popred}{✗} Ne JAMAIS ouvrir bloc alim / écran
    \end{itemize}
};

% Column 2: ESD Protection
\node[box, anchor=north, align=center] (esd) at (0, 3.5){
    \\[4pt]
    \textbf{Antistatique (ESD)}\\
    \begin{itemize}
        \item[] \textcolor{popgreen}{✓} Bracelet antistatique (idéal)
        \item[] \textcolor{popgreen}{✓} Sinon : toucher châssis métal nu
        \item[] \textcolor{popgreen}{✓} Travail sur tapis antistatique
        \item[] \textcolor{popred}{✗} Pas de tapis synthétique, pas de laine
    \end{itemize}
};

% Column 3: Environment & Data
\node[box, anchor=north east, align=center] (env) at (7, 3.5){
    \\[4pt]
    \textbf{Environnement \& Données}\\
    \begin{itemize}
        \item[] \textcolor{popgreen}{✓} Mains sèches, pièce sèche
        \item[] \textcolor{popgreen}{✓} Pas de boisson à proximité
        \item[] \textcolor{popgreen}{✓} \textbf{Sauvegarde} avant action logicielle
        \item[] \textcolor{popgreen}{✓} Cahier de suivi à jour
    \end{itemize}
};

% Bottom: Protection Réseau
\node[box, anchor=north, fill=popblueL!50, text width=10cm] (reseau) at (0, -0.5) {
    \hspace{5pt} \textbf{Protection Réseau (Contexte Cameroun / ENEO)} \\
    \begin{tabular}{>{\bfseries}l l}
        \textcolor{popgreen}{✓} \textbf{Onduleur (UPS)} : Coupures + Surtensions + Régulation & (Postes critiques, Serveurs) \\
        \textcolor{poporange}{✓} \textbf{Régulateur AVR} : Régulation tension seule & (PC fixes standards, Cybercafés) \\
        \textcolor{popred}{✗} \textbf{Multiprise simple} : Aucune protection & (Interdit pour UC / Serveur)
    \end{tabular}
};

% Arrows from title
\draw[arrow] (titre.south) -- ++(0,-0.3) -| (elec.north);
\draw[arrow] (titre.south) -- ++(0,-0.3) -| (esd.north);
\draw[arrow] (titre.south) -- ++(0,-0.3) -| (env.north);
\draw[arrow] (titre.south) -- ++(0,-0.5) -| (reseau.north);

% Legend icons
\node[font=\footnotesize, text=popmuted, anchor=north west] at (-7, -3.5) {
    ⚡ = Danger électrique \quad
    🛡 = Bouclier / Protection \quad
    💾 = Disque / Données \quad
    🏠 = Bâtiment / Infrastructure
};
\end{tikzpicture}
\legende{Synthèse visuelle des consignes de sécurité : électricité, antistatique, environnement/données, protection réseau.}
\end{popfigure}

\courssub{Bilan de la section}

Tu as maintenant toutes les cartes en main pour intervenir \textbf{proprement} et \textbf{en sécurité}. Résumons les points non négociables :
\begin{itemize}
    \item \textbf{Zéro tension :} Éteint $\rightarrow$ Débranché $\rightarrow$ Attente (condensateurs).
    \item \textbf{Zéro statique :} Bracelet ou main sur châssis métallique nu.
    \item \textbf{Zéro liquide :} Mains sèches, pas de boisson, pièce saine.
    \item \textbf{Zéro risque données :} Sauvegarde \emph{avant} toute manipulation logicielle.
    \item \textbf{Zéro improvisation :} On n'ouvre pas le bloc alim, ni l'écran. On note tout dans le cahier de suivi.
    \item \textbf{Zéro « nu » sur le réseau :} Onduleur ou Régulateur AVR obligatoire au Cameroun.
\end{itemize}

Dans la prochaine section, nous passerons à la pratique pure : \textbf{la maintenance matérielle de premier niveau} (nettoyage, remplacement pièces faciles : RAM, disque, pile, carte graphique, ventilateurs). Prépare ton tournevis et ta bombe à air !

\coursec{Maintenance matérielle de premier niveau}

Après avoir assimilé les \emph{règles de sécurité indispensables} — débrancher le cordon d'alimentation, se décharger de l'électricité statique en touchant une pièce métallique mise à la terre (châssis du boîtier), travailler dans un espace éclairé, sec et dégagé — nous passons à l'action concrète. La maintenance matérielle de premier niveau, c'est l'ensemble des gestes simples, réguliers et sans risque que tu peux effectuer toi-même pour prolonger la durée de vie du matériel, prévenir les pannes et garantir un confort d'utilisation optimal. Au Cameroun, où la poussière rouge de la saison sèche et l'humidité de la saison des pluies mettent les équipements à rude épreuve, ces gestes ne sont pas optionnels : ils sont vitaux pour ton ordinateur, qu'il soit au cybercafé du quartier, dans la salle informatique du lycée ou à la maison.

\courssub{Nettoyage des périphériques externes : écran, clavier, souris}

Ces trois éléments sont en contact direct avec l'utilisateur. Ils accumulent poussière, miettes, traces de doigts et, hélas, microbes. Un nettoyage hebdomadaire est l'idéal.

\courssub{Nettoyage de l'écran : douceur et produits adaptés}

L'écran (LCD, LED ou OLED) est fragile. Sa surface possède souvent un traitement antireflet ou anti-traces que les produits agressifs détruisent irrémédiablement.

\begin{definition}[Nettoyage d'un écran]
Opération consistant à retirer la poussière et les traces de doigts sur la dalle d'affichage à l'aide d'un \textbf{chiffon microfibre doux et sec} (type lunettes), légèrement humidifié avec de l'eau distillée ou un \textbf{produit spécifique « écran »} sans alcool, sans ammoniaque, sans acétone. Le mouvement doit être circulaire ou en lignes droites, sans appuyer excessivement.
\end{definition}

\begin{attention}[Produits interdits sur l'écran]
N'utilise \textbf{jamais} : eau du robinet (calcaire = traces blanches), nettoyant vitres (ammoniaque), alcool à 90\%, acétone, essuie-tout ou mouchoirs en papier (micro-rayures), ni spray directement sur la dalle (risque d'infiltration par les bords). Vaporise toujours sur le chiffon, jamais sur l'écran.
\end{attention}

\courssub{Nettoyage du clavier : le nid à microbes}

Le clavier est souvent l'objet le plus sale d'un poste de travail. Entre les touches s'accumulent peaux mortes, cheveux, poussière, résidus alimentaires... et des colonies de bactéries.

\begin{savaistu}[Clavier vs cuvette de toilettes]
Des études microbiologiques ont montré qu'un clavier d'ordinateur partagé et jamais nettoyé peut abriter \textbf{400 fois plus de bactéries} qu'une cuvette de toilettes (staphylocoques, E. coli, etc.). Dans un cybercafé ou une salle de TP, c'est un enjeu de santé publique ! Nettoyer son clavier, c'est aussi se protéger soi-même.
\end{savaistu}

\begin{methode}[Procédure de nettoyage d'un clavier en 4 étapes]
\begin{enumerate}
    \item \textbf{Débrancher} le clavier (ou éteindre l'ordinateur portable) et le retourner au-dessus d'une poubelle ou d'un journal. Secouer doucement pour faire tomber les grosses particules (miettes, cheveux, poussière).
    \item \textbf{Souffler} l'intérieur avec une \textbf{bombe d'air sec} (air comprimé) en maintenant la bombe à la verticale pour éviter le jet de liquide propulseur. Passer entre chaque rangée de touches. Alternative : un pinceau souple antistatique ou un aspirateur à main très faible puissance (embout brosse).
    \item \textbf{Nettoyer les touches} : passer un chiffon microfibre légèrement imbibé d'alcool isopropylique à 70\% (ou lingette désinfectante pour électronique bien essorée) sur le dessus et les flancs des touches. Ne pas laisser de liquide s'infiltrer sous les touches.
    \item \textbf{Laisser sécher} 1 à 2 minutes à l'air libre avant de rebrancher.
\end{enumerate}
\end{methode}

\begin{exemplebox}[Nettoyage en profondeur (avancé, optionnel)]
Si le clavier est très encrassé (boisson renversée, touches qui collent), on peut \emph{retirer les touches} une par une (photo à l'appui pour le remontage !) à l'aide d'un extracteur de touches ou d'un petit tournevis plat, laver les touches à l'eau savonneuse (bien sécher 24h), nettoyer la base au pinceau/air sec, puis remonter. \textbf{Attention :} ne pas faire cela sur un clavier de portable (mécanisme ciseau fragile) ni sous garantie.
\end{exemplebox}

\courssub{Nettoyage de la souris : précision et hygiène}

\begin{definition}[Entretien de la souris optique/laser]
Nettoyage du \textbf{capteur optique} (sous la souris) avec un coton-tige sec ou légèrement humide (alcool 70\%) pour garantir le suivi du curseur, et nettoyage des \textbf{patins de glisse} (pieds en téflon) pour une fluidité parfaite. Nettoyage de la coque et de la molette au chiffon microfibre légèrement humide.
\end{definition}

\begin{propriete}[Surface de travail]
Une souris optique fonctionne mal sur surfaces brillantes (verre, miroir), transparentes ou trop irrégulières. Utilise un \textbf{tapis de souris} propre (tissu ou plastique dur). Un tapis sale use les patins et encrasse le capteur.
\end{propriete}

\begin{popfigure}
\begin{tikzpicture}[
    scale=0.85,
    every node/.style={font=\small},
    box/.style={draw=popblue, thick, rounded corners, fill=popblueL, inner sep=6pt},
    arrow/.style={-Stealth, thick, popdark},
    labelbox/.style={draw=poporange, thick, rounded corners, fill=poporangeL, inner sep=4pt, font=\small\bfseries}
]
% Poste de travail simplifié
% Écran
\node[box, minimum width=6cm, minimum height=4cm] (ecran) at (0,0) {};
\node[labelbox] at (0, 2.5) {ÉCRAN};
\node[align=center, text width=5cm] at (0,0) {Dalle fragile\\Chiffon microfibre\\+ produit spécifique};

% Clavier
\node[box, minimum width=7cm, minimum height=1.5cm] (clavier) at (0,-3.5) {};
\node[labelbox] at (0, -2.5) {CLAVIER};
\node[align=center, text width=6cm] at (0,-3.5) {Secouer $\rightarrow$ Air sec $\rightarrow$ Lingette 70\%};

% Souris
\node[box, minimum width=2.5cm, minimum height=1.5cm] (souris) at (4.5,-3.5) {};
\node[labelbox] at (4.5, -2.5) {SOURIS};
\node[align=center, text width=2.2cm] at (4.5,-3.5) {Capteur (coton-tige)\\Patins + coque};

% Unité centrale
\node[box, minimum width=3.5cm, minimum height=5cm] (uc) at (-5.5, -0.5) {};
\node[labelbox] at (-5.5, 2.2) {UNITÉ CENTRALE};
\node[align=center, text width=3cm] at (-5.5, 0.5) {Grilles aération\\Ventilateurs\\Alim + Câbles};

% Flèches d'annotation
\draw[arrow] (2.5, 1.5) -- ++(1.5,0) node[right, labelbox, anchor=west] {Nettoyage hebdo};
\draw[arrow] (2.5, -3.5) -- ++(1.5,0) node[right, labelbox, anchor=west] {Nettoyage hebdo};
\draw[arrow] (-7.5, 1.5) -- ++(-1.5,0) node[left, labelbox, anchor=east] {Dépoussiérage mensuel};
\draw[arrow] (-7.5, -2) -- ++(-1.5,0) node[left, labelbox, anchor=east] {Vérif. câbles};

% Légende globale
\node[draw=popdark, fill=popcream, rounded corners, inner sep=8pt, align=center, font=\small\bfseries] at (0, -5.8) {Points de maintenance de premier niveau sur un poste de travail fixe};
\end{tikzpicture}
\legende{Poste de travail : zones de nettoyage et fréquence recommandée (hebdomadaire pour périphériques, mensuel pour l'unité centrale).}
\end{popfigure}

\courssub{Entretien de l'unité centrale : le cœur de la machine}

L'unité centrale (tour) aspire l'air pour refroidir les composants (processeur, carte graphique, alimentation). Elle aspire aussi la poussière. Un \emph{ordinateur encrassé = surchauffe = instabilité = panne}.

\begin{exemplebox}[Conséquence d'un ventilateur encrassé]
Dans un cybercafé de Douala, un PC s'éteint brutalement au bout de 20 minutes de jeu vidéo. Le gérant pense à un virus. En ouvrant la tour, on découvre le \textbf{ventilateur du processeur (CPU) bloqué par une « feutre » de poussière} de 5 mm d'épaisseur. Le processeur monte à 95°C, la protection thermique coupe l'alimentation. \textbf{Solution :} soufflage du ventilateur et du radiateur à l'air sec $\rightarrow$ température redescend à 45°C $\rightarrow$ PC stable. Coût : 0 FCFA, 10 minutes.
\end{exemplebox}

\begin{methode}[Dépoussiérage de l'unité centrale (mensuel ou trimestriel selon environnement)]
\begin{enumerate}
    \item \textbf{Éteindre et débrancher} complètement (câble secteur + périphériques). Appuyer 5 sec sur le bouton Power pour vider les condensateurs (décharge résiduelle).
    \item \textbf{Ouvrir le boîtier} (généralement panneau latéral gauche vu de face : 2 vis pouce à l'arrière).
    \item \textbf{Souffler l'intérieur} avec une bombe d'air sec : \emph{grilles d'aération avant/arrière/haut}, \emph{ventilateurs} (CPU, boîtier, alimentation, carte graphique). \textbf{Astuce pro :} bloquer les pales du ventilateur avec un doigt ou un cure-dent en plastique pendant le soufflage pour éviter qu'il ne tourne trop vite (risque de générer du courant inverse ou d'user le roulement).
    \item \textbf{Pinceau antistatique} pour les recoins (slots RAM, PCIe, connecteurs SATA, batterie CMOS). Ne pas toucher les circuits imprimés (puces, pistes) avec les doigts (électricité statique, acidité sueur).
    \item \textbf{Refermer}, rebrancher, démarrer. Écouter : les ventilateurs doivent tourner silencieusement.
\end{enumerate}
\end{methode}

\begin{attention}[Interdits en unité centrale]
\begin{itemize}
    \item \textbf{Aspirateur domestique :} fort champ électromagnétique + électricité statique = carte mère grillée.
    \item \textbf{Chiffon / coton-tige sur composants :} fibres qui accrochent les pattes de composants CMS.
    \item \textbf{Souffler avec la bouche :} humidité (condensation) + projections salivaires = corrosion.
    \item \textbf{Démonter le ventirad (ventilateur + radiateur) du CPU :} nécessite pâte thermique neuve et savoir-faire (niveau supérieur, technicien).
\end{itemize}
\end{attention}

\courssub{Vérification des connexions et environnement d'exploitation}

La moitié des « pannes » signalées en maintenance de premier niveau sont de simples faux contacts ou des conditions d'environnement défavorables.

\begin{propriete}[Check-list connexions physiques (visuel + tactile)]
\begin{itemize}
    \item \textbf{Alimentation :} Câble secteur bien enfoncé dans le bloc alim ET dans la prise murale / multiprise. Interrupteur multiprise sur ON. Témoin LED alim allumé ?
    \item \textbf{Vidéo :} Câble HDMI / VGA / DisplayPort bien vissé (VGA/DVI) ou encliqueté (HDMI/DP) des deux côtés (écran + UC/portable).
    \item \textbf{USB :} Clavier, souris, clé USB, imprimante : bien insérés. Ports USB 3.0 (bleus) pour disques durs externes.
    \item \textbf{Réseau :} Câble RJ45 « cliqueté » dans le port LAN (voyant vert/orange clignotant = lien actif).
\end{itemize}
\end{propriete}

\begin{definition}[Bonnes pratiques d'environnement]
\begin{itemize}
    \item \textbf{Aération :} Laisser 10--15 cm derrière et sur les côtés de la tour pour le flux d'air. Ne pas coller au mur ni enfermer dans un meuble fermé.
    \item \textbf{Propreté :} Sol balayé, pas de tapis épais sous la tour (aspire la poussière). Éviter la proximité d'une fenêtre ouverte (poussière, pluie) ou d'un radiateur/climatiseur direct.
    \item \textbf{Liquides :} \textbf{Zéro boisson} sur le bureau à côté du clavier/unité centrale. Un verre renversé = court-circuit instantané = perte de données + matériel HS.
    \item \textbf{Onduleur / Parafoudre :} Indispensable au Cameroun (coupures fréquentes, surtensions orageuses). Protège l'alimentation et permet un arrêt propre.
\end{itemize}
\end{definition}

\courssub{Diagnostic simple des pannes matérielles courantes}

Face à un symptôme, ne panique pas. Applique une démarche structurée : \textbf{Observer $\rightarrow$ Vérifier les connexions $\rightarrow$ Tester/Isoler $\rightarrow$ Conclure}. C'est la base du dépannage logique.

\begin{methode}[Démarche de diagnostic d'une panne simple (4 étapes)]
\begin{enumerate}
    \item \textbf{Observer et qualifier le symptôme :} Quoi exactement ? (Rien ne s'allume ? Ventilateurs tournent mais écran noir ? Bips au démarrage ? Message d'erreur ?). Depuis quand ? Qu'est-ce qui a changé (déplacement, orage, installation logicielle, chute) ?
    \item \textbf{Vérifier les connexions (source d'énergie + liaison) :} Prise murale OK ? Multiprise ON ? Câble secteur des deux côtés ? Câble vidéo des deux côtés ? Luminosité écran pas au minimum ?
    \item \textbf{Tester / Isoler le composant suspect :} Échanger le câble vidéo avec un autre connu fonctionnel. Brancher l'écran sur un autre PC (ou le PC sur un autre écran). Tester la prise murale avec une lampe. Écouter les bips (codes BIOS).
    \item \textbf{Conclure et agir :} Identifier la cause racine (câble HS, multiprise grillée, alimentation PC morte, écran HS, carte graphique...). Appliquer la correction (remplacer câble, reset CMOS, appeler technicien si composant interne HS).
\end{enumerate}
\end{methode}

\begin{popfigure}
\begin{tikzpicture}[
    scale=0.85,
    every node/.style={font=\small, align=center},
    start/.style={draw=popgreen, thick, rounded corners, fill=popgreenL, inner sep=6pt},
    process/.style={draw=popblue, thick, rectangle, fill=popblueL, inner sep=6pt, minimum width=3.5cm, text width=3.2cm},
    decision/.style={draw=poporange, thick, diamond, aspect=2, fill=poporangeL, inner sep=4pt, text width=2.5cm},
    stop/.style={draw=poppink, thick, rounded corners, fill=poppinkL, inner sep=6pt, text width=3cm},
    arrow/.style={-Stealth, thick, popdark},
    yes/.style={midway, fill=white, inner sep=2pt, font=\small\bfseries, text=popgreen},
    no/.style={midway, fill=white, inner sep=2pt, font=\small\bfseries, text=poppink}
]
\node[start, align=center] (start){ORDINATEUR NE DÉMARRE PAS\\(Aucun voyant, aucun bruit)};
\node[process, below=1.8cm of start, align=center] (check_prise){1. Vérifier la prise murale\\(Autre appareil fonctionne ?)};
\node[decision, below=1.8cm of check_prise] (prise_ok) {Prise OK ?};
\node[process, below=1.8cm of prise_ok, align=center] (check_multi){2. Vérifier multiprise / onduleur\\(Interrupteur ON ? Fusible ?)};
\node[decision, below=1.8cm of check_multi] (multi_ok) {Multiprise OK ?};
\node[process, below=1.8cm of multi_ok, align=center] (check_cable_alim){3. Vérifier câble alimentation PC\\(Bien enfoncé ? Test avec autre câble)};
\node[decision, below=1.8cm of check_cable_alim] (cable_ok) {Câble / Connexion OK ?};
\node[process, below=1.8cm of cable_ok, align=center] (check_bouton){4. Vérifier interrupteur arrière alim (I/O)\\+ Bouton façade (coincé ?)};
\node[decision, below=1.8cm of check_bouton] (bouton_ok) {Interrupteurs OK ?};
\node[stop, below=1.8cm of bouton_ok, align=center] (conclu_alim){PANNE ALIMENTATION PC\\(Bloc alim HS ou Carte mère)\\$\rightarrow$ Technicien requis};

% Branches NON
\node[stop, right=3.5cm of check_prise, align=center] (sol_prise){Prise défaillante\\$\rightarrow$ Électricien};
\node[stop, right=3.5cm of check_multi, align=center] (sol_multi){Multiprise / Onduleur HS\\$\rightarrow$ Remplacer};
\node[stop, right=3.5cm of check_cable_alim, align=center] (sol_cable){Câble alimentation HS\\$\rightarrow$ Remplacer (câble standard)};
\node[stop, right=3.5cm of check_bouton, align=center] (sol_bouton){Interrupteur I/O sur O (0)\\ou Bouton façade bloqué\\$\rightarrow$ Corriger};

\draw[arrow] (start) -- (check_prise);
\draw[arrow] (check_prise) -- (prise_ok);
\draw[arrow] (prise_ok) -- node[yes] {OUI} (check_multi);
\draw[arrow] (prise_ok) -| node[no] {NON} (sol_prise);
\draw[arrow] (check_multi) -- (multi_ok);
\draw[arrow] (multi_ok) -- node[yes] {OUI} (check_cable_alim);
\draw[arrow] (multi_ok) -| node[no] {NON} (sol_multi);
\draw[arrow] (check_cable_alim) -- (cable_ok);
\draw[arrow] (cable_ok) -- node[yes] {OUI} (check_bouton);
\draw[arrow] (cable_ok) -| node[no] {NON} (sol_cable);
\draw[arrow] (check_bouton) -- (bouton_ok);
\draw[arrow] (bouton_ok) -- node[yes] {OUI} (conclu_alim);
\draw[arrow] (bouton_ok) -| node[no] {NON} (sol_bouton);
\end{tikzpicture}
\legende{Arbre de décision : diagnostic « L'ordinateur ne démarre pas » (aucun signe de vie).}
\end{popfigure}

\courssub{Cas particulier : L'écran reste noir (ordinateur allumé)}

C'est le scénario classique : ventilateurs tournent, voyants UC allumés, mais écran noir (ou « Pas de signal »).

\begin{exoresolu}[Diagnostic écran noir]
\textbf{Énoncé :} Tu allumes l'ordinateur de la bibliothèque. Les ventilateurs de la tour tournent, la LED « Power » est bleue, mais l'écran affiche « No Signal » puis se met en veille. Que fais-tu dans l'ordre ?

\textbf{Solution détaillée :}
\begin{enumerate}
    \item \textbf{Observer :} UC sous tension (ventilateurs, LED). Écran sous tension (LED veille allumée souvent orange/jaune) mais « No Signal ». \emph{Conclusion partielle :} Le problème est la liaison vidéo ou la carte graphique, pas l'alimentation générale.
    \item \textbf{Vérifier connexions vidéo :} Câble HDMI/VGA bien vissé/encliqueté des \textbf{deux côtés} (écran + tour). Si tour a carte graphique dédiée + carte mère : le câble \textbf{doit} être sur la carte graphique (ports horizontaux bas), pas sur la carte mère (ports verticaux haut).
    \item \textbf{Tester source écran :} Bouton « Source/Input » sur l'écran $\rightarrow$ sélectionner la bonne entrée (HDMI 1, VGA, DP).
    \item \textbf{Isoler (Test croisé) :}
        \begin{itemize}
            \item Brancher un \textbf{autre câble vidéo} connu bon.
            \item Brancher l'écran sur un \textbf{autre PC} (ou portable) : si ça marche $\rightarrow$ Écran OK, câble OK $\rightarrow$ Problème PC (Carte graphique, RAM, CMOS).
            \item Brancher un \textbf{autre écran} sur ce PC : si ça marche $\rightarrow$ Écran d'origine HS.
        \end{itemize}
    \item \textbf{Reset CMOS (action technicien) :} Si test croisé pointe vers PC : éteindre, débrancher secteur, retirer pile bouton CR2032 carte mère 30 sec (ou cavalier CLR\_CMOS), remettre, redémarrer. Résout souvent un mauvais réglage BIOS ou conflit matériel. \emph{À réserver au personnel formé en 3ème.}
    \item \textbf{Conclure :} Câble HS / Mauvaise source / Carte graphique non alimentée (câble PCIe 6/8 pins manquant) / RAM mal enfichée (bip codes) / Écran HS / Carte mère HS.
\end{enumerate}
\end{exoresolu}

\begin{aretenir}[Bilan : Maintenance matérielle de premier niveau]
\begin{itemize}
    \item \textbf{Régularité :} Périphériques (hebdo) $>$ Unité centrale (mensuel/trimestriel).
    \item \textbf{Outils de base :} Chiffon microfibre, bombe air sec, pinceau antistatique, alcool isopropylique 70\%, tournevis cruciforme.
    \item \textbf{Règle d'or :} \textbf{JAMAIS} de liquide direct, JAMAIS d'aspirateur, JAMAIS de force brutale.
    \item \textbf{Diagnostic :} Logique, progressif, par élimination (test croisé). Note tes observations.
    \item \textbf{Environnement :} Aéré, sec, propre, protégé (onduleur), sans liquide à proximité.
\end{itemize}
\end{aretenir}

\begin{exercice}[Mise en situation : Le PC du secrétariat]
Au secrétariat du lycée, le PC du proviseur adjoint ne démarre plus lundi matin après un week-end d'orage. La multiprise est allumée, la lampe de bureau branchée dessus fonctionne. Le câble d'alimentation du PC est bien enfoncé. L'interrupteur arrière de l'alimentation est sur « I ». Tu appuies sur le bouton Power : \textbf{rien ne se passe, silence total, aucune LED}.
\begin{enumerate}
    \item Propose une démarche de diagnostic en 3 étapes supplémentaires avant d'appeler le technicien de maintenance du lycée.
    \item Quelle est la panne la plus probable au vu du contexte (orage) ?
    \item Quelle protection aurait pu l'éviter ?
\end{enumerate}
\end{exercice}

\begin{corrige}[Exercice n°1 : Le PC du secrétariat]
\begin{enumerate}
    \item \textbf{Étape 1 :} Tester le câble d'alimentation du PC avec un autre câble connu fonctionnel (ou tester ce câble sur l'écran si même connectique IEC C13). \textbf{Étape 2 :} Observer s'il y a une LED allumée sur la carte mère (visible par une grille d'aération) indiquant que le 5V standby arrive. \textbf{Étape 3 :} Sentir l'odeur à l'arrière du bloc alim (odeur de brûlé = composant grillé).
    \item \textbf{Panne probable :} \textbf{Bloc d'alimentation (PSU) grillé} par une surtension due à la foudre (orage). La varistance (protection interne) a claqué. La carte mère peut être indemne si l'alim a bien fait son travail de « fusible ».
    \item \textbf{Protection :} Un \textbf{onduleur (UPS) avec protection contre les surtensions (parafoudre intégré)} ou au minimum une multiprise \textbf{parafoudre normée NF/CE} (pas une simple multiprise 200 FCFA du marché). L'onduleur permet aussi d'éteindre proprement lors de coupure.
\end{enumerate}
\end{corrige}

\coursec{Maintenance logicielle de premier niveau}

Après avoir appris à dépoussiérer l'unité centrale, vérifier les câbles et assurer une bonne ventilation dans la section précédente, nous devons maintenant nous occuper de la « partie invisible » : le logiciel. Un ordinateur dont le matériel est impeccable mais dont le disque est saturé de fichiers inutiles, fragmenté à outrance ou infecté par des logiciels malveillants deviendra tout aussi inutilisable qu'une machine surchauffée. La maintenance logicielle de premier niveau, c'est l'hygiène numérique quotidienne qui garantit la réactivité, la stabilité et la longévité de votre poste de travail.

\courssub{Nettoyage du disque : vider la corbeille et supprimer les fichiers temporaires}

Chaque action que vous effectuez sur l'ordinateur laisse des traces. Lorsque vous supprimez un fichier, il n'est pas effacé immédiatement : il est déplacé dans la \cle{Corbeille}, une zone de stockage temporaire qui permet de restaurer un fichier supprimé par erreur. Tant que la corbeille n'est pas vidée, ces fichiers continuent d'occuper de l'espace disque.

\begin{savaistu}[Les fichiers temporaires : des gigaoctets oubliés]
Saviez-vous que Windows et vos applications créent en permanence des \cle{fichiers temporaires} (fichiers \texttt{.tmp}, caches navigateurs, journaux d'installation, rapports d'erreurs) ? Après quelques mois d'utilisation intensive (navigation web, bureautique, mises à jour), ces fichiers peuvent occuper \textbf{plusieurs gigaoctets} (souvent 5 à 15 Go, parfois davantage sur un poste de cybercafé ou d'administration). Ils ne servent plus à rien une fois la tâche terminée, mais l'ordinateur ne les supprime pas toujours automatiquement.
\end{savaistu}

Pour récupérer cet espace, deux actions simples s'imposent :
\begin{enumerate}
    \item \textbf{Vider la Corbeille :} Clic droit sur l'icône Corbeille $\rightarrow$ \emph{Vider la Corbeille}.
    \item \textbf{Supprimer les fichiers temporaires :} Utilisez l'outil intégré \texttt{Nettoyage de disque} (\texttt{cleanmgr.exe}) ou le raccourci clavier \texttt{Win + R}, tapez \texttt{\%temp\%} et supprimez le contenu du dossier ouvert (ignorez les fichiers en cours d'utilisation).
\end{enumerate}

\begin{lstlisting}[language=bash]
:: Ouvrir l'outil de nettoyage de disque graphique
cleanmgr

:: Supprimer les fichiers temporaires de l'utilisateur courant (prudence : /q = sans confirmation, /s = récursif, /f = forcer lecture-seule)
del /q /f /s %temp%\*

:: Supprimer les fichiers temporaires système (nécessite droits administrateur)
del /q /f /s C:\Windows\Temp\*

:: Vider la corbeille de tous les utilisateurs (Windows 10/11)
rd /s /q C:\$Recycle.Bin
\end{lstlisting}

\courssub{Désinstallation propre d'une application}

Il arrive qu'un logiciel ne serve plus, qu'il soit obsolète ou qu'il cause des conflits. L'erreur classique d'un débutant est de se rendre dans \texttt{C:\textbackslash Program Files} ou \texttt{C:\textbackslash Program Files (x86)}, de supprimer le dossier du logiciel et de croire l'affaire réglée.

\begin{attention}[Supprimer le dossier $\neq$ Désinstaller]
\textbf{Ne supprimez jamais manuellement le dossier d'un logiciel dans \texttt{Program Files}.}
Cela laisse des « résidus » partout dans le système :
\begin{itemize}
    \item Des entrées dans la \cle{Base de registre} (clés de configuration, associations de fichiers, démarrage automatique) qui ralentissent le démarrage et polluent le système.
    \item Des raccourcis orphelins dans le menu Démarrer et sur le Bureau.
    \item Des bibliothèques partagées (\texttt{DLL}) qui ne seront pas désenregistrées.
    \item L'entrée dans la liste « Applications installées » persiste, faussant l'inventaire logiciel.
\end{itemize}
Une désinstallation propre utilise le \cle{programme de désinstallation} fourni par l'éditeur (souvent \texttt{unins000.exe}, \texttt{uninstall.exe} ou une entrée MSI) qui sait nettoyer *tous* ces endroits.
\end{attention}

\begin{methode}[Désinstaller proprement une application sous Windows (4 étapes)]
\begin{enumerate}
    \item Ouvrez \textbf{Paramètres} (\texttt{Win + I}) $\rightarrow$ \textbf{Applications} $\rightarrow$ \textbf{Applications installées} (ou \texttt{Programmes et fonctionnalités} via \texttt{appwiz.cpl}).
    \item Repérez le logiciel dans la liste (utilisez la barre de recherche). Vérifiez l'éditeur et la date d'installation pour ne pas supprimer un pilote ou un composant système (ex: \emph{Microsoft Visual C++ Redistributable}).
    \item Cliquez sur les \textbf{trois points (...)} à droite de l'entrée $\rightarrow$ \textbf{Désinstaller}.
    \item Suivez l'assistant de désinstallation du logiciel. Acceptez la suppression des données utilisateur si vous êtes sûr de ne plus en avoir besoin. Redémarrez l'ordinateur si c'est demandé.
\end{enumerate}
\end{methode}

\courssub{Défragmentation du disque dur : principe et utilité}

Sur un disque dur mécanique (HDD), les données sont écrites sur des plateaux magnétiques tournants. Lorsque vous créez, modifiez et supprimez des fichiers, l'espace libre se morcelle. Un même fichier peut être éclaté en plusieurs \cle{fragments} dispersés à différents endroits du plateau. La tête de lecture doit alors « sauter » d'un fragment à l'autre pour reconstituer le fichier : c'est la \cle{fragmentation}. Cela use la mécanique et ralentit considérablement l'accès aux données.

\begin{definition}[Défragmentation]
La \textbf{défragmentation} est l'opération qui réorganise les données sur un disque dur mécanique (HDD) pour regrouper les fragments de chaque fichier en blocs contigus et rassembler l'espace libre en une seule zone continue. Cela réduit les mouvements de la tête de lecture et accélère les temps d'accès.
\end{definition}

\begin{popfigure}
\begin{tikzpicture}[scale=0.85, every node/.style={font=\small}]
    % Styles
    \tikzset{
        block/.style={rectangle, draw=popdark, minimum height=0.6cm, minimum width=0.6cm, outer sep=0},
        sys/.style={block, fill=popblue!60},
        app/.style={block, fill=poporange!60},
        doc/.style={block, fill=popgreen!60},
        free/.style={block, fill=popline, pattern=north east lines, pattern color=popmuted},
        frag/.style={block, fill=poppink!70},
        label/.style={anchor=north, yshift=-0.2cm, font=\small\bfseries}
    }
    
    % AVANT DÉFRAGMENTATION
    \node[label] at (5, 3.2) {\textbf{AVANT défragmentation (Disque fragmenté)}};
    \foreach \i in {0,...,19} {
        \pgfmathsetmacro{\x}{mod(\i,10)*0.65}
        \pgfmathsetmacro{\y}{-floor(\i/10)*0.8}
        \ifnum\i=0 \node[sys] (b\i) at (\x,\y) {}; \fi
        \ifnum\i=1 \node[sys] (b\i) at (\x,\y) {}; \fi
        \ifnum\i=2 \node[app] (b\i) at (\x,\y) {}; \fi
        \ifnum\i=3 \node[doc] (b\i) at (\x,\y) {}; \fi
        \ifnum\i=4 \node[free] (b\i) at (\x,\y) {}; \fi
        \ifnum\i=5 \node[app] (b\i) at (\x,\y) {}; \fi
        \ifnum\i=6 \node[frag] (b\i) at (\x,\y) {}; \fi % Frag 1
        \ifnum\i=7 \node[free] (b\i) at (\x,\y) {}; \fi
        \ifnum\i=8 \node[frag] (b\i) at (\x,\y) {}; \fi % Frag 2
        \ifnum\i=9 \node[doc] (b\i) at (\x,\y) {}; \fi
        \ifnum\i=10 \node[free] (b\i) at (\x,\y) {}; \fi
        \ifnum\i=11 \node[frag] (b\i) at (\x,\y) {}; \fi % Frag 3
        \ifnum\i=12 \node[app] (b\i) at (\x,\y) {}; \fi
        \ifnum\i=13 \node[free] (b\i) at (\x,\y) {}; \fi
        \ifnum\i=14 \node[frag] (b\i) at (\x,\y) {}; \fi % Frag 4
        \ifnum\i=15 \node[free] (b\i) at (\x,\y) {}; \fi
        \ifnum\i=16 \node[doc] (b\i) at (\x,\y) {}; \fi
        \ifnum\i=17 \node[free] (b\i) at (\x,\y) {}; \fi
        \ifnum\i=18 \node[frag] (b\i) at (\x,\y) {}; \fi % Frag 5
        \ifnum\i=19 \node[free] (b\i) at (\x,\y) {}; \fi
    }
    % Légende Avant
    \begin{scope}[yshift=-2.2cm, xshift=0.5cm]
        \node[sys] at (0,0) {}; \node[right] at (0.4,0) {Système};
        \node[app] at (2,0) {}; \node[right] at (2.4,0) {Applications};
        \node[doc] at (4,0) {}; \node[right] at (4.4,0) {Documents};
        \node[frag] at (6,0) {}; \node[right] at (6.4,0) {Fragments (1 fichier)};
        \node[free] at (8.5,0) {}; \node[right] at (8.9,0) {Espace libre};
    \end{scope}
    
    % APRÈS DÉFRAGMENTATION
    \node[label] at (5, -4.5) {\textbf{APRÈS défragmentation (Disque optimisé)}};
    \foreach \i in {0,...,19} {
        \pgfmathsetmacro{\x}{mod(\i,10)*0.65}
        \pgfmathsetmacro{\y}{-floor(\i/10)*0.8 - 5.5}
        \ifnum\i=0 \node[sys] (a\i) at (\x,\y) {}; \fi
        \ifnum\i=1 \node[sys] (a\i) at (\x,\y) {}; \fi
        \ifnum\i=2 \node[app] (a\i) at (\x,\y) {}; \fi
        \ifnum\i=3 \node[app] (a\i) at (\x,\y) {}; \fi
        \ifnum\i=4 \node[doc] (a\i) at (\x,\y) {}; \fi
        \ifnum\i=5 \node[doc] (a\i) at (\x,\y) {}; \fi
        \ifnum\i=6 \node[frag] (a\i) at (\x,\y) {}; \fi % Fichier regroupé
        \ifnum\i=7 \node[frag] (a\i) at (\x,\y) {}; \fi
        \ifnum\i=8 \node[frag] (a\i) at (\x,\y) {}; \fi
        \ifnum\i=9 \node[frag] (a\i) at (\x,\y) {}; \fi
        \ifnum\i=10 \node[free] (a\i) at (\x,\y) {}; \fi
        \ifnum\i=11 \node[free] (a\i) at (\x,\y) {}; \fi
        \ifnum\i=12 \node[free] (a\i) at (\x,\y) {}; \fi
        \ifnum\i=13 \node[free] (a\i) at (\x,\y) {}; \fi
        \ifnum\i=14 \node[free] (a\i) at (\x,\y) {}; \fi
        \ifnum\i=15 \node[free] (a\i) at (\x,\y) {}; \fi
        \ifnum\i=16 \node[free] (a\i) at (\x,\y) {}; \fi
        \ifnum\i=17 \node[free] (a\i) at (\x,\y) {}; \fi
        \ifnum\i=18 \node[free] (a\i) at (\x,\y) {}; \fi
        \ifnum\i=19 \node[free] (a\i) at (\x,\y) {}; \fi
    }
    % Flèche symbolique
    \draw[poporange, very thick, -{Stealth[length=3mm]}] (5, -2.5) -- (5, -3.5) node[midway, right, font=\small] {Optimisation};
\end{tikzpicture}
\legende{Schéma simplifié de la défragmentation : les fragments dispersés (rose) sont regroupés contiguëment ; l'espace libre (hachuré) est consolidé en fin de disque.}
\end{popfigure}

\begin{aretenir}[Cas des disques SSD]
Sur un \textbf{SSD (Solid State Drive)}, il n'y a pas de pièce mécanique mobile. L'accès à n'importe quelle cellule mémoire prend le même temps. La défragmentation est donc \textbf{inutile} et même \textbf{nuisible} car elle use les cycles d'écriture limités de la mémoire flash. Windows détecte automatiquement les SSD et remplace la défragmentation par l'optimisation \texttt{TRIM} (signalement des blocs inutilisés au contrôleur). Ne lancez jamais manuellement une défragmentation classique sur un SSD.
\end{aretenir}

Sous Windows, l'outil \texttt{Optimiser les lecteurs} (\texttt{dfrgui}) planifie cette tâche automatiquement (hebdomadaire). Vous pouvez l'analyser et la lancer manuellement si le taux de fragmentation dépasse 10\%.

\courssub{Analyse antivirus régulière et mise à jour de la base virale}

Un antivirus sans base de signatures à jour est comme un gardien qui ne connaîtrait pas les visages des voleurs recherchés. La maintenance logicielle impose deux rigueurs :
\begin{enumerate}
    \item \textbf{Mise à jour quotidienne (automatique) :} Vérifiez que l'antivirus (Windows Defender ou solution tierce) télécharge les nouvelles signatures plusieurs fois par jour.
    \item \textbf{Analyse complète planifiée :} Une analyse rapide (dossiers système, démarrage) tous les jours ; une \textbf{analyse complète} (tous les fichiers, tous les disques) au moins une fois par semaine, de préférence hors heures d'activité (nuit, pause déjeuner).
\end{enumerate}

\begin{methode}[Vérifier l'état de Windows Defender (intégré à Windows 10/11)]
\begin{enumerate}
    \item \texttt{Win + I} $\rightarrow$ \textbf{Mise à jour et sécurité} $\rightarrow$ \textbf{Sécurité Windows} $\rightarrow$ \textbf{Ouvrir Sécurité Windows}.
    \item \textbf{Protection contre les virus et menaces} $\rightarrow$ Vérifiez « Aucune action nécessaire » et « Dernière analyse : aujourd'hui ».
    \item Cliquez \textbf{Options d'analyse} $\rightarrow$ Cochez \textbf{Analyse complète} $\rightarrow$ \textbf{Lancer l'analyse}.
\end{enumerate}
\end{methode}

\courssub{Gestion de l'espace disque : surveiller le remplissage}

Un disque système (lecteur \texttt{C:}) rempli à plus de 85-90\% provoque des dysfonctionnements graves : impossibilité d'installer les mises à jour de sécurité, échec de la création de fichiers temporaires (plantages d'applications), ralentissement extrême du système de fichiers (NTFS a besoin d'espace libre pour son journal et la défragmentation), risque de corruption du profil utilisateur.

\begin{exemplebox}[Exemple chiffré : le disque saturé à 95\%]
Imaginons un disque \texttt{C:} de \SI{256}{Go} (typique d'un SSD d'entrée de gamme sur un portable élève).
\begin{itemize}
    \item Espace total : \SI{256}{Go} ($\approx$ \SI{238}{Gio} utilisables).
    \item Espace utilisé : \SI{226}{Gio} (95\%).
    \item Espace libre : \SI{12}{Gio}.
\end{itemize}
\underline{Conséquences concrètes :}
\begin{enumerate}
    \item \textbf{Mise à jour Windows bloquée :} Une mise à jour cumulative nécessite souvent 8 à 15 Gio d'espace temporaire pour télécharger, décompresser et installer. \textbf{Échec garanti.}
    \item \textbf{Fichier d'échange (pagefile.sys) :} Si la RAM (ex: 8 Gio) sature, Windows ne peut pas agrandir le fichier d'échange $\rightarrow$ gel de l'interface, plantage des applications lourdes (navigateur avec 20 onglets, montage vidéo).
    \item \textbf{Défragmentation/TRIM impossible :} L'optimisateur a besoin d'au moins 15-20\% d'espace libre pour déplacer les blocs. Il s'arrête.
    \item \textbf{Usure SSD :} Le contrôleur SSD n'a plus de blocs vierges pour l'usure nivelée (\emph{wear leveling}), réduisant drastiquement la durée de vie du disque.
\end{enumerate}
\textbf{Règle d'or :} Maintenez \textbf{au minimum 20\% d'espace libre} sur le disque système (soit $\ge$ \SI{47}{Gio} sur 238 Gio).
\end{exemplebox}

\begin{popfigure}
\begin{tikzpicture}
    % Pie chart data: System=40, Apps=25, Docs=20, Free=15 (Total 100)
    % Colors: popblue, poporange, popgreen, popline (hatched)
    \def\angleSystem{144}   % 40%
    \def\angleApps{90}      % 25%
    \def\angleDocs{72}      % 20%
    \def\angleFree{54}      % 15%
    
    \begin{scope}[shift={(0,0)}]
        % System (Blue) - Start at 90deg (top)
        \fill[popblue!70] (0,0) -- (90:3) arc (90:90-\angleSystem:3) -- cycle;
        % Apps (Orange)
        \fill[poporange!70] (0,0) -- (90-\angleSystem:3) arc (90-\angleSystem:90-\angleSystem-\angleApps:3) -- cycle;
        % Docs (Green)
        \fill[popgreen!70] (0,0) -- (90-\angleSystem-\angleApps:3) arc (90-\angleSystem-\angleApps:90-\angleSystem-\angleApps-\angleDocs:3) -- cycle;
        % Free (Hatched Grey)
        \fill[popline, pattern=north east lines, pattern color=popmuted] (0,0) -- (90-\angleSystem-\angleApps-\angleDocs:3) arc (90-\angleSystem-\angleApps-\angleDocs:90-\angleSystem-\angleApps-\angleDocs-\angleFree:3) -- cycle;
        
        % Labels with lines
        % System label (mid angle 18deg)
        \draw[popblue, thick] (18:3.2) -- (18:4.2) node[right, align=left, font=\small] {\textbf{Système (40\%)}\\\texttt{Windows, pilotes, pagefile}};
        % Apps label (mid angle -45deg)
        \draw[poporange, thick] (-45:3.2) -- (-45:4.2) node[right, align=left, font=\small] {\textbf{Applications (25\%)}\\\texttt{Office, Navigateurs, Antivirus}};
        % Docs label (mid angle -144deg)
        \draw[popgreen, thick] (-144:3.2) -- (-144:4.5) node[left, align=right, font=\small] {\textbf{Documents (20\%)}\\\texttt{Photos, Vidéos, Travaux}};
        % Free label (mid angle -207deg = 153deg)
        \draw[popmuted, thick] (153:3.2) -- (153:4.5) node[left, align=right, font=\small] {\textbf{Libre (15\%)}\\\texttt{\textit{Seuil d'alerte !}}};
    \end{scope}
    
    % Title
    \node[above, font=\bfseries\large, yshift=1cm] at (0,3) {Répartition typique d'un disque système \texttt{C:} (\SI{256}{Go})};
    % Warning box
    \node[draw=poporange, fill=poporange!10, rounded corners, align=center, font=\small, below, yshift=-4.5cm] at (0,-3) {\textbf{Attention :} 15\% libre est insuffisant.\\Objectif : $\ge$ 20\% (47 Gio) pour santé optimale.};
\end{tikzpicture}
\legende{Diagramme circulaire de l'occupation disque. La zone hachurée (Espace libre) est critique.}
\end{popfigure}

Pour vérifier l'espace : Ouvrez l'\textbf{Explorateur de fichiers} (\texttt{Win + E}) $\rightarrow$ \textbf{Ce PC}. La barre de progression sous le lecteur \texttt{C:} change de couleur (bleu $\rightarrow$ rouge) quand l'espace devient critique. Clic droit $\rightarrow$ \textbf{Propriétés} donne les chiffres exacts (octets / Go).

\courssub{Sauvegarde préventive, redémarrage et fermeture correcte}

La maintenance, c'est aussi anticiper la panne. La \cle{sauvegarde} (backup) n'est pas une option, c'est une assurance-vie pour vos données.

\begin{itemize}
    \item \textbf{Règle 3-2-1 (adaptée élève) :} Gardez vos fichiers importants (devoirs, photos, projets) en \textbf{3 exemplaires} : sur le PC, sur une \textbf{clé USB} ou disque dur externe (déconnecté après copie), et si possible sur un \textbf{cloud} (OneDrive, Google Drive - gratuit pour l'éducation).
    \item \textbf{Fréquence :} Au minimum \textbf{hebdomadaire}, ou après chaque travail important non reproductible (ex: devoir maison finalisé, montage vidéo rendu).
\end{itemize}

Deux habitudes logicielles simples mais puissantes :
\begin{enumerate}
    \item \textbf{Fermer correctement les applications} (Fichier $\rightarrow$ Quitter / \texttt{Ctrl+Q} / \texttt{Alt+F4}) avant d'éteindre. Cela force l'écriture des buffers sur le disque, la sauvegarde des fichiers temporaires et la libération des verrous sur les fichiers. « Tuer » le processus via le Gestionnaire des tâches ou couper l'alimentation brutalement corrompt souvent les fichiers ouverts (base de données, documents Word/Excel non sauvegardés, profil navigateur).
    \item \textbf{Redémarrer l'ordinateur au moins une fois par semaine} (pas seulement « Mettre en veille »). Le redémarrage :
    \begin{itemize}
        \item Vide la RAM (fuites mémoire, processus zombies).
        \item Applique les mises à jour Windows en attente (nécessitent un redémarrage pour finaliser le remplacement des fichiers système verrouillés).
        \item Réinitialise les connexions réseau, les files d'attente d'impression, les services bloqués.
    \end{itemize}
\end{enumerate}

\begin{experience}[TP Guidé : Nettoyer et optimiser son poste Windows (Durée : 20 min)]
\textbf{Objectif :} Appliquer une routine complète de maintenance logicielle de premier niveau.
\textbf{Matériel :} Poste Windows 10/11 (session administrateur ou utilisateur standard avec droits UAC), clé USB pour sauvegarde.

\textbf{Manipulation :}
\begin{enumerate}
    \item \textbf{Sauvegarde :} Copiez votre dossier \texttt{Documents} (ou \texttt{Bureau}) vers votre clé USB. Vérifiez la taille copiée (Propriétés).
    \item \textbf{Fermeture apps :} Fermez tous les programmes ouverts (navigateur, Word, jeux). Vérifiez dans la zone de notification (près de l'horloge) qu'aucun logiciel ne tourne en arrière-plan (Steam, OneDrive, antivirus tiers $\rightarrow$ Quitter si possible).
    \item \textbf{Nettoyage disque :} \texttt{Win + S} $\rightarrow$ tapez \texttt{Nettoyage de disque} $\rightarrow$ Entrée. Sélectionnez \texttt{C:}. Cochez \textbf{Tous les éléments} (Fichiers temporaires, Corbeille, Miniatures, Fichiers d'installation Windows précédents, Nettoyage Windows Update). Cliquez \textbf{Nettoyer les fichiers système} (nécessite admin) pour accéder aux mises à jour supprimées. Validez.
    \item \textbf{Vérification espace :} \texttt{Win + E} $\rightarrow$ Ce PC. Notez l'espace libre sur \texttt{C:} après nettoyage.
    \item \textbf{Défragmentation (si HDD) / Optimisation (si SSD) :} \texttt{Win + S} $\rightarrow$ \texttt{Défragmenter et optimiser les lecteurs}. Sélectionnez \texttt{C:} $\rightarrow$ \textbf{Analyser}. Si HDD et fragmentation $> 10\%$ $\rightarrow$ \textbf{Optimiser}. Si SSD $\rightarrow$ \textbf{Optimiser} (lance TRIM).
    \item \textbf{Analyse Antivirus :} Sécurité Windows $\rightarrow$ Protection virus $\rightarrow$ Options analyse $\rightarrow$ \textbf{Analyse complète}. Lancez-la (peut durer 30-60 min, à faire en fin de séance ou à la maison).
    \item \textbf{Redémarrage :} Menu Démarrer $\rightarrow$ Alimentation $\rightarrow$ \textbf{Redémarrer} (pas Arrêter, pas Veille).
\end{enumerate}

\textbf{Observations attendues :}
\begin{itemize}
    \item Gain d'espace disque significatif (souvent 2 à 10 Go récupérés).
    \item Barre du disque \texttt{C:} repassée en bleu (zone saine).
    \item Ordinateur plus réactif au redémarrage (menu Démarrer, ouverture explorateur).
\end{itemize}

\textbf{Interprétation :} La maintenance logicielle n'est pas une action unique mais un \textbf{cycle} (Nettoyer $\rightarrow$ Protéger $\rightarrow$ Optimiser $\rightarrow$ Sauvegarder $\rightarrow$ Redémarrer). L'automatisation (Nettoyage stockage Windows, Analyses planifiées, OneDrive) en réduit la charge mentale, mais la vérification humaine reste indispensable.
\end{experience}

\begin{exoresolu}[Diagnostic d'un poste ralenti]
\textbf{Énoncé :} Un élève signale que son ordinateur portable (Windows 11, SSD \SI{256}{Go}, RAM \SI{8}{Gio}) est « très lent » depuis deux semaines. Le démarrage prend 2 minutes, l'ouverture du navigateur 30 secondes. Vous ouvrez l'Explorateur de fichiers et voyez la barre du disque \texttt{C:} \textbf{entièrement rouge}. L'espace libre affiché est \textbf{\SI{3,2}{Gio}} sur \SI{237}{Gio}. L'antivirus n'a pas fait d'analyse complète depuis 45 jours. La corbeille contient 120 éléments. Le dossier \texttt{Téléchargements} pèse \SI{45}{Gio} (installateurs anciens, PDF, vidéos).

\textbf{Diagnostic et plan d'action priorisé :}
\begin{enumerate}
    \item \textbf{Cause racine :} \textbf{Saturation critique du disque système (1,3\% libre $<$ 20\% seuil).} Le SSD n'a plus de blocs libres pour le TRIM/wear-leveling $\rightarrow$ écritures lentes $\rightarrow$ gel système. Le fichier d'échange ne peut pas s'agrandir.
    \item \textbf{Action 1 (Immédiate - Gain d'espace massif) :} Vider la Corbeille. Supprimer les installateurs obsolètes dans \texttt{Téléchargements} (garder seulement les derniers pilotes/logiciels utiles). Lancer \texttt{Nettoyage de disque} + \texttt{Nettoyer les fichiers système} (coche \emph{Précédentes installations Windows} si présent). Objectif : remonter $> 50$ Gio libre.
    \item \textbf{Action 2 (Stabilité) :} Redémarrer l'ordinateur après nettoyage pour libérer RAM, appliquer MAJ en attente, reset fichier d'échange.
    \item \textbf{Action 3 (Sécurité) :} Lancer une \textbf{Analyse complète} Windows Defender (45 jours sans analyse = risque élevé).
    \item \textbf{Action 4 (Prévention) :} Activer \textbf{Stockage intelligent} (Paramètres $\rightarrow$ Système $\rightarrow$ Stockage) pour vider automatiquement corbeille et temporaires > 30 jours. Configurer OneDrive pour sauvegarder \texttt{Bureau}, \texttt{Documents}, \texttt{Images}.
    \item \textbf{Vérification :} Après redémarrage, vérifier que l'espace libre $> 20\%$ et que la réactivité est revenue.
\end{enumerate}
\tcblower
\textbf{Solution détaillée :}
Le diagnostic met en évidence la violation de la règle des 20\% d'espace libre. Sur SSD, la saturation dégrade les performances d'écriture de façon exponentielle (le contrôleur doit effacer des blocs avant d'écrire = \emph{write amplification}). La corbeille et le dossier Téléchargements sont des « poubelles » utilisateur faciles à vider. L'antivirus en retard est un risque secondaire mais réel. La priorité absolue est la libération d'espace (Action 1), condition \textit{sine qua non} du retour à la normale. Le redémarrage (Action 2) valide le nettoyage. L'automatisation (Action 4) évite la récidive.
\end{exoresolu}

\begin{aretenir}[Check-list hebdomadaire de maintenance logicielle (Niveau 3ème)]
\begin{itemize}
    \item [$\square$] Vider la Corbeille.
    \item [$\square$] Lancer \texttt{Nettoyage de disque} (fichiers temporaires, miniatures).
    \item [$\square$] Vérifier l'espace libre sur \texttt{C:} ($> 20\%$).
    \item [$\square$] Lancer une analyse antivirus \textbf{complète} (ou vérifier qu'elle a tourné).
    \item [$\square$] Vérifier que Windows Update n'a pas d'erreurs / redémarrage en attente.
    \item [$\square$] Sauvegarder les nouveaux fichiers importants sur clé USB / Cloud.
    \item [$\square$] \textbf{Redémarrer} l'ordinateur (pas veille).
\end{itemize}
\end{aretenir}

\coursec{Installation d'une application simple}

Après avoir vu comment entretenir le système d'exploitation et les utilitaires de base (\emph{maintenance logicielle de premier niveau}), nous abordons maintenant une compétence indispensable : \textbf{installer une nouvelle application} pour répondre à un besoin précis (regarder une vidéo, traiter un texte, programmer, etc.). Installer, ce n'est pas seulement « cliquer sur Suivant » : c'est vérifier la compatibilité, choisir une source sûre, respecter la licence et valider le bon fonctionnement.

\courssub{Qu'est-ce qu'installer une application ?}

\begin{definition}[Installation d'une application]
\textbf{Installer une application}, c'est copier ses fichiers programmes, ses bibliothèques et ses ressources sur le disque dur de l'ordinateur, puis les configurer (clés de registre, raccourcis, associations de fichiers) pour que le système d'exploitation puisse la lancer et l'utiliser correctement.
\end{definition}

Contrairement à un fichier document (texte, image) qu'on ouvre simplement, une application a besoin d'être \emph{intégrée} à l'environnement : création d'un dossier dans \texttt{Program Files} (Windows), inscription dans le menu Démarrer, éventuelle installation de prérequis (ex. \texttt{.NET Framework}, \texttt{Visual C++ Redistributable}). Une installation réussie rend l'application immédiatement accessible et fonctionnelle.

\courssub{Les sources d'installation : d'où vient le programme ?}

On distingue trois sources principales, de la plus sûre à la plus risquée :

\begin{enumerate}
    \item \textbf{Le site officiel de l'éditeur} (ex. \texttt{videolan.org} pour VLC, \texttt{libreoffice.org} pour LibreOffice, \texttt{adobe.com} pour Acrobat Reader). C'est la source de référence : version à jour, sans logiciel indésirable, signature numérique valide.
    \item \textbf{Les magasins d'applications officiels} : \texttt{Microsoft Store} (Windows), \texttt{Mac App Store}, \texttt{Google Play Store} / \texttt{Apple App Store}. Les applications y sont vérifiées, sandboxées et mises à jour automatiquement.
    \item \textbf{Supports physiques} : CD/DVD d'origine, clé USB fournie par l'établissement ou l'entreprise. Fiables si le support n'est pas endommagé, mais la version n'est pas toujours la plus récente.
\end{enumerate}

\begin{attention}[Sources non fiables : danger !]
Ne télécharge \textbf{jamais} une application depuis :
\begin{itemize}
    \item Un site de « téléchargement gratuit » type \texttt{01net}, \texttt{Softonic}, \texttt{Clubic} (souvent des installateurs modifiés qui ajoutent des barres d'outils, des publiciels ou des logiciels espions).
    \item Un lien reçu par WhatsApp, Facebook ou e-mail non sollicité (\emph{phishing}).
    \item Un « crack », « keygen » ou « patch » pour contourner une licence payante : c'est illégal et c'est le vecteur n°1 des ransomwares et chevaux de Troie.
\end{itemize}
\textbf{Règle d'or :} Si l'application est payante, achète-la ou utilise une alternative libre. Si elle est gratuite, télécharge-la \emph{uniquement} sur le site de l'éditeur.
\end{attention}

\courssub{Vérifications préalables : mon ordinateur peut-il l'accueillir ?}

Avant de lancer le moindre \texttt{setup.exe}, vérifie trois points critiques :

\begin{propriete}[Check-list pré-installation]
\begin{enumerate}
    \item \textbf{Compatibilité OS :} L'application existe-t-elle pour \emph{ton} système (Windows 10/11 64 bits, macOS Ventura, Linux Ubuntu 22.04) ? Un \texttt{.exe} Windows ne s'installe pas sur Mac ni sur Linux (sauf couche de compatibilité type Wine, hors programme).
    \item \textbf{Architecture :} 32 bits (x86) ou 64 bits (x64) ? Sur un Windows 64 bits, privilégie la version 64 bits (plus performante). La version 32 bits fonctionne aussi (via WoW64), mais l'inverse est impossible.
    \item \textbf{Configuration minimale :} RAM (ex. 4 Go minimum), espace disque (ex. 2 Go libres), processeur (ex. double cœur 2 GHz), carte graphique (pour les jeux/CAO). Elles sont indiquées sur la page de téléchargement.
    \item \textbf{Espace disque réel :} Vérifie dans l'Explorateur de fichiers (\texttt{Win+E} $\to$ \texttt{Ce PC}) que le disque \texttt{C:} a au moins 20\% d'espace libre \emph{après} l'installation prévue.
    \item \textbf{Droits d'administration :} L'installation écrit dans des dossiers protégés (\texttt{Program Files}, registre \texttt{HKLM}). Tu dois être \emph{administrateur} de la machine ou connaître le mot de passe admin (contrôle de compte utilisateur UAC).
\end{enumerate}
\end{propriete}

\courssub{Étapes types d'une installation sous Windows}

La grande majorité des installateurs Windows suivent le même assistant (wizard). Voici la démarche standard en 6 étapes.

\begin{methode}[Installer une application en 6 étapes]
\begin{enumerate}
    \item \textbf{Télécharger le fichier d'installation} depuis le site officiel. Enregistre-le dans un dossier connu (ex. \texttt{Téléchargements} ou un dossier \texttt{Install} créé à la racine de \texttt{C:}). Le fichier s'appelle souvent \texttt{Setup.exe}, \texttt{Install.exe} ou \texttt{NomAppli\_Version.exe}.
    \item \textbf{Lancer l'installateur} : double-clic sur le fichier. Si l'UAC demande « Voulez-vous autoriser cette application à apporter des modifications à votre appareil ? », clique \textbf{Oui} (si tu fais confiance à la source).
    \item \textbf{Accepter le contrat de licence (CLUF/EULA)} : lis le résumé (ou coche « J'accepte les termes du contrat »). Sans accord, l'installation s'arrête.
    \item \textbf{Choisir le type d'installation} :
    \begin{itemize}
        \item \textbf{Standard / Express / Recommandé :} installe tout par défaut (dossier \texttt{Program Files}, raccourcis, associations). \emph{Choisis cette option si tu es débutant.}
        \item \textbf{Personnalisée / Avancée :} permet de changer le dossier de destination, de ne pas créer de raccourci Bureau, de refuser les composants optionnels (ex. barre d'outils navigateur, télémétrie). \emph{Utile pour éviter les logiciels parasites.}
    \end{itemize}
    \item \textbf{Laisser l'assistant copier les fichiers} (barre de progression). Ne pas éteindre l'ordinateur ni tuer le processus.
    \item \textbf{Terminer et tester} : décoche « Lancer l'application maintenant » si tu veux le faire toi-même. Ouvre le menu Démarrer, tape le nom de l'application, lance-la et teste une fonction de base (ex. ouvrir un fichier, lire une vidéo).
\end{enumerate}
\end{methode}

\begin{popfigure}
\begin{tikzpicture}[
    node distance=1.2cm and 2.5cm,
    startstop/.style={rectangle, rounded corners, minimum width=3cm, minimum height=1cm, text centered, draw=popblue, fill=popblueL, thick},
    process/.style={rectangle, minimum width=3.5cm, minimum height=1cm, text centered, draw=popdark, fill=popcream, thick},
    decision/.style={diamond, aspect=2, minimum width=3cm, minimum height=1.2cm, text centered, draw=poporange, fill=poporangeL, thick},
    arrow/.style={-Stealth, thick, popdark},
    check/.style={rectangle, rounded corners, minimum width=3.5cm, minimum height=1cm, text centered, draw=popgreen, fill=popgreenL, thick}
]
\node (start) [startstop] {Source : Site officiel / Store / CD};
\node (check1) [check, below of=start] {Vérifications : OS, Arch, Config, Espace, Droits admin};
\node (proc1) [process, below of=check1] {Lancer setup.exe (Double-clic $\to$ Oui UAC)};
\node (proc2) [process, below of=proc1] {Accepter licence (CLUF)};
\node (dec1) [decision, below of=proc2] {Type d'installation ?};
\node (proc3a) [process, left of=dec1, xshift=-3.5cm] {Standard / Express (Recommandé)};
\node (proc3b) [process, right of=dec1, xshift=3.5cm] {Personnalisée (Dossier, options, pas de parasites)};
\node (proc4) [process, below of=dec1, yshift=-1.5cm] {Copie des fichiers / Configuration registre / Raccourcis};
\node (check2) [check, below of=proc4] {Test : Lancer l'appli $\to$ Fonction de base OK ?};
\node (end) [startstop, below of=check2] {Installation réussie !};

\draw [arrow] (start) -- (check1);
\draw [arrow] (check1) -- (proc1);
\draw [arrow] (proc1) -- (proc2);
\draw [arrow] (proc2) -- (dec1);
\draw [arrow] (dec1) -- node[above, sloped, font=\small] {Débutant} (proc3a);
\draw [arrow] (dec1) -- node[above, sloped, font=\small] {Avancé} (proc3b);
\draw [arrow] (proc3a) |- (proc4);
\draw [arrow] (proc3b) |- (proc4);
\draw [arrow] (proc4) -- (check2);
\draw [arrow] (check2) -- node[right, font=\small] {Oui} (end);
\draw [arrow] (check2.west) -- ++(-1.5,0) node[left, font=\small] {Non} |- (proc1.west);
\end{tikzpicture}
\legende{Organigramme des étapes d'installation d'une application sous Windows}
\end{popfigure}

\courssub{La notion de licence : ce qu'on a le droit de faire}

\begin{definition}[Types de licences logicielles]
\begin{itemize}
    \item \textbf{Logiciel libre (Open Source) :} Code source accessible, libertés d'utilisation, d'étude, de modification, de redistribution (ex. \textbf{LibreOffice}, \textbf{VLC}, \textbf{Firefox}, \textbf{Linux}, \textbf{Python}). Licences : GPL, MIT, Apache. \emph{Gratuit et légal.}
    \item \textbf{Logiciel gratuit (Freeware / Gratuit) :} Utilisation gratuite, mais code source \emph{fermé}, pas de droit de modification ni redistribution commerciale (ex. \textbf{Adobe Acrobat Reader}, \textbf{7-Zip}, \textbf{TeamViewer} usage personnel). \emph{Gratuit mais propriétaire.}
    \item \textbf{Logiciel payant (Propriétaire / Commercial) :} Nécessite l'achat d'une licence (perpétuelle ou abonnement) pour l'utiliser légalement (ex. \textbf{Microsoft Office}, \textbf{Adobe Photoshop}, \textbf{Windows} hors mise à jour gratuite). \emph{Interdit de cracker.}
    \item \textbf{Shareware / Version d'essai :} Gratuit pour une durée limitée (30 jours) ou avec fonctions limitées. Au-delà, paiement obligatoire.
\end{itemize}
\end{definition}

\begin{savaistu}[Le logiciel libre, une aubaine pour les écoles camerounaises]
Au Cameroun, équiper un laboratoire de 50 postes en licences Microsoft Office + Windows coûte des millions de FCFA par an. De nombreux établissements (lycées, universités, cybercafés éducatifs) migrent vers \textbf{LibreOffice} (traitement de texte, tableur, présentation) et \textbf{Linux Ubuntu} ou \textbf{Linux Mint} : c'est \textbf{gratuit, légal, sécurisé, en français}, et ça tourne sur des machines anciennes (4 Go RAM, 128 Go SSD). Le MINESEC encourage l'usage des logiciels libres dans les programmes pour réduire la fracture numérique et former aux standards ouverts (ODF, HTML, Python).
\end{savaistu}

\courssub{Exemple concret : installer VLC media player}

VLC est le lecteur multimédia de référence : gratuit, libre (GPL), sans publicité, lit \emph{tous} les formats (MKV, MP4, AVI, FLAC, DVD, flux réseau). C'est l'exemple parfait pour un premier TP.

\begin{exemplebox}[Installer VLC depuis videolan.org]
\textbf{Objectif :} Installer VLC 3.0.x (version stable) sur Windows 10/11 64 bits.

\textbf{Étapes détaillées :}
\begin{enumerate}
    \item Ouvre ton navigateur (Edge, Chrome, Firefox). Va sur \texttt{https://www.videolan.org/vlc/}.
    \item Le site détecte ton OS automatiquement. Clique sur le gros bouton \textbf{« Télécharger VLC »} (fichier \texttt{vlc-3.0.x-win64.exe}, ~40 Mo).
    \item Une fois le téléchargement fini (barre en bas du navigateur ou \texttt{Ctrl+J}), clique sur le fichier pour l'ouvrir (ou double-clique dans \texttt{Explorateur > Téléchargements}).
    \item \textbf{UAC :} « Voulez-vous autoriser... VideoLAN ? » $\to$ \textbf{Oui}.
    \item \textbf{Langue :} Français $\to$ OK.
    \item \textbf{Assistant :} Suivant $\to$ J'accepte $\to$ \textbf{Installation standard} (laisse tout coché : raccourcis, associations de fichiers, interface web) $\to$ Installer.
    \item Attends la barre verte. Décoche « Exécuter VLC » $\to$ Terminer.
    \item \textbf{Test :} Menu Démarrer $\to$ tape « VLC » $\to$ Entrée. Glisse-dépose un fichier \texttt{.mp4} ou \texttt{.mkv} dans la fenêtre : la vidéo doit se lire.
\end{enumerate}
\textbf{Résultat :} VLC est installé dans \texttt{C:\textbackslash Program Files\textbackslash VideoLAN\textbackslash VLC}, raccourci dans Menu Démarrer, associations \texttt{.mp4}, \texttt{.mkv}, \texttt{.avi} modifiées pour s'ouvrir avec VLC par défaut.
\end{exemplebox}

\begin{exoresolu}[Installation en ligne de commande (administrateur) pour un parc informatique]
\textbf{Énoncé :} Tu gères une salle de 30 PC identiques (Windows 10 Pro 64 bits). Tu dois installer VLC sur tous les postes sans intervention manuelle (pas de clics « Suivant »). Propose une commande unique à lancer en tant qu'administrateur (via script, stratégie de groupe, ou PSExec).

\textbf{Solution détaillée :}
La plupart des installateurs Windows (Inno Setup, NSIS, MSI) acceptent des paramètres en ligne de commande pour l'installation silencieuse. VLC utilise \textbf{Inno Setup} : les options sont \texttt{/VERYSILENT} (pas d'interface), \texttt{/SUPPRESSMSGBOXES} (pas de boîtes de message), \texttt{/NORESTART} (pas de redémarrage forcé), \texttt{/L=1036} (langue française).
\begin{lstlisting}[language=bash]
vlc-3.0.20-win64.exe /VERYSILENT /SUPPRESSMSGBOXES /NORESTART /L=1036
\end{lstlisting}
\textbf{Explication :}
\begin{itemize}
    \item \texttt{/VERYSILENT} : Installation totalement invisible pour l'utilisateur.
    \item \texttt{/SUPPRESSMSGBOXES} : Masque les éventuelles alertes (ex. fichier en cours d'utilisation).
    \item \texttt{/NORESTART} : Empêche le redémarrage automatique de la machine.
    \item \texttt{/L=1036} : Force la langue française (code LCID 1036).
\end{itemize}
Tu places l'installeur sur un partage réseau (ex. \texttt{\textbackslash\textbackslash SERVEUR\textbackslash Install\textbackslash VLC\textbackslash vlc-3.0.20-win64.exe}) et tu lances la commande à distance via \texttt{psexec} ou une stratégie de groupe (GPO) « Script de démarrage » (ordinateur). Vérification : clé de registre \texttt{HKLM\textbackslash SOFTWARE\textbackslash Microsoft\textbackslash Windows\textbackslash CurrentVersion\textbackslash Uninstall\textbackslash VLC} existe et version = 3.0.20.
\end{exoresolu}

\courssub{Vérification post-installation : ça marche vraiment ?}

L'installation n'est finie que quand tu as \textbf{validé le service rendu}.

\begin{methode}[Check-list post-installation]
\begin{enumerate}
    \item \textbf{Lancement :} Menu Démarrer $\to$ nom de l'appli $\to$ la fenêtre principale s'ouvre sans message d'erreur (DLL manquante, erreur 0xc000007b, etc.).
    \item \textbf{Raccourcis :} Présence dans Menu Démarrer (et Bureau si demandé). Clic droit $\to$ « Ouvrir l'emplacement du fichier » $\to$ cible valide vers \texttt{.exe} dans \texttt{Program Files}.
    \item \textbf{Association de fichiers (si applicable) :} Double-clic sur un fichier géré par l'appli (ex. \texttt{.odt} pour LibreOffice Writer, \texttt{.mp4} pour VLC) $\to$ l'appli s'ouvre et charge le fichier.
    \item \textbf{Fonction de base :} Effectue l'action principale : créer/enregistrer un document, lire une vidéo, ouvrir une page web, compiler un « Hello World ».
    \item \textbf{Mise à jour immédiate :} Menu Aide $\to$ « Vérifier les mises à jour » (ou paramètres $\to$ Mise à jour). Installe la dernière version stable si proposée (correction de failles de sécurité).
\end{enumerate}
\end{methode}

\begin{attention}[Pièges fréquents après installation]
\begin{itemize}
    \item \textbf{« L'icône est là mais rien ne se passe » :} Souvent un antivirus a mis en quarantaine le \texttt{.exe} principal (faux positif). Vérifie l'historique de l'antivirus (Windows Defender : Protection contre les virus $\&$ menaces $\to$ Historique de protection).
    \item \textbf{« Erreur DLL manquante (MSVCP140.dll, VCRUNTIME140.dll) » :} L'application a besoin des \textbf{Visual C++ Redistributables}. Télécharge \texttt{vc\_redist.x64.exe} sur \texttt{learn.microsoft.com} (site officiel Microsoft) et installe-le.
    \item \textbf{Barres d'outils / Moteur de recherche changé :} Tu as laissé coché « Installer la barre XYZ » ou « Définir Yahoo comme page d'accueil » pendant l'installation personnalisée. \emph{Solution :} Désinstalle le parasite via Paramètres $\to$ Applications $\to$ Applications installées, répare les raccourcis navigateur.
    \item \textbf{Application 32 bits sur OS 64 bits mais plantage :} Parfois la version 32 bits bugue sur Win64. Désinstalle, télécharge la version 64 bits (souvent un lien « Other systems » ou « 64-bit » sur la page de téléchargement).
\end{itemize}
\end{attention}

\courssub{Exercice d'application}

\begin{exercice}[Installer et valider LibreOffice Writer]
Tu dois rédiger un rapport de stage au format \texttt{.odt} (OpenDocument). L'ordinateur de la bibliothèque n'a que WordPad.
\begin{enumerate}
    \item Indique l'URL exacte du site officiel pour télécharger LibreOffice (version stable, français, Windows 64 bits).
    \item Liste les 3 vérifications à faire \emph{avant} de lancer l'installeur \texttt{LibreOffice\_24.2.x\_Win\_x64\_fr.msi}.
    \item Pendant l'installation, tu arrives à l'écran « Type d'installation ». Quelle option choisis-tu pour \emph{ne pas} installer le module Base (base de données) dont tu n'as pas besoin, et comment fais-tu ?
    \item Une fois installé, comment vérifies-tu que les fichiers \texttt{.odt} s'ouvrent bien par défaut avec LibreOffice Writer ?
\end{enumerate}
\end{exercice}

\begin{corrige}[Exercice n°1]
\begin{enumerate}
    \item \texttt{https://fr.libreoffice.org/download/download-libreoffice/} (choisir « Windows (64-bit) » $\to$ « Français » $\to$ « Télécharger »).
    \item (1) OS = Windows 10/11 64 bits (compatibilité). (2) Espace disque : $\approx$ 1,5 Go requis + marge $\to$ vérifier \texttt{C:} libre. (3) Droits admin : être connecté en admin ou connaître le mot de passe (UAC).
    \item Choisir \textbf{Personnalisée}. Dans l'arborescence des composants, déplier « Base de données » (ou « Base »), cliquer sur l'icône disque $\to$ « Cette fonctionnalité ne sera pas disponible ». Laisser Writer, Calc, Impress, Draw, Math cochés. Continuer.
    \item Double-cliquer sur un fichier \texttt{.odt} existant (ou en créer un vide via Clic droit $\to$ Nouveau $\to$ Document OpenDocument Text). Il doit s'ouvrir dans Writer. Sinon : Clic droit sur le \texttt{.odt} $\to$ Ouvrir avec $\to$ Choisir une autre application $\to$ LibreOffice Writer $\to$ cocher « Toujours utiliser cette application ».
\end{enumerate}
\end{corrige}

\courssub{Transition vers la mise à jour}

Tu sais maintenant installer une application proprement, depuis une source sûre, en respectant la licence et en validant le résultat. Mais un logiciel installé \emph{aujourd'hui} peut devenir une faille de sécurité \emph{demain} si on ne le met pas à jour. La section suivante — \textbf{Mise à jour des applications et démarche de diagnostic} — t'apprendra à maintenir tes applications à jour (manuellement, automatiquement, en ligne de commande) et à diagnostiquer une installation qui ne fonctionne pas.

\coursec{Mise à jour des applications et démarche de diagnostic}

Après avoir vu comment installer une application proprement, il est indispensable de comprendre que la vie d'un logiciel ne s'arrête pas à son installation. Un logiciel, comme un véhicule, nécessite un entretien régulier pour rester performant, compatible avec les autres outils et, surtout, sécurisé. Dans cette section, nous abordons les mises à jour — correctifs vitaux — et la méthode rigoureuse pour diagnostiquer une panne logicielle ou matérielle de premier niveau.

\courssub{Qu'est-ce qu'une mise à jour et pourquoi est-elle cruciale ?}

\begin{definition}[Mise à jour logicielle (Update)]
Une \cle{mise à jour} est une nouvelle version d'un logiciel ou d'un système d'exploitation, distribuée par son éditeur, qui corrige des erreurs de programmation (\emph{bogues}), comble des failles de sécurité, améliore les performances ou ajoute de nouvelles fonctionnalités. Elle se présente souvent sous la forme d'un fichier exécutable ou d'un paquet téléchargé automatiquement.
\end{definition}

On distingue généralement trois types de mises à jour :
\begin{itemize}
    \item \textbf{Correctives (Patchs) :} Corrigent un dysfonctionnement précis (ex. : un plantage lors de l'impression).
    \item \textbf{De sécurité (Critical Updates) :} Colmatent une faille exploitable par un virus ou un pirate. C'est le type le plus urgent.
    \item \textbf{Évolutives (Feature Updates) :} Apportent de nouvelles fonctions ou une nouvelle interface (ex. : passage de Windows 10 à Windows 11, ou mise à jour majeure d'Android).
\end{itemize}

\begin{propriete}[Rôles fondamentaux des mises à jour]
Les mises à jour assurent trois piliers de la santé informatique :
\begin{enumerate}
    \item \textbf{Sécurité :} Elles ferment les portes d'entrée des logiciels malveillants (ransomwares, chevaux de Troie). Un logiciel non mis à jour est une cible facile.
    \item \textbf{Stabilité et performances :} Elles corrigent les fuites de mémoire, les conflits de pilotes et optimisent l'utilisation du processeur et de la RAM.
    \item \textbf{Compatibilité :} Elles permettent au logiciel de fonctionner avec les nouveaux formats de fichiers, les nouveaux périphériques (imprimantes, clés 4G) et les nouvelles versions du système d'exploitation.
\end{enumerate}
\end{propriete}

\begin{attention}[Logiciels obsolètes : la porte ouverte aux virus]
\emph{Des logiciels non mis à jour constituent la principale porte d'entrée des virus et des rançongiciels.} L'attaque mondiale \textbf{WannaCry} (mai 2017) a paralysé des centaines de milliers d'ordinateurs (hôpitaux, entreprises, administrations) en exploitant une faille de Windows (\texttt{MS17-010}) pour laquelle un correctif existait depuis deux mois. Les victimes étaient celles qui n'avaient pas appliqué la mise à jour. Au Cameroun, de nombreux postes de cybercafés ou d'écoles tournent encore sur d'anciennes versions de Windows 7 ou 10 non patchées : c'est un risque majeur.
\end{attention}

\courssub{Comment mettre à jour une application ?}

La procédure dépend du mode de distribution du logiciel. Voici les deux cas les plus fréquents en environnement scolaire ou professionnel.

\begin{methode}[Mise à jour via le menu intégré de l'application (Méthode recommandée)]
La plupart des logiciels modernes (navigateurs, suites bureautiques, antivirus, VLC, 7-Zip) intègrent un module de vérification automatique.
\begin{enumerate}
    \item Ouvrir l'application.
    \item Aller dans le menu \textbf{Aide} (ou \texttt{Help}) puis cliquer sur \textbf{Rechercher les mises à jour} / \texttt{Check for updates} / \textbf{À propos de...}.
    \item Si une mise à jour est trouvée : cliquer sur \textbf{Télécharger} puis \textbf{Installer}.
    \item Accepter les droits administrateur (UAC) si demandés.
    \item Redémarrer l'application (ou l'ordinateur si requis) pour finaliser.
\end{enumerate}
\end{methode}

\begin{methode}[Mise à jour par téléchargement manuel (Site officiel uniquement)]
Nécessaire pour les logiciels sans mise à jour automatique ou si la mise à jour automatique a échoué.
\begin{enumerate}
    \item Identifier la version installée : Menu \texttt{Aide} $\to$ \texttt{À propos de} (noter le numéro de version, ex. \texttt{3.4.1}).
    \item Se rendre \textbf{uniquement} sur le site officiel de l'éditeur (ex. \texttt{www.videolan.org} pour VLC, \texttt{www.7-zip.org} pour 7-Zip). \textbf{Éviter les sites de téléchargement tiers (Softonic, Clubic, 01net) qui injectent souvent des publiciels (adware).}
    \item Télécharger la version \textbf{stable} (pas \emph{Beta} ni \emph{Nightly}) correspondant à l'architecture du processeur (\texttt{64-bit} / \texttt{x64} pour les PC récents, \texttt{32-bit} / \texttt{x86} pour les anciens).
    \item Fermer l'ancienne version de l'application.
    \item Lancer l'installateur téléchargé : il détecte souvent l'installation existante et propose une mise à jour \emph{par-dessus} (conservant les préférences), sinon désinstaller l'ancienne version d'abord (via \texttt{Paramètres} $\to$ \texttt{Applications}).
\end{enumerate}
\end{methode}

\begin{exemplebox}[Mise à jour de l'antivirus Avast Free avant analyse d'une clé USB]
\textbf{Contexte :} Tu insères une clé USB venue d'un cybercafé à Yaoundé pour récupérer un devoir. Avant d'ouvrir le moindre fichier, tu dois t'assurer que ton antivirus connaît les toutes dernières menaces.

\begin{lstlisting}[language=text, caption=Étapes de mise à jour Avast (interface graphique)]
1. Clic droit sur l'icône Avast (zone de notification, près de l'horloge)
   > "Ouvrir l'interface utilisateur Avast"
2. Menu "Menu" (trois barres en haut à droite) > "Paramètres"
3. Onglet "Général" > "Mise à jour"
4. Section "Définitions de virus" > Cliquer sur "Mettre à jour"
   // Attendre : "Les définitions de virus sont à jour"
5. Section "Programme" > Cliquer sur "Vérifier les mises à jour"
   // Si dispo : "Mettre à jour" > Redémarrer si demandé
6. Retour écran principal > "Protection" > "Analyses antivirus"
   > "Analyse personnalisée" > Sélectionner la lettre de la clé USB (ex. E:)
   > "Démarrer l'analyse"
\end{lstlisting}

\textbf{Règle d'or :} \textbf{Jamais d'analyse antivirus sans mise à jour préalable des définitions de virus.} Une base de signatures vieille de 15 jours ne détectera pas le virus du jour.
\end{exemplebox}

\courssub{La mise à jour du système d'exploitation : Windows Update}

Le système d'exploitation (Windows 10/11 en milieu scolaire camerounais) est la couche la plus critique. Sa mise à jour gère le noyau, les pilotes matériels (carte graphique, Wi-Fi, son, chipset) et les composants de sécurité (Windows Defender, BitLocker, Pare-feu).

\begin{methode}[Lancer et gérer Windows Update (Windows 10/11)]
\begin{enumerate}
    \item \texttt{Touche Windows + I} (Paramètres) $\to$ \textbf{Mise à jour et sécurité} (Win10) / \textbf{Windows Update} (Win11).
    \item Cliquer sur \textbf{Rechercher les mises à jour}.
    \item Laisser télécharger et installer. \textbf{Ne pas éteindre l'ordinateur.}
    \item Si le bouton \textbf{Redémarrer maintenant} apparaît : enregistrer ton travail et redémarrer.
    \item Après redémarrage, revenir dans Windows Update et cliquer à nouveau sur \textbf{Rechercher les mises à jour} (parfois des mises à jour « facultatives » ou des pilotes apparaissent seulement après le premier redémarrage).
    \item Vérifier \texttt{Options avancées} $\to$ \texttt{Mises à jour facultatives} pour installer les pilotes spécifiques (ex. pilote carte graphique NVIDIA/AMD, pilote Wi-Fi Realtek).
\end{enumerate}
\end{methode}

\begin{attention}[Coupure de courant pendant mise à jour système = Risque de « briquer » le PC]
\emph{Interrompre une mise à jour Windows en cours (coupure de courant, batterie vide, appui long sur le bouton Power) peut corrompre les fichiers système critiques (Registre, \texttt{winload.efi}, \texttt{ntoskrnl.exe}).} L'ordinateur peut ne plus démarrer (écran bleu \texttt{INACCESSIBLE\_BOOT\_DEVICE} ou boucle de réparation automatique).
\begin{itemize}
    \item \textbf{Sur secteur :} Brancher sur un \textbf{onduleur} (UPS) si possible, ou au minimum une multiprise parafoudre.
    \item \textbf{Sur portable :} Vérifier que la batterie est chargée > 20\% \textbf{ET} laisser le chargeur branché pendant toute l'opération.
    \item \textbf{Sauvegarde préalable :} Avant une mise à jour majeure (Feature Update type 22H2 $\to$ 23H2), faire une sauvegarde des documents importants sur clé USB ou Google Drive / OneDrive.
\end{itemize}
\end{attention}

\begin{savaistu}[WannaCry et le Cameroun]
En mai 2017, le rançongiciel \textbf{WannaCry} a infecté plus de 200 000 ordinateurs dans 150 pays en quelques heures. Il exploitait la faille \texttt{EternalBlue} (port SMB 445) dans Windows. Microsoft avait publié le correctif (\texttt{MS17-010}) en \textbf{mars 2017}, soit deux mois avant l'attaque. Au Cameroun, plusieurs administrations, banques et établissements scolaires ont été touchés car leurs parcs informatiques n'étaient pas configurés pour les mises à jour automatiques. Cette attaque a coûté des milliards de FCFA en pertes de données et temps de réparation. \textbf{La leçon : activer les mises à jour automatiques, c'est gratuit et ça protège.}
\end{savaistu}

\courssub{La démarche de diagnostic : une méthode scientifique pour le dépannage}

Face à un ordinateur qui « ne marche pas » (lent, écran bleu, imprimante muette, Wi-Fi absent), l'amateur devine et essaie au hasard. Le technicien de premier niveau applique une \cle{démarche de diagnostic} structurée, traçable et justifiée. C'est une compétence évaluée au BEPC : « Rédiger et justifier une démarche, une réponse ou une conclusion ».

\begin{methode}[Rédiger une démarche de diagnostic en 5 étapes (Standard professionnel)]
Cette méthode s'applique à tout incident matériel ou logiciel de premier niveau.

\begin{enumerate}
    \item \textbf{Description précise du problème (Constat) :} Qu'observe-t-on exactement ? Message d'erreur (copier le code), comportement anormal, contexte (quand, quoi, quoi fait juste avant). \emph{Ex. : « Au démarrage, écran bleu \texttt{CRITICAL\_PROCESS\_DIED} après logo Windows. Survenu après installation pilote imprimante HP. »}
    \item \textbf{Formulation d'hypothèses (Ordre de probabilité) :} Lister les causes possibles, de la plus simple/fréquente à la plus complexe. \emph{Ex. : H1: Conflit pilote imprimante (récent). H2: Corruption fichiers système. H3: Disque dur défaillant. H4: RAM défectueuse.}
    \item \textbf{Vérifications / Tests (Élimination) :} Appliquer des tests ciblés pour valider ou invalider chaque hypothèse. Noter le résultat de chaque test. \emph{Ex. : Test 1: Démarrage en Mode Sans Échec $\to$ OK $\Rightarrow$ H1 probable. Test 2: Désinstaller pilote imprimante en Mode Sans Échec $\to$ Redémarrage normal $\Rightarrow$ H1 confirmée.}
    \item \textbf{Solution mise en œuvre :} Action corrective définitive. \emph{Ex. : Télécharger pilote latest version sur site HP, installer, redémarrer.}
    \item \textbf{Conclusion et validation (Non-régression) :} Vérifier que le problème initial est résolu ET qu'aucun nouveau problème n'est apparu. Documenter la solution pour mémoire future. \emph{Ex. : « PC redémarre normalement 3 fois de suite. Imprimante imprime page de test. Problème résolu. »}
\end{enumerate}
\end{methode}

\begin{popfigure}
\begin{tikzpicture}[
    node distance=1.2cm and 2.5cm,
    box/.style={rectangle, draw=popdark, thick, rounded corners=6pt, minimum width=3.2cm, minimum height=1.4cm, text centered, font=\small, fill=popcream},
    arrow/.style={-Stealth, thick, popblue},
    startstop/.style={ellipse, draw=popgreen, thick, minimum width=3cm, minimum height=1.3cm, text centered, font=\small\bfseries, fill=popgreenL},
    decision/.style={diamond, draw=poporange, thick, aspect=1.5, minimum width=3.5cm, minimum height=1.5cm, text centered, font=\small, fill=poporangeL, inner sep=2pt}
]
% Nœuds
\node (start) [startstop, align=center]{DÉBUT\\Problème signalé};
\node (etape1) [box, below of=start, align=center]{1. DESCRIPTION\\Précise du problème\\(Symptômes, code erreur, contexte)};
\node (etape2) [box, below of=etape1, align=center]{2. HYPOTHÈSES\\Classées par probabilité\\ (Simple $\to$ Complexe)};
\node (etape3) [box, below of=etape2, align=center]{3. TESTS / VÉRIFICATIONS\\Élimination systématique\\ (Mode sans échec, logs, swap matériel)};
\node (decision) [decision, below of=etape3, yshift=-0.5cm, align=center]{Hypothèse\\validée ?};
\node (etape4) [box, below of=decision, yshift=-0.5cm, align=center]{4. SOLUTION\\Action corrective\\ (Réparation, MAJ, Remplacement)};
\node (etape5) [box, below of=etape4, align=center]{5. VALIDATION\\Test non-régression\\ + Documentation};
\node (end) [startstop, below of=etape5, align=center]{FIN\\Incident clos};

% Flèches principales
\draw [arrow] (start) -- (etape1);
\draw [arrow] (etape1) -- (etape2);
\draw [arrow] (etape2) -- (etape3);
\draw [arrow] (etape3) -- (decision);

% Branche OUI
\draw [arrow] (decision) -- node[left, font=\small, text=popgreen] {OUI} (etape4);
\draw [arrow] (etape4) -- (etape5);
\draw [arrow] (etape5) -- (end);

% Branche NON (retour)
\draw [arrow] (decision.east) -- ++(1.5,0) node[above, font=\small, text=poporange] {NON} |- (etape2.east);
\draw [arrow] (decision.west) -- ++(-1.5,0) node[above, font=\small, text=poporange] {NON} |- (etape2.west);

% Annotations
\node[align=center, font=\tiny, text=popmuted] at ($(etape1.east)+(1.8,0)$) {Observation\\Factuelle};
\node[align=center, font=\tiny, text=popmuted] at ($(etape2.east)+(1.8,0)$) {Raisonnement\\Déductif};
\node[align=center, font=\tiny, text=popmuted] at ($(etape3.east)+(1.8,0)$) {Expérimentation\\Contrôlée};
\node[align=center, font=\tiny, text=popmuted] at ($(etape4.east)+(1.8,0)$) {Action\\Corrective};
\node[align=center, font=\tiny, text=popmuted] at ($(etape5.east)+(1.8,0)$) {Preuve\\Qualité};
\end{tikzpicture}
\legende{Cycle de la démarche de diagnostic en 5 étapes (norme ITIL simplifiée pour le premier niveau). La boucle « NON » sur les tests est la clé : on ne devine pas, on élimine.}
\end{popfigure}

\courssub{Exemple résolu : Diagnostic d'un PC « très lent » au lycée}

\begin{exoresolu}[Diagnostic lenteur poste salle informatique]
\textbf{Énoncé :} Un poste de la salle informatique (Windows 10, 4 Go RAM, HDD 500 Go) devient anormalement lent : ouverture de Word > 45 s, gel souris fréquent, disque dur à 100\% d'activité dans le Gestionnaire des tâches. L'élève n'a rien installé de nouveau. Proposer une démarche de diagnostic complète et la solution.

\tcblower
\textbf{Solution détaillée :}

\textbf{1. Description du problème :}
Symptômes : Latence extrême interface utilisateur, disque système (C:) à 100\% d'activité continue (colonne « Disque » Gestionnaire des tâches), RAM utilisée 3,8/4 Go. Contexte : Poste partagé, pas d'installation récente connue. Pas de message d'erreur explicite.

\textbf{2. Hypothèses (classées) :}
\begin{itemize}
    \item H1 (Probable) : Processus système « Superfetch / SysMain » ou « Windows Update » ou « Antimalware Service Executable » saturant le disque HDD lent.
    \item H2 (Possible) : Infection malicielle (miner crypto, botnet) masquée.
    \item H3 (Possible) : Disque dur en fin de vie (secteurs défectueux $\to$ réessais lecture).
    \item H4 (Moins probable) : RAM insuffisante (swapping excessif) ou barrette défectueuse.
\end{itemize}

\textbf{3. Vérifications / Tests :}
\begin{center}
\begin{tabularx}{\linewidth}{|l|X|X|}
\hline
\textbf{Test} & \textbf{Action} & \textbf{Résultat / Interprétation} \\
\hline
T1 & Ouvrir Gestionnaire des tâches (Ctrl+Maj+Echap) $\to$ Onglet Détail $\to$ Trier par « Disque » & Processus \texttt{Antimalware Service Executable} (MsMpEng.exe) à 2-3 Mo/s constant. \textbf{H1 validée (partiellement)}. \\
\hline
T2 & Vérifier Windows Update : Paramètres $\to$ Mise à jour & Téléchargement en arrière-plan d'une mise à jour cumulative (2 Go). \textbf{H1 confirmée : cause principale identifiée}. \\
\hline
T3 (si T2 négatif) & Lancer analyse rapide Windows Defender & Propre. \textbf{H2 invalidée}. \\
\hline
T4 (si T2 négatif) & \texttt{chkdsk C: /f /r} au redémarrage & Aucune erreur. \textbf{H3 invalidée}. \\
\hline
\end{tabularx}
\end{center}

\textbf{4. Solution mise en œuvre :}
\begin{enumerate}
    \item Laisser Windows Update terminer le téléchargement et l'installation (redémarrage requis).
    \item \textbf{Optimisation post-MAJ :} Désactiver \texttt{SysMain} (ancien Superfetch) inadapté aux HDD : \texttt{services.msc} $\to$ SysMain $\to$ Type de démarrage : \texttt{Désactivé} $\to$ Arrêter.
    \item Configurer les « Heures d'activité » dans Windows Update pour éviter les téléchargements en cours de cours (ex. 07h-18h).
    \item \textbf{Recommandation matériel :} Passer à 8 Go RAM + SSD 240 Go (coût modique, gain massif).
\end{enumerate}

\textbf{5. Conclusion et validation :}
Après redémarrage post-MAJ et désactivation SysMain : Disque < 5\% au repos, ouverture Word < 5 s. Problème résolu. Fiche d'intervention remplie : « Lenteur due MAJ Windows + SysMain sur HDD. MAJ appliquée, service désactivé. Préconisation upgrade SSD/RAM. »
\end{exoresolu}

\courssub{Calendrier de maintenance préventive : l'antidote aux pannes}

La maintenance de premier niveau ne se fait pas « quand ça casse », mais \textbf{avant}. Un calendrier préventif simple, affiché dans la salle informatique ou noté dans le cahier de maintenance du responsable informatique de l'établissement, garantit la longévité du parc.

\begin{popfigure}
\begin{tikzpicture}
\begin{scope}[
    every node/.style={font=\small},
    header/.style={fill=popblue, text=white, font=\small\bfseries, minimum height=0.9cm, align=center},
    rowodd/.style={fill=popblueL, minimum height=0.85cm, align=center},
    roweven/.style={fill=white, minimum height=0.85cm, align=center},
    rowhead/.style={fill=poppurple, text=white, font=\small\bfseries, minimum height=0.85cm, align=center},
    cell/.style={minimum height=0.85cm, align=left, draw=popline, thin},
    cellc/.style={minimum height=0.85cm, align=center, draw=popline, thin},
]
% En-têtes
\node[header, minimum width=3.2cm] (h1) at (0,0) {Fréquence};
\node[header, minimum width=5.5cm] (h2) at (3.2,0) {Opérations Logicielles};
\node[header, minimum width=5.5cm] (h3) at (8.7,0) {Opérations Matérielles / Environnement};
% Lignes
% Hebdomadaire
\node[rowodd, minimum width=3.2cm, align=center] (f1) at (0,-0.9){HEBDOMADAIRE\\(Chaque vendredi)};
\node[rowodd, cell, minimum width=5.5cm, text width=5.1cm, align=left] (l1) at (3.2,-0.9) {
\begin{itemize}\item Windows Update : Rechercher \& installer
\item Antivirus : MAJ définitions + Analyse rapide
\item Vider Corbeille + Nettoyage disque (fichiers temporaires)
\item Vérifier sauvegardes élèves (OneDrive/Clé USB)
\end{itemize}
};
\node[rowodd, cell, minimum width=5.5cm, text width=5.1cm, align=left] (m1) at (8.7,-0.9) {
\begin{itemize}\item Dépoussiérage écrans, claviers, souris (chiffon microfibre)
\item Vérifier câblage (alimentation, RJ45) bien enclenché
\item Ranger câbles, souris, casques aux places
\end{itemize}
};
% Mensuel
\node[roweven, minimum width=3.2cm, align=center] (f2) at (0,-2.3){MENSUEL\\(1er du mois)};
\node[roweven, cell, minimum width=5.5cm, text width=5.1cm, align=left] (l2) at (3.2,-2.3) {
\begin{itemize}\item Windows Update : Vérifier « Mises à jour facultatives » (Pilotes)
\item MAJ navigateurs (Chrome/Firefox/Edge), VLC, 7-Zip, PDF Reader
\item Désinstaller logiciels inutiles / barres d'outils (Panneau config)
\item Vérifier espace disque C: > 20\% libre
\item Vérifier points de restauration système actifs
\end{itemize}
};
\node[roweven, cell, minimum width=5.5cm, text width=5.1cm, align=left] (m2) at (8.7,-2.3) {
\begin{itemize}\item Nettoyage ventilateurs / grilles boîtier (bombe air sec \textbf{PC éteint \& débranché})
\item Test démarrage / extinction normal
\item Vérifier état piles souris/clavier sans fil
\item Test imprimante (page de test) si locale
\end{itemize}
};
% Semestriel
\node[rowhead, minimum width=3.2cm, align=center] (f3) at (0,-4.1){SEMESTRIEL\\(Vacances grandes / petites)};
\node[rowhead, cell, minimum width=5.5cm, text width=5.1cm, align=left] (l3) at (3.2,-4.1) {
\begin{itemize}\item \textbf{Nettoyage approfondi :} CCleaner (registre + fichiers) ou scripts PowerShell
\item \textbf{Analyse antivirus COMPLÈTE} (hors ligne si possible : Windows Defender Offline)
\item Vérification intégrité fichiers système : \texttt{sfc /scannow} + \texttt{DISM /Online /Cleanup-Image /RestoreHealth} (Admin)
\item Revue comptes utilisateurs : supprimer comptes élèves sortants, réinitialiser mots de passe
\item Sauvegarde image système (Macrium Reflect Free / Clonezilla) sur disque externe
\end{itemize}
};
\node[rowhead, cell, minimum width=5.5cm, text width=5.1cm, align=left] (m3) at (8.7,-4.1) {
\begin{itemize}\item \textbf{Ouverture boîtier :} Soufflage complet (alimentation, CPU, GPU, RAM, carte mère) \textbf{PC débranché, bracelet antistatique}
\item Remplacement pâte thermique CPU si T° > 75°C en charge (HWMonitor)
\item Vérification visuelle condensateurs (bombés = HS)
\item Test alimentation (testeur ATX ou multimètre) ou swap
\item Vérification onduleurs (test autonomie batterie)
\item Inventaire matériel mis à jour (numéro de série, specs, état)
\end{itemize}
};
% Grille
\draw[popline, thin] (h1.north west) rectangle (h3.north east);
\draw[popline, thin] (f1.south west) rectangle (m3.south east);
% Lignes verticales
\draw[popline, thin] (h1.north east) -- (m3.south east);
\draw[popline, thin] (h2.north east) -- (m3.south east);
% Lignes horizontales
\foreach \y in {-0.9, -2.3, -4.1} {
    \draw[popline, thin] (-1.6, \y) -- (11.45, \y);
}
\end{scope}
\end{tikzpicture}
\legende{Calendrier type de maintenance préventive pour un parc scolaire (adaptable). La régularité hebdomadaire (MAJ + Antivirus) est le minimum vital. Le semestriel nécessite l'arrêt des cours (vacances).}
\end{popfigure}

\begin{aretenir}[Synthèse : Les 3 piliers de la maintenance de premier niveau]
\begin{enumerate}
    \item \textbf{Mises à jour systématiques :} OS (Windows Update), Antivirus (Définitions + Moteur), Applications (Menu Aide ou Site officiel). \textbf{Jamais sans sauvegarde ni alimentation stable pour l'OS.}
    \item \textbf{Diagnostic structuré :} Constat $\to$ Hypothèses $\to$ Tests $\to$ Solution $\to$ Validation. On écrit ce qu'on fait, ce qu'on voit, ce qu'on conclut. C'est la trace professionnelle.
    \item \textbf{Prévention planifiée :} Hebdo (MAJ, AV, Nettoyage), Mensuel (Pilotes, Logiciels, Air sec), Semestriel (Ouverture boîtier, Image système, Inventaire). Un calendrier affiché = des pannes évitées.
\end{enumerate}
\end{aretenir}

\begin{experience}[TP Guidé : Simuler une démarche de diagnostic complète]
\textbf{Objectif :} Appliquer la méthode en 5 étapes sur un scénario réaliste.
\textbf{Matériel :} Poste Windows (VM ou physique), accès administrateur, clé USB test.

\textbf{Scénario :} « L'ordinateur de la bibliothécaire ne peut plus imprimer sur l'imprimante réseau HP LaserJet Pro M404dn. Le message « Imprimante hors ligne » s'affiche. Elle a besoin d'imprimer les listes de classe pour demain 8h. »

\textbf{Manipulation :}
\begin{enumerate}
    \item \textbf{Constat (Étape 1) :} Reproduire le défaut. Noter l'adresse IP de l'imprimante (page de configuration imprimée depuis le panneau de l'imprimante). Noter l'IP du PC (\texttt{ipconfig}). Pinguer l'IP imprimante depuis le PC.
    \item \textbf{Hypothèses (Étape 2) :} Écrire 3 hypothèses classées (ex. : H1: IP imprimante changée par DHCP / H2: Service « Spouleur d'impression » arrêté / H3: Pare-feu bloque port 9100).
    \item \textbf{Tests (Étape 3) :} Exécuter : \texttt{services.msc} (vérifier Spouleur), \texttt{ping <IP\_imprimante>}, \texttt{telnet <IP\_imprimante> 9100} (si activé), vérifier file d'attente imprimante (clic droit $\to$ Voir ce qui s'imprime $\to$ Imprimante $\to$ Décocher « Utiliser l'imprimante hors ligne »).
    \item \textbf{Solution (Étape 4) :} Appliquer le correctif (ex. : Redémarrer service Spouleur, ou réinstaller port TCP/IP standard avec IP fixe, ou réserver bail DHCP sur routeur).
    \item \textbf{Validation (Étape 5) :} Imprimer page de test Windows + document Word réel. Vérifier que d'autres postes impriment encore. Rédiger la fiche d'intervention (Date, Heure, Poste, Problème, Solution, Signature).
\end{enumerate}

\textbf{Observations attendues :} L'élève doit produire un compte-rendu écrit structuré selon les 5 étapes. L'évaluation porte sur la clarté du constat, la logique des hypothèses, la pertinence des tests (pas de « tout réinstaller » d'emblée), l'efficacité de la solution et la preuve de validation.
\end{experience}

\begin{exercice}[Entraînement BEPC - Maintenance \& Diagnostic]
\begin{enumerate}
    \item \textbf{QCM :} Quelle action est \textbf{la plus risquée} pendant l'installation d'une mise à jour majeure de Windows (Feature Update) ?
    \begin{itemize}
        \item[A] Laisser l'ordinateur portable sur batterie à 80\% sans chargeur.
        \item[B] Éteindre l'écran pour économiser l'énergie.
        \item[C] Couper l'alimentation secteur (ou batterie vide) pour forcer le redémarrage car « ça prend trop de temps ».
        \item[D] Lancer une analyse antivirus en parallèle.
    \end{itemize}
    \item \textbf{Mise en situation :} Tu insères une clé USB d'un camarade. Windows Defender affiche : « Menace détectée : Trojan:Script/Wacatac.B!ml ». L'antivirus n'a pas été mis à jour depuis 3 semaines.
    \begin{enumerate}
        \item Quelle est la \textbf{première action} à faire avant de cliquer sur « Supprimer » ou « Mettre en quarantaine » ? Justifie.
        \item Pourquoi l'analyse a-t-elle pu détecter ce virus malgré des définitions anciennes ? (Indice : heuristique / cloud).
    \end{enumerate}
    \item \textbf{Rédaction (Compétence « Justifier une démarche ») :} Un poste de la salle info affiche un écran bleu \texttt{VIDEO\_TDR\_FAILURE} (nvlddmkm.sys) au lancement d'un jeu éducatif 3D. Le pilote graphique NVIDIA date de 2020. Rédige la démarche de diagnostic complète en 5 étapes (Constat, Hypothèses, Tests, Solution, Validation) pour résoudre ce problème durablement.
\end{enumerate}
\end{exercice}

\begin{corrige}[Exercice n°1]
\begin{enumerate}
    \item \textbf{Bonne réponse : C.} Couper l'alimentation pendant une mise à jour système corrompt les fichiers critiques en cours d'écriture (Registre, BCD, fichiers noyau), rendant le système inutilisable (bootloop, écran bleu). C'est la cause n°1 de « briquage » logiciel évitable. A est déconseillé mais moins fatal (risque extinction si batterie finit). B est sans conséquence. D ralentit mais ne casse pas.
    \item \begin{enumerate}
        \item \textbf{Mettre à jour les définitions de virus (et le moteur si dispo) AVANT toute action de nettoyage.} Justification : Les définitions de 3 semaines ne contiennent pas la signature exacte de cette variante précise. La détection actuelle est probablement « heuristique » (comportement suspect) ou « cloud » (réputation), moins fiable qu'une signature à jour. La mise à jour garantit la suppression complète et la détection d'éventuels compagnons (droppers, persistence).
        \item Windows Defender utilise la \textbf{protection cloud (MAPS - Microsoft Active Protection Service)} et l'\textbf{analyse heuristique/comportementale} (surveillance API suspectes, scripts obfusqués, injection mémoire). Cela permet de bloquer des menaces « zero-day » non encore dans la base de signatures locale, mais la base locale reste indispensable pour la désinfection précise et hors-ligne.
    \end{enumerate}
    \item \textbf{Proposition de démarche :}
    \textbf{1. Constat :} Écran bleu \texttt{VIDEO\_TDR\_FAILURE} (module \texttt{nvlddmkm.sys}) au lancement application 3D. Pilote NVIDIA version 2020 (4 ans). OS Windows 10 22H2 à jour.
    \textbf{2. Hypothèses :} H1 (Très probable) : Incompatibilité pilote graphique ancien / Windows 10 récent / DirectX 12 du jeu. H2 : Surchauffe GPU (pâte thermique sèche, ventilateur bloqué). H3 : Carte graphique défectueuse (VRAM). H4 : Alimentation insuffisante (pic conso jeu).
    \textbf{3. Tests :} T1: \texttt{dxdiag} $\to$ Onglet Affichage $\to$ Vérifier version pilote / Notes (problèmes). T2: Désinstaller pilote via DDU (Display Driver Uninstaller) en Mode Sans Échec $\to$ Redémarrer. T3: Installer pilote latest « Game Ready » depuis \texttt{nvidia.com/Download} (détection auto ou manuelle GPU). T4: Lancer jeu $\to$ Surveiller T° GPU (HWMonitor/MSI Afterburner) < 83°C. T5 (si T4 échoue) : Test carte sur autre PC / Test autre carte sur ce PC.
    \textbf{4. Solution :} Mise à jour pilote graphique via DDU (clean install) + vérification T°/poussière. Si H2 : Nettoyage + pâte thermique. Si H3/H4 : Remplacement matériel (hors premier niveau $\to$ escalade).
    \textbf{5. Validation :} Jeu lancé 30 min sans crash. T° stable 65°C. Bureau fluide. Fiche intervention : « BSOD VIDEO\_TDR\_FAILURE $\to$ Pilote legacy 2020. Clean install DDU + Driver 550.xx. OK. »
\end{enumerate}
\end{corrige}

\newpage

\coursec{Activité d'intégration}

\courssub{Objectifs de l'activité}
Cette activité d'intégration vise à mobiliser l'ensemble des compétences acquises lors de la leçon « Maintenance de premier niveau » à travers une mise en situation professionnelle réaliste. À l'issue de ce travail, tu seras capable de :
\begin{itemize}
    \item \textbf{Diagnostiquer} de manière structurée les dysfonctionnements courants d'un poste de travail (lenteur, saturation du disque, sécurité obsolète, saleté physique, connexions défaillantes).
    \item \textbf{Différencier} et \textbf{classer} les opérations de maintenance en \cle{préventive} (anticiper la panne) et \cle{corrective} (réparer la panne), en respectant scrupuleusement les \cle{règles de sécurité électrique et matérielle}.
    \item \textbf{Exécuter} les interventions logicielles de premier niveau autorisées en environnement scolaire : nettoyage de disque, mise à jour des signatures antivirus et du système, désinstallation propre de programmes indésirables, gestion des programmes au démarrage.
    \item \textbf{Rédiger} une \cle{fiche de maintenance} tracée, justifiant chaque action technique par sa cause et son effet attendu, et proposer un \cle{calendrier de maintenance préventive} pérenne pour un parc informatique scolaire.
    \item \textbf{Communiquer} oralement une démarche technique structurée devant un public (la classe), en utilisant le vocabulaire précis du métier (symptôme, cause probable, action corrective, validation).
\end{itemize}

\courssub{Situation}
\textbf{Contexte :} Tu es élève en classe de 3ème au \textbf{Collège Bilingue de Mvog-Ada} (Yaoundé). La salle informatique compte 20 postes de travail sous \texttt{Windows 10 Pro} (architecture 64 bits, 8 Go de RAM, SSD 256 Go pour le système + disque dur 1 To pour les données). Le professeur d'informatique, M. \textsc{Fouda}, a constaté une dégradation générale des performances depuis la rentrée de septembre. Il te désigne, avec ton binôme, comme « \emph{Techniciens de premier niveau} » pour la semaine. Votre mission : remettre en état le \textbf{Poste n°07}, le plus critique, et établir un protocole pour l'ensemble de la salle.

\textbf{État des lieux du Poste n°07 (constaté le 15 octobre 2024) :}
\begin{enumerate}
    \item \textbf{Démarrage anormalement long :} plus de \textbf{3 minutes 30 secondes} entre l'appui sur le bouton \texttt{Power} et l'affichage du Bureau utilisable. Le disque (voyant rouge clignotant) reste sollicité à 100\% pendant 2 minutes après l'ouverture de session.
    \item \textbf{Espace disque critique :} sur le disque système \texttt{C:} (SSD 256 Go), il reste \textbf{12,4 Go libres} (environ 5\%). L'outil \texttt{Stockage} signale : \texttt{Fichiers temporaires : 38 Go}, \texttt{Corbeille : 4,2 Go}, \texttt{Anciennes versions de Windows : 15 Go}.
    \item \textbf{Sécurité compromise :} \texttt{Windows Defender} affiche « \emph{Les définitions de virus ne sont pas à jour} ». Dernière mise à jour des signatures : \textbf{12 février 2024} (il y a 8 mois). Aucune analyse planifiée n'a été effectuée depuis mars. L'icône du Centre de sécurité est jaune (alerte).
    \item \textbf{Clavier et souris encrassés :} les touches \texttt{Espace}, \texttt{Entrée} et \texttt{Maj} gauche accrochent. Poussière visible entre les touches. Capteur optique de la souris obstrué par des fibres (déplacement saccadé du pointeur).
    \item \textbf{Logiciels parasites :} dans \texttt{Paramètres > Applications > Applications installées}, on trouve : \texttt{WebCompanion} (installé le 03/10/2024), \texttt{Candy Crush Saga} (Microsoft Store, 1,2 Go), \texttt{Driver Updater Pro} (version d'essai expirée, 08/10/2024). Ces programmes se lancent au démarrage (onglet \texttt{Démarrage} du \texttt{Gestionnaire des tâches} : impact « Élevé » pour \texttt{Driver Updater Pro}).
    \item \textbf{Affichage instable :} l'écran (Dell 24 pouces, entrée HDMI) scintille par intermittence (écran noir d'une seconde toutes les 30 à 40 secondes). Le câble HDMI, vérifié côté écran, est bien enfoncé ; côté unité centrale, il est simplement \emph{posé} dans le port, mal enfoncé, et bouge au moindre contact.
\end{enumerate}

\courssub{Consignes}
Travaillant par binôme, tu réaliseras les quatre étapes suivantes. Chaque étape fait l'objet d'une production écrite (sur papier ou au traitement de texte) qui sera évaluée.

\begin{enumerate}
    \item \textbf{Diagnostic structuré (15 min) :} remplis le tableau de diagnostic ci-dessous en identifiant pour chaque symptôme : la \cle{cause probable} (technique), le \cle{type de maintenance} requis (Préventive / Corrective / Les deux), et la \cle{priorité} (Urgent / Important / Secondaire). Justifie chaque choix en une phrase.
    \begin{center}
    \begin{tabularx}{\linewidth}{|X|l|l|l|X|}
    \hline
    \rowcolor{popblueL}
    \textbf{Symptôme observé} & \textbf{Cause probable} & \textbf{Type maintenance} & \textbf{Priorité} & \textbf{Justification technique} \\
    \hline
    Démarrage > 3 min 30 s & & & & \\
    \hline
    Disque C: plein à 95\% & & & & \\
    \hline
    Antivirus périmé (8 mois) & & & & \\
    \hline
    Clavier/Souris sales & & & & \\
    \hline
    Logiciels parasites au démarrage & & & & \\
    \hline
    Scintillement écran (câble HDMI) & & & & \\
    \hline
    \end{tabularx}
    \end{center}

    \item \textbf{Plan d'intervention ordonné et sécurisé (10 min) :} rédige la liste numérotée des opérations à effectuer, \textbf{dans l'ordre chronologique strict}, en indiquant pour chacune : (a) l'action précise, (b) la règle de sécurité associée (ex. : « Poste éteint et débranché », « Session administrateur requise »), (c) le type (Matériel / Logiciel). \textbf{Attention :} la sécurité physique précède toujours toute manipulation logicielle.

    \item \textbf{Exécution des opérations logicielles autorisées (30 min) :} sur le poste (ou en simulation guidée par le professeur), réalise \textbf{uniquement} les actions logicielles de premier niveau (nettoyage de disque, mise à jour de Defender, désinstallation des programmes parasites, désactivation des programmes au démarrage). Pour chaque action effectuée, note : l'outil Windows utilisé (chemin d'accès exact), la durée approximative, le résultat observé (ex. : « Espace libéré : 42 Go »).

    \item \textbf{Fiche de maintenance et calendrier préventif (15 min) :}
    \begin{itemize}
        \item Rédige la \textbf{Fiche de maintenance du Poste n°07} (en-tête : Date, Techniciens, Poste, Système d'exploitation, Durée totale). Corps : tableau récapitulatif \textit{Problème $\rightarrow$ Action $\rightarrow$ Résultat $\rightarrow$ Validation}.
        \item Propose un \textbf{Calendrier de maintenance préventive mensuel/trimestriel} pour les 20 postes de la salle, sous forme de tableau (Périodicité | Tâche | Responsable | Outils/Indicateurs). Pense à la poussière, aux mises à jour Windows/Defender, au nettoyage de disque, à la vérification des câbles, à la sauvegarde des documents des élèves.
    \end{itemize}
\end{enumerate}

\courssub{Ressources à mobiliser}
Pour réussir cette activité, tu t'appuieras sur les savoirs essentiels de la leçon :
\begin{itemize}
    \item \textbf{Distinction maintenance préventive / corrective :} la maintenance \cle{préventive} regroupe les actions planifiées pour éviter la panne (nettoyage physique, mises à jour, sauvegardes). La maintenance \cle{corrective} regroupe les actions réalisées après une panne avérée (nettoyage d'un disque saturé, remise en place d'un câble, désinfection d'un virus). \textit{Une même intervention peut avoir les deux aspects.}
    \item \textbf{Règles de sécurité impératives (ordre strict) :}
    \begin{enumerate}
        \item \textbf{Éteindre} le système proprement (Démarrer $\rightarrow$ Arrêter) ; si le poste est gelé, forcer l'arrêt par un appui long (environ 5 s) sur le bouton Power.
        \item \textbf{Débrancher} le cordon d'alimentation secteur (prise murale \textbf{puis} côté unité centrale).
        \item \textbf{Appuyer} 2 à 3 fois sur le bouton Power pour évacuer l'électricité résiduelle (condensateurs).
        \item \textbf{Se décharger de l'électricité statique} (bracelet antistatique, ou toucher le châssis métallique nu) avant de toucher un composant interne.
        \item \textbf{Ne jamais pulvériser} de produit directement sur le matériel : appliquer le produit sur un chiffon microfibre.
        \item \textbf{Remonter} dans l'ordre inverse : brancher $\rightarrow$ allumer $\rightarrow$ vérifier le démarrage (BIOS/POST) $\rightarrow$ ouvrir la session.
    \end{enumerate}
    \item \textbf{Outils Windows 10 de premier niveau :}
    \begin{itemize}
        \item \texttt{Nettoyage de disque} (\texttt{cleanmgr}, ou Paramètres $\rightarrow$ Système $\rightarrow$ Stockage $\rightarrow$ Fichiers temporaires) : supprime les fichiers temporaires, vide la corbeille, efface les fichiers de mise à jour et le dossier \texttt{Windows.old}.
        \item \texttt{Paramètres > Applications > Applications installées} : désinstallation propre (menu \texttt{...} $\rightarrow$ \texttt{Désinstaller}).
        \item \texttt{Gestionnaire des tâches} (\texttt{Ctrl+Maj+Échap}) $\rightarrow$ onglet \texttt{Démarrage} : \texttt{Désactiver} les programmes inutiles (impact sur le temps de démarrage).
        \item \texttt{Sécurité Windows} (Defender) $\rightarrow$ \texttt{Protection contre les virus et menaces} $\rightarrow$ \texttt{Rechercher des mises à jour} puis \texttt{Analyse rapide}.
        \item \texttt{Paramètres > Mise à jour et sécurité > Windows Update} : \texttt{Rechercher des mises à jour}.
    \end{itemize}
    \item \textbf{Traçabilité :} une intervention non documentée n'a pas existé. La fiche de maintenance est une preuve professionnelle et un historique pour le prochain technicien.
\end{itemize}

\courssub{Démarche et corrigé commenté}
\begin{exoresolu}[Corrigé type de l'activité « Technicien d'un jour »]
\textbf{ÉTAPE 1 : Diagnostic structuré (tableau complété)}

\begin{center}
\begin{tabularx}{\linewidth}{|X|X|l|l|X|}
\hline
\rowcolor{popblueL}
\textbf{Symptôme observé} & \textbf{Cause probable} & \textbf{Type maintenance} & \textbf{Priorité} & \textbf{Justification technique} \\
\hline
Démarrage > 3 min 30 s & Trop de programmes lancés au démarrage (Driver Updater, WebCompanion) + disque système presque saturé & Corrective + Préventive & \textbf{Urgent} & Temps de démarrage inacceptable pour un cours (perte de 10 min par séance). \\
\hline
Disque C: plein à 95\% (12,4 Go libres) & Accumulation de fichiers temporaires (38 Go), corbeille non vidée (4,2 Go), dossier \texttt{Windows.old} (15 Go) conservé après une mise à jour & Corrective (libérer l'espace vital) & \textbf{Urgent} & Sous le seuil critique de 10 à 15\% d'espace libre : risque de plantage du système et impossibilité d'installer les mises à jour Windows. \\
\hline
Antivirus périmé (8 mois) & Échec des mises à jour automatiques (problème réseau ou service arrêté) ; aucune analyse planifiée & Corrective (MAJ immédiate) + Préventive (réactiver la planification) & \textbf{Urgent} & 8 mois sans signatures = exposition aux virus récents ; risque de propagation à tout le réseau de la salle. \\
\hline
Clavier/Souris sales & Poussière, résidus alimentaires, cheveux entre les touches ; capteur optique de la souris obstrué & Préventive (hygiène) + Corrective (restaurer le fonctionnement) & \textbf{Important} & Inconfort de saisie, erreurs de frappe, hygiène collective (postes partagés par plusieurs classes). \\
\hline
Logiciels parasites au démarrage & Installations « rapides » (Suivant/Suivant/Terminer) de logiciels gratuits qui ajoutent des programmes indésirables au démarrage & Corrective (désinstallation) + Préventive (éducation des utilisateurs) & \textbf{Important} & Ils consomment mémoire, processeur et disque en arrière-plan. \texttt{Driver Updater Pro} est un scareware potentiel (fausses alertes pour forcer un achat). \\
\hline
Scintillement écran (câble HDMI) & Connecteur HDMI côté unité centrale mal enfoncé $\rightarrow$ micro-coupures du signal vidéo & Corrective (remise en place physique) & \textbf{Urgent} & Fatigue visuelle immédiate (maux de tête) ; risque de détérioration du port à terme. \\
\hline
\end{tabularx}
\end{center}

\textbf{ÉTAPE 2 : Plan d'intervention ordonné et sécurisé (chronologie)}

\begin{enumerate}
    \item \textbf{[SÉCURITÉ / MATÉRIEL] Éteindre le poste proprement} (Démarrer $\rightarrow$ Arrêter). \textit{Règle : arrêt logiciel avant toute coupure secteur, pour ne pas perdre de données.}
    \item \textbf{[SÉCURITÉ / MATÉRIEL] Débrancher le cordon d'alimentation} (prise murale puis côté unité centrale). Appuyer 2 fois sur le bouton Power (évacuation de l'électricité résiduelle). \textit{Règle : risque zéro d'électrocution et de court-circuit pendant le nettoyage et la manipulation des câbles.}
    \item \textbf{[MATÉRIEL - CORRECTIVE] Nettoyer clavier et souris :} bombe à air sec (poussière entre les touches) $\rightarrow$ chiffon microfibre légèrement imbibé d'alcool isopropylique à 70\% (touches, coque) $\rightarrow$ souffler sur le capteur optique de la souris. \textit{Règle : jamais de liquide pulvérisé directement ; attendre le séchage complet (2 min) avant rebranchement.}
    \item \textbf{[MATÉRIEL - CORRECTIVE] Rebrancher fermement le câble HDMI :} l'enfoncer à fond dans le port de la carte graphique côté unité centrale ; vérifier aussi le côté écran. \textit{Règle : manipulation douce, sans forcer latéralement sur le port.}
    \item \textbf{[SÉCURITÉ / MATÉRIEL] Rebrancher l'alimentation} (côté unité centrale puis prise murale). Allumer le poste. \textit{Validation : démarrage normal (bip unique du POST), pas de message « No Signal » sur l'écran, image stable.}
    \item \textbf{[LOGICIEL - CORRECTIVE] Ouvrir la session. Lancer le \texttt{Gestionnaire des tâches} (\texttt{Ctrl+Maj+Échap}) $\rightarrow$ onglet \texttt{Démarrage}. Sélectionner \texttt{Driver Updater Pro} et \texttt{WebCompanion} $\rightarrow$ clic droit $\rightarrow$ \texttt{Désactiver}.} \textit{Justification : action immédiate qui réduit la charge au prochain démarrage.}
    \item \textbf{[LOGICIEL - CORRECTIVE] Nettoyage de disque (Paramètres > Système > Stockage > Fichiers temporaires) :} cocher \texttt{Fichiers temporaires} (38 Go), \texttt{Corbeille} (4,2 Go), \texttt{Installations Windows précédentes} / \texttt{Windows.old} (15 Go), \texttt{Fichiers d'optimisation de livraison}. Cliquer \texttt{Supprimer les fichiers}. \textit{Durée : 5 à 10 min. Résultat attendu : environ 57 Go libérés.}
    \item \textbf{[LOGICIEL - CORRECTIVE] Désinstallation des programmes parasites (Paramètres > Applications > Applications installées) :} rechercher \texttt{WebCompanion} $\rightarrow$ \texttt{...} $\rightarrow$ \texttt{Désinstaller} (suivre l'assistant, refuser les offres additionnelles). Idem pour \texttt{Driver Updater Pro} et \texttt{Candy Crush Saga}. \textit{Note : la désinstallation peut demander une confirmation administrateur (UAC) $\rightarrow$ demander au professeur. En cas de refus, se limiter à la désactivation au démarrage (étape 6).}
    \item \textbf{[LOGICIEL - CORRECTIVE + PRÉVENTIVE] Mise à jour de Windows Defender :} \texttt{Sécurité Windows} > \texttt{Protection contre les virus et menaces} > \texttt{Rechercher des mises à jour}. Attendre la fin du téléchargement. Lancer ensuite une \texttt{Analyse rapide}. \textit{Durée : 3 à 5 min. Validation : « Aucune menace détectée », date des définitions = aujourd'hui.}
    \item \textbf{[LOGICIEL - PRÉVENTIVE] Lancer Windows Update :} \texttt{Paramètres > Mise à jour et sécurité > Windows Update} > \texttt{Rechercher des mises à jour}. Installer les mises à jour de qualité et de sécurité. Planifier le redémarrage hors des heures de cours. \textit{Justification : corriger les failles de sécurité exploitées par les virus.}
    \item \textbf{[VALIDATION FINALE] Redémarrer le poste.} Chronométrer le démarrage. Vérifier : espace libre sur C: supérieur à 20\%, Defender au vert, pas de scintillement, clavier fluide, aucune fenêtre parasite au démarrage.
\end{enumerate}

\textbf{ÉTAPE 3 : Exécution logicielle (journal de bord simulé)}

\begin{center}
\begin{tabularx}{\linewidth}{|l|X|l|X|}
\hline
\rowcolor{popgreenL}
\textbf{Heure} & \textbf{Action réalisée (outil Windows exact)} & \textbf{Durée} & \textbf{Résultat / Observation} \\
\hline
10:05 & \texttt{Paramètres > Système > Stockage > Fichiers temporaires} : suppression (temporaires, corbeille, Windows.old, optimisation de livraison) & 8 min & \textbf{+57,2 Go libérés.} Le disque C: passe de 12,4 Go libres (5\%) à 69,6 Go libres (27\%). \\
\hline
10:15 & \texttt{Paramètres > Applications} : désinstallation de \texttt{WebCompanion}, \texttt{Driver Updater Pro}, \texttt{Candy Crush Saga} (validation administrateur par le professeur) & 5 min & 3 programmes supprimés. Espace gagné : +1,5 Go. Leurs entrées disparaissent de l'onglet Démarrage. \\
\hline
10:22 & \texttt{Gestionnaire des tâches > Démarrage} : vérification. Seuls \texttt{OneDrive} et l'icône de \texttt{Sécurité Windows} restent activés. & 1 min & Impact au démarrage passé de « Élevé » à « Faible ». \\
\hline
10:25 & \texttt{Sécurité Windows > Protection contre les virus et menaces > Rechercher des mises à jour} & 2 min & \textbf{MAJ réussie.} Date des définitions = 15/10/2024. Icône verte. \\
\hline
10:28 & \texttt{Analyse rapide} (Defender) & 4 min & \textbf{0 menace détectée.} 12 450 fichiers analysés. \\
\hline
10:33 & \texttt{Paramètres > Mise à jour et sécurité > Windows Update} : recherche puis installation de 3 mises à jour cumulatives & 12 min & Redémarrage planifié à 12h00 (pause déjeuner). \\
\hline
12:00 & Redémarrage après mises à jour + chronométrage du démarrage & 1 min 15 s & \textbf{Démarrage : 1 min 15 s (contre 3 min 30 s avant).} Poste opérationnel. \\
\hline
\end{tabularx}
\end{center}

\textbf{ÉTAPE 4 : Fiche de maintenance et calendrier préventif}

\textbf{FICHE DE MAINTENANCE — POSTE N°07 — COLLÈGE BILINGUE DE MVOG-ADA}\par
\begin{center}
\begin{tabularx}{\linewidth}{|l|X|}
\hline
\rowcolor{poporangeL}
\textbf{Champ} & \textbf{Information} \\
\hline
Date d'intervention & 15 octobre 2024 \\
Techniciens & \textit{[Noms des deux élèves]} \\
Poste concerné & POSTE-07 (salle informatique, 2\up{e} rangée, près de la fenêtre) \\
Système d'exploitation & Windows 10 Pro 22H2 — 64 bits \\
Durée totale d'intervention & 55 minutes (matériel : 10 min / logiciel : 45 min) \\
\hline
\end{tabularx}
\end{center}

\begin{center}
\begin{tabularx}{\linewidth}{|X|X|X|l|}
\hline
\rowcolor{poppurpleL}
\textbf{Problème identifié} & \textbf{Action effectuée} & \textbf{Résultat mesuré} & \textbf{Validation} \\
\hline
Démarrage excessif (3 min 30 s) & Désactivation des programmes parasites au démarrage + nettoyage de disque + mises à jour Windows & Démarrage = 1 min 15 s & OK \\
\hline
Disque C: saturé (95\%) & Nettoyage des fichiers temporaires, corbeille, Windows.old & Espace libre = 69,6 Go (27\%) & OK \\
\hline
Antivirus périmé (8 mois) & Mise à jour des définitions Defender + analyse rapide & Définitions du 15/10/2024, 0 menace & OK \\
\hline
Clavier/Souris sales & Air comprimé + chiffon microfibre et alcool isopropylique (extérieur) & Touches fluides, capteur précis & OK \\
\hline
Logiciels parasites (3) & Désinstallation propre via Paramètres (validation administrateur) & Programmes absents de la liste & OK \\
\hline
Scintillement écran & Rebranchement ferme du câble HDMI côté unité centrale & Image stable pendant 30 min de test & OK \\
\hline
\end{tabularx}
\end{center}
\textbf{Observations finales :} le poste est de nouveau utilisable pour les cours. \textbf{Recommandation :} sensibiliser les élèves au mode d'installation « Personnalisée » afin de refuser les logiciels sponsorisés, et prévoir une sauvegarde régulière des documents sur le disque de données (D:).

\textbf{CALENDRIER DE MAINTENANCE PRÉVENTIVE — SALLE INFORMATIQUE (20 POSTES)}\par
\begin{center}
\begin{tabularx}{\linewidth}{|l|X|l|X|}
\hline
\rowcolor{popgoldL}
\textbf{Périodicité} & \textbf{Tâche de maintenance préventive} & \textbf{Responsable} & \textbf{Outils / Indicateurs de validation} \\
\hline
\textbf{Quotidien} (fin de journée) & Extinction propre des postes (pas de veille prolongée). Vider la corbeille de la session. & Élèves / Professeur & Postes éteints, corbeille vide. \\
\hline
\textbf{Hebdomadaire} (vendredi 15h) & \textbf{Nettoyage physique :} claviers (air sec), souris, écrans (chiffon sec), entrées d'air des unités centrales. Vérification visuelle des câbles (HDMI, alimentation, réseau). & Binôme « Techniciens de la semaine » (rotation des élèves) + professeur & Fiche « État de propreté / Câbles » signée. Aucune touche collante. \\
\hline
\textbf{Hebdomadaire} (lundi 7h30) & \textbf{Mise à jour de Windows Defender :} « Rechercher des mises à jour » + « Analyse rapide » sur 4 postes différents chaque semaine (cycle complet en 5 semaines). & Professeur & Centre de sécurité au vert, date de MAJ récente, rapport d'analyse : 0 menace. \\
\hline
\textbf{Mensuel} (1\up{er} mardi du mois) & \textbf{Nettoyage de disque} sur chaque poste : fichiers temporaires, corbeille, fichiers d'optimisation de livraison. & Professeur + binôme de la semaine & Espace libre sur C: supérieur à 20\% sur tous les postes. \\
\hline
\textbf{Mensuel} (2\up{e} mardi) & \textbf{Windows Update :} installation des mises à jour cumulatives + redémarrage planifié à la pause de midi. & Professeur & Historique Windows Update : mises à jour « réussies » sur les 20 postes. \\
\hline
\textbf{Trimestriel} (début de chaque trimestre) & \textbf{Audit complet :} vérification de l'état des disques, dépoussiérage interne des unités centrales (ventilateurs), sauvegarde des documents des élèves vers un support externe ou un serveur. & Professeur + technicien externe si besoin & Rapport signé : état des disques, sauvegarde datée, postes dépoussiérés. \\
\hline
\end{tabularx}
\end{center}
\end{exoresolu}

\courssub{Bilan de l'activité}
Cette activité d'intégration a permis de mettre en évidence plusieurs acquis majeurs du programme de 3ème :
\begin{itemize}
    \item \textbf{La rigueur de la démarche diagnostique :} un symptôme unique (la lenteur) peut avoir des causes multiples (logicielles : programmes au démarrage, disque plein ; matérielles : poussière, câble défectueux). Le tableau de diagnostic oblige à isoler chaque cause pour ne pas « soigner » le mauvais organe.
    \item \textbf{La primauté de la sécurité :} l'ordre strict (Éteindre $\rightarrow$ Débrancher $\rightarrow$ Évacuer l'électricité résiduelle $\rightarrow$ Se protéger de l'électricité statique) n'est pas une formalité : c'est une condition de survie du matériel et de l'opérateur. L'omettre expose à l'électrocution, aux courts-circuits et à la perte de données.
    \item \textbf{L'efficacité des outils natifs de Windows :} pas besoin de logiciels tiers payants (souvent porteurs de publicités ou de programmes indésirables). Le \texttt{Nettoyage de disque}, le \texttt{Gestionnaire des tâches}, les \texttt{Paramètres}, la \texttt{Sécurité Windows} et \texttt{Windows Update} couvrent l'essentiel des besoins de premier niveau. Leur maîtrise rend l'élève autonome et critique face aux scarewares comme \texttt{Driver Updater Pro}.
    \item \textbf{La distinction opératoire préventif / correctif :} nettoyer le clavier est préventif (hygiène, longévité) ; rebrancher le câble HDMI est correctif (panne avérée) ; vider la corbeille est correctif (saturation), mais le planifier chaque mois devient préventif. Cette double lecture structure la pensée « maintenance ».
    \item \textbf{La traçabilité professionnelle :} la fiche de maintenance n'est pas un simple exercice scolaire : c'est le contrat de confiance entre le technicien, l'utilisateur et l'établissement. Elle permet la continuité du service (le prochain technicien sait ce qui a été fait) et la capitalisation des connaissances (historique des pannes récurrentes).
    \item \textbf{La gestion d'un parc à l'échelle :} le calendrier préventif transforme l'approche « réactive » (je répare quand ça casse) en approche « proactive » (j'empêche la panne). Il introduit les notions de \cle{planification}, de rotation des responsabilités et d'\cle{indicateurs} de suivi (espace disque libre, date des mises à jour).
\end{itemize}
\textbf{Compétence validée :} « \emph{Appliquer les notions de maintenance de premier niveau à une situation concrète de la vie courante (la salle informatique du collège) et justifier sa démarche par écrit et à l'oral.} » \cle{Bravo, techniciens !} Le Poste n°07 est sauvé, la salle a son planning, et vous avez prouvé que l'informatique, c'est d'abord de la méthode et de la rigueur.

\newpage

\coursec{Méthodes et exercices résolus}

\courssub{Notions de base sur la maintenance informatique}

\begin{experience}[Observation : Le parc informatique du cybercafé scolaire]
    \textbf{Contexte :} L'établissement gère un cybercafé de 20 postes sous Windows 10. En début d'année, tout fonctionne. Au bout de trois mois, les plaintes s'accumulent : « Le poste 7 plante au démarrage », « La souris du poste 1 ne marche plus », « Le poste 10 rame affreusement », « Le ventilateur du poste 3 fait un bruit d'enfer ».
    \textbf{Problème :} Face à cette diversité de pannes, comment l'élève référent maintenance (niveau 3ème) doit-il réagir ? Quelle intervention est de son ressort, et quand doit-il alerter le technicien de l'établissement ?
    \textbf{Objectif :} Comprendre la classification des maintenances (type, nature, niveau) pour prendre la bonne décision immédiate.
\end{experience}

\begin{definition}[Maintenance informatique]
    Ensemble des actions techniques (matérielles et logicielles) visant à maintenir ou rétablir un bien (ordinateur, périphérique, logiciel, réseau) dans un état spécifié pour qu'il puisse assurer un service déterminé.
\end{definition}

\begin{definition}[Types de maintenance]
    \begin{itemize}
        \item \cle{Maintenance préventive :} Actions planifiées pour éviter l'apparition de pannes (nettoyage poussière, mises à jour, sauvegardes, vérification connectique).
        \item \cle{Maintenance curative :} Actions menées après une panne pour rétablir le service (remplacement pièce, réparation système, désinfection virus).
    \end{itemize}
\end{definition}

\begin{definition}[Niveaux de maintenance (Cadre 3ème / BEPC)]
    \begin{itemize}
        \item \cle{Niveau 1 (Premier niveau) :} Interventions sans ouverture du boîtier (sauf nettoyage air comprimé si autorisé) ou actions logicielles standards : nettoyage disque, défragmentation/TRIM, installation/mise à jour logiciels, scan antivirus, restauration système, vérification câbles/piles, sauvegarde utilisateur. \textbf{Autorisé à l'élève formé de 3ème.}
        \item \cle{Niveau 2 :} Remplacement de composants internes (RAM, disque, alimentation, ventilateur, carte graphique), réinstallation OS, configuration réseau avancée, récupération données sur disque défaillant. \textbf{Réservé au technicien qualifié de l'établissement.}
        \item \cle{Niveau 3 :} Réparation carte mère (microsoudure), expertise judiciaire (forensic), administration serveur, architecture réseau. \textbf{Spécialiste externe / Expert.}
    \end{itemize}
\end{definition}

\begin{aretenir}[Règle de décision en 3ème]
    \begin{enumerate}
        \item Analyser le symptôme (bloquant ? ralenti ? bruit ?).
        \item Déterminer la nature (matériel ? logiciel ?).
        \item Évaluer le niveau requis.
        \item \textbf{Si Niveau 1 :} Appliquer la procédure adaptée (nettoyage, MAJ, scan, câbles).
        \item \textbf{Si Niveau 2 ou 3 :} \textbf{STOP.} Éteindre/Débrancher $\rightarrow$ Étiquette « En panne » $\rightarrow$ Fiche d'incident $\rightarrow$ Transmettre au responsable maintenance. \textbf{Ne jamais forcer, ne jamais ouvrir sous tension, ne jamais improviser au-delà du Niv 1.}
    \end{enumerate}
\end{aretenir}

\begin{methode}[Identifier le type et le niveau de maintenance adaptés à une situation]
    \textbf{Objectif :} Classer une intervention (préventive/curative, matériel/logiciel, premier/niveau supérieur) pour choisir la bonne procédure.
    \begin{enumerate}
        \item \textbf{Analyser le symptôme :} Panne bloquante (écran noir, ne démarre plus) $\rightarrow$ \emph{curative}. Ralentissement, bruit anormal, mise à jour due $\rightarrow$ \emph{préventive}.
        \item \textbf{Déterminer la nature :} Composant physique (ventilateur, RAM, disque, câble) $\rightarrow$ \emph{matérielle}. Système, application, virus, pilote $\rightarrow$ \emph{logicielle}.
        \item \textbf{Évaluer le niveau (compétence 3ème) :}
              \begin{itemize}
                  \item \cle{Niveau 1 :} Nettoyage poussière, vérification câbles/connectiques, nettoyage disque, défragmentation (HDD) / Optimisation (SSD), installation/mise à jour logiciels standards, scan antivirus, restauration système, sauvegarde utilisateur. \textbf{Action autorisée en 3ème.}
                  \item \cle{Niveau 2 :} Remplacement composants (RAM, disque, alimentation, écran), réinstallation OS, configuration réseau avancée, récupération données sur disque défaillant. \textbf{Nécessite technicien qualifié.}
                  \item \cle{Niveau 3 :} Réparation carte mère, microsoudure, forensic, administration serveur. \textbf{Spécialiste/expert.}
              \end{itemize}
        \item \textbf{Décider :} Si Niveau 1 $\rightarrow$ Appliquer la méthode correspondante. Si Niveau 2 ou 3 $\rightarrow$ \textbf{Arrêter}, protéger le poste (éteindre, débrancher), étiqueter « En panne », remplir fiche d'incident, transmettre au responsable maintenance.
    \end{enumerate}
\end{methode}

\begin{exoresolu}[Classification d'interventions en cybercafé géré par l'établissement]
Dans le cybercafé de l'établissement, vous observez les situations suivantes. Pour chacune, précisez : (a) Type (préventive/curative), (b) Nature (matérielle/logicielle), (c) Niveau (1, 2, 3), (d) Action immédiate à mener par un élève de 3ème formé.
\begin{enumerate}
    \item Le poste n°7 affiche un écran bleu (BSOD) au démarrage avec l'erreur \texttt{CRITICAL\_PROCESS\_DIED}.
    \item Le poste n°3 fonctionne mais le ventilateur de l'unité centrale fait un bruit de grattage violent.
    \item Le poste n°10 est très lent à l'ouverture de LibreOffice ; l'analyseur de stockage signale 92\% d'occupation sur le disque C: (SSD 256 Go).
    \item La souris sans fil du poste n°1 ne répond plus ; le témoin lumineux du récepteur USB ne clignote pas.
    \item Une alerte Windows Defender signale une menace \texttt{Trojan:Win32/Wacatac} sur le poste n°5 mise en quarantaine.
\end{enumerate}
\tcblower
\textbf{Solution détaillée :}

\begin{center}
\begin{tabularx}{\linewidth}{|c|X|X|X|X|}
\hline
\rowcolor{popblueL}
\textbf{Poste} & \textbf{(a) Type} & \textbf{(b) Nature} & \textbf{(c) Niveau} & \textbf{(d) Action élève 3ème} \\
\hline
n°7 (BSOD) & Curative (panne bloquante) & Logicielle (fichier système/pilote corrompu) & 1 (tentative) / 2 (si persiste) & 1. Redémarrer en mode sans échec (F8 ou Maj+Redémarrer). 2. Lancer \texttt{sfc /scannow} en invite admin. 3. Si échec $\rightarrow$ Restauration système à point antérieur. 4. Si échec $\rightarrow$ \textbf{Arrêter, étiqueter, signaler (Niv 2 : réinstall OS)}. \\
\hline
n°3 (Ventilateur bruit) & Préventive (risque surchauffe imminent) & Matérielle (roulement HS / poussière) & 2 (remplacement) / 1 (nettoyage si juste poussière) & 1. \textbf{Éteindre et débrancher immédiatement} (risque CPU grillé). 2. Ouvrir boîtier (si autorisé localement) $\rightarrow$ souffler poussière à l'air comprimé (bombe). 3. Si bruit persiste $\rightarrow$ \textbf{Ne pas remplacer le ventilo soi-même (Niv 2)}. Étiqueter « Ventilateur HS », signaler. \\
\hline
n°10 (Lent, disque 92\%) & Préventive (dégradation performance) & Logicielle (espace disque, fragmentation si HDD, fichiers temporaires) & 1 & 1. Lancer \texttt{cleanmgr} (Nettoyage de disque) : cocher fichiers temporaires, corbeille, logs, miniatures. 2. Vider corbeille. 3. Désinstaller logiciels inutiles (Panneau config). 4. \textbf{SSD : pas de défragmentation} (Optimiser lecteurs = TRIM). 5. Vérifier espace $> 20\%$ libre. \\
\hline
n°1 (Souris sans fil) & Curative (panne périphérique) & Matérielle (piles / récepteur / driver) & 1 & 1. Vérifier piles (remplacer par neuves). 2. Vérifier récepteur USB bien enfoncé (changer port USB). 3. Test sur autre poste $\rightarrow$ si OK, réinstaller driver (Gestionnaire périphériques). 4. Si KO sur autre poste $\rightarrow$ Matériel HS $\rightarrow$ Remplacer (stock établissement) ou signaler. \\
\hline
n°5 (Trojan quarantaine) & Curative (incident sécurité) & Logicielle (malware) & 1 (traitement standard) & 1. Ne pas restaurer le fichier. 2. Lancer \textbf{analyse complète} Windows Defender (ou Malwarebytes en complément). 3. Vérifier mises à jour Windows (Definitions virus). 4. Sensibiliser l'utilisateur (pièce jointe ? clé USB ?). 5. Si récidive $\rightarrow$ Signalement admin réseau (Niv 2). \\
\hline
\end{tabularx}
\end{center}

\textbf{Conclusion :} L'élève de 3ème gère les niveaux 1 (nettoyage, mises à jour, scans, vérifications câbles/piles, nettoyage disque). Tout démontage composant interne (ventilateur, RAM, disque) ou réinstallation OS relève du Niveau 2 : il sécurise le poste et transmet une fiche d'incident précise.
\end{exoresolu}

\courssub{Règles de sécurité avant toute intervention matériel}

\begin{experience}[Observation : Pourquoi tant de précautions ?]
    \textbf{Constat :} Un élève veut « juste » ajouter une barrette RAM. Il ouvre le boîtier PC allumé, tire la barrette à mains nues sur la moquette, la remet, referme. Résultat : PC ne démarre plus (bip mémoire) ou instabilité aléatoire 3 semaines plus tard.
    \textbf{Causes invisibles :} 1. Tension résiduelle dans les condensateurs (risque court-circuit / choc). 2. Décharge électrostatique (ESD) de 3000 V (imperceptible pour l'homme, fatale pour les puces CMOS). 3. Poussière/poils sur connecteurs dorés.
    \textbf{Leçon :} La sécurité n'est pas une option, c'est la condition \emph{sine qua non} de toute intervention matérielle.
\end{experience}

\begin{methode}[Appliquer les règles de sécurité électrique et ESD avant toute intervention matériel]
    \textbf{Objectif :} Protéger l'opérateur (risque électrique) et les composants (risque décharge électrostatique - ESD) avant d'ouvrir un boîtier ou manipuler des composants.
    \begin{enumerate}
        \item \textbf{Sécurité électrique (Priorité absolue) :}
              \begin{enumerate}
                  \item Éteindre l'ordinateur via l'OS (Arrêter), pas le bouton physique.
                  \item \textbf{Débrancher le cordon d'alimentation secteur (prise murale ET arrière bloc alim).}
                  \item Appuyer 5 secondes sur le bouton \textbf{Power} de la tour (vidange condensateurs / \emph{flea power}).
                  \item Vérifier l'absence de tension sur l'alimentation (pas de LED allumée sur carte mère).
              \end{enumerate}
        \item \textbf{Protection ESD (Décharge Électrostatique) :}
              \begin{enumerate}
                  \item Porter un \textbf{bracelet antistatique} relié à la terre (prise de terre murale via cordon 1 M$\Omega$ ou partie métallique non peinte du boîtier \textbf{après} mise à la terre du boîtier).
                  \item \textbf{Alternative si pas de bracelet :} Toucher systématiquement une partie métallique non peinte du boîtier (alimentation) \textbf{avant} de toucher tout composant, et répéter l'opération régulièrement. Travailler sur un \textbf{tapis antistatique} si possible.
                  \item Manipuler les composants (RAM, CPU, carte mère, disque) \textbf{uniquement par les bords} (carte PCB), jamais par les connecteurs dorés ni les puces.
                  \item Stocker les composants retirés dans des \textbf{sachets antistatiques} (rose/noir), jamais sur plastique, polystyrène, carton ou tissu.
              \end{enumerate}
        \item \textbf{Outillage :} Tournevis cruciforme (Phillips \#2) aimanté (pratique pour vis), pince à bec long (sauteurs), bombe air comprimé (tenue verticale), pinceau antistatique. \textbf{Interdit :} Aspirateur classique (génère ESD fort), chiffon synthétique, force excessive.
        \item \textbf{Environnement :} Sol non moquetté, taux d'humidité $> 40\%$ si possible, pas de boisson à proximité.
    \end{enumerate}
\end{methode}

\begin{attention}[Piège classique : La coupure secteur incomplète]
    « Je débranche juste la tour. » $\rightarrow$ \textbf{FAUX}. Il faut débrancher \textbf{la prise murale} (ou couper l'onduleur). Le bloc d'alimentation reste sous tension secteur (condensateurs chargés) tant qu'il est branché, même PC éteint. Le bouton Power 5s \textbf{après} débranchement est obligatoire pour vider l'énergie résiduelle (\emph{flea power}).
\end{attention}

\begin{exoresolu}[Sécurisation d'un poste avant remplacement d'une barrette RAM (Niveau 2 simulé)]
Un élève de 3ème doit aider le technicien à préparer le poste n°12 pour un ajout de RAM (DDR4). Le technicien demande à l'élève de « sécuriser le poste complètement » avant son arrivée. L'élève propose la séquence suivante :
\begin{enumerate}
    \item Éteindre Windows via le menu Démarrer.
    \item Débrancher le câble réseau (RJ45) et la souris.
    \item Ouvrir le boîtier en enlevant les deux vis pouce à l'arrière.
    \item Toucher le radiateur du processeur pour se décharger.
    \item Sortir la nouvelle RAM de sa boîte en carton et la poser sur la moquette.
    \item Attendre le technicien.
\end{enumerate}
\textbf{Corrigez la procédure en listant les erreurs et en donnant la séquence correcte, justifiée.}
\tcblower
\textbf{Analyse des erreurs de l'élève :}
\begin{enumerate}
    \item \textbf{Étape 2 (Incomplète) :} Débrancher RJ45/souris ne coupe pas l'alimentation. \textbf{Manque : Débrancher le cordon secteur (prise murale + arrière tour).}
    \item \textbf{Étape 3 (Prématurée/Dangereuse) :} Ouvrir le boîtier \textbf{avant} d'avoir vidé les condensateurs (bouton Power 5s) et sans protection ESD. Risque : choc résiduel ou ESD sur carte mère.
    \item \textbf{Étape 4 (Insuffisante) :} Toucher le radiateur CPU \textbf{après} ouverture ne protège pas l'ouverture elle-même. Il faut se décharger \textbf{avant} de toucher le boîtier/composants. De plus, le radiateur n'est pas nécessairement à la terre si le PC n'est pas branché secteur (prise de terre coupée).
    \item \textbf{Étape 5 (Critique ESD) :} Sortir la RAM de son sachet antistatique pour la poser sur du \textbf{carton} (isolant, charge statique) puis sur de la \textbf{moquette} (générateur ESD majeur). \textbf{Destruction quasi certaine de la barrette.}
    \item \textbf{Manque :} Pas de bracelet antistatique, pas de tapis, pas de bouton Power 5s.
\end{enumerate}

\textbf{Procédure correcte (Séquence rigoureuse) :}
\begin{enumerate}
    \item \textbf{Arrêt logiciel :} Menu Démarrer $\rightarrow$ Arrêter (attendre extinction complète ventilateurs/LED).
    \item \textbf{Coupure secteur :} Débrancher le cordon d'alimentation \textbf{à la prise murale} (ou éteindre l'onduleur), puis débrancher côté bloc alim.
    \item \textbf{Vidange condensateurs :} Appuyer \textbf{5 secondes} sur le bouton \texttt{Power} de la tour (LED clignote souvent une fois).
    \item \textbf{Mise à la terre opérateur :} Enfiler bracelet antistatique $\rightarrow$ connecter pince crocodile sur \textbf{prise de terre murale} (ou partie métallique non peinte du boîtier \textbf{si} le cordon secteur est branché \textbf{uniquement côté prise murale avec terre} — \emph{préférence : prise murale dédiée}). \textbf{Alternative :} Toucher prise de terre murale (vis plaque prise) ou radiateur chauffage central.
    \item \textbf{Préparation composant :} Laisser la nouvelle RAM \textbf{dans son sachet antistatique scellé} posé sur le tapis antistatique (ou sur le sachet lui-même).
    \item \textbf{Ouverture boîtier :} Retirer vis (tournevis), glisser panneau latéral. Poser panneau sur surface plane.
    \item \textbf{Attente technicien :} Ne rien toucher d'autre. Informer le technicien : « Poste sécurisé : secteur coupé, vidangé, bracelet porté, RAM en sachet sur tapis. »
\end{enumerate}
\textbf{Conclusion :} La sécurité électrique (coupure + vidange) précède la sécurité ESD (bracelet/tapis/manipulation bords). Un composant ne quitte jamais son emballage antistatique avant l'insertion immédiate.
\end{exoresolu}

\courssub{Maintenance matérielle de premier niveau}

\begin{experience}[Observation : La poussière, ennemie silencieuse]
    \textbf{Constat :} Un poste en salle de TP (environnement poussiéreux) s'éteint brutalement au bout de 20 min de TP Python. Température CPU : 98°C. Après ouverture : radiateur CPU bouché par un « feutre » de poussière, ventilateur bloqué.
    \textbf{Action Niv 1 :} Nettoyage à l'air comprimé (pales bloquées), vérification flux d'air, rangement câbles. Résultat : 45°C en charge stable.
    \textbf{Enseignement :} La maintenance matérielle de premier niveau, c'est 80\% de nettoyage préventif et 20\% de vérification connectique. C'est l'action à plus fort retour sur investissement (coût quasi nul, gain durée de vie énorme).
\end{experience}

\begin{methode}[Réaliser la maintenance matérielle de premier niveau (Nettoyage \& Vérification connectique)]
    \textbf{Objectif :} Assurer la longévité du matériel par un nettoyage physique efficace et une vérification des connexions sans risque.
    \begin{enumerate}
        \item \textbf{Prérequis :} Appliquer la \textbf{Méthode Sécurité} (Secteur coupé, vidangé, ESD géré).
        \item \textbf{Nettoyage extérieur (Boîtier, Écran, Périphériques) :}
              \begin{itemize}
                  \item Éteindre/Débrancher écran. Nettoyer dalle avec \textbf{chiffon microfibre sec} (ou légèrement humidifié eau/dédié écran \textbf{sur le chiffon}, jamais direct). Pas de produit vitre (alcool/ammoniaque abîme traitements antireflet).
                  \item Clavier : Retourner, taper doucement $\rightarrow$ souffler bombe air comprimé (buses courtes, bombe \textbf{verticale} pour éviter projection propulseur liquide). Passer chiffon microfibre légèrement humide (eau + savon doux) sur touches.
                  \item Souris (optique) : Nettoyer capteur (soufflette), patins glisse (chiffon sec).
                  \item Boîtier : Aspirer/souffler grilles d'aération (avant, arrière, haut, fond). Chiffon humide sur faces externes.
              \end{itemize}
        \item \textbf{Nettoyage intérieur (Unité centrale) — \emph{Si autorisé par l'établissement} :}
              \begin{enumerate}
                  \item Ouvrir boîtier (panneau latéral verre/acier).
                  \item \textbf{Souffler la poussière} (bombe air comprimé verticale, bursts courts) : Ventilateurs (CPU, GPU, Boîtier, Alim), Radiateurs, Carte mère, Slots RAM/PCIe. \textbf{Bloquer les ventilateurs avec un doigt (ou pince en plastique)} pour ne pas les faire tourner trop vite (génère courant inverse / casse roulement).
                  \item Pinceau antistatique doux pour résidus tenaces sur carte mère.
                  \item \textbf{Ne pas aspirer} à l'intérieur (ESD, statique, risque arrachement composants CMS).
              \end{enumerate}
        \item \textbf{Vérification connectique (Visuel + Tactile) :}
              \begin{itemize}
                  \item Câbles alimentation (24 broches ATX, 4/8 broches CPU, PCIe GPU, SATA) : Bien enfoncés, clips verrouillés.
                  \item Câbles données (SATA SSD/HDD, M.2 vissé) : Bien insérés.
                  \item Barrettes RAM : Clips latéraux fermés (clic audible).
                  \item Carte graphique : Vissée au châssis, connecteur PCIe verrouillé.
                  \item Câbles façade (Power SW, Reset SW, LED, USB, Audio) : Repères \texttt{+}/\texttt{-} respectés, bien sur broches (manuel carte mère).
              \end{itemize}
        \item \textbf{Refermeture :} Panneau latéral, vis. Rebrancher secteur $\rightarrow$ Allumage $\rightarrow$ Vérifier BIOS (températures, détection RAM/Disques) $\rightarrow$ Démarrage OS.
    \end{enumerate}
\end{methode}

\begin{aretenir}[Fréquence recommandée]
    Environnement bureau/école propre : \textbf{Tous les 6 mois}. Environnement poussiéreux (atelier, cybercafé ouvert rue) : \textbf{Tous les 3 mois}. Écran/Clavier/Souris : \textbf{Hebdomadaire} (hygiène).
\end{aretenir}

\begin{exoresolu}[Diagnostic d'une surchauffe après nettoyage « maison »]
Après un nettoyage intérieur effectué par un élève non formé (aspirateur utilisé, ventilateur CPU fait tourner à l'air comprimé sans bloquer, pâte thermique non refaite), un poste présente : Température CPU 95°C au repos, ventilateur CPU à 3000 tr/min (bruyant), extinction aléatoire au bout de 10 min.
\textbf{Questions :}
\begin{enumerate}
    \item Identifiez 3 erreurs de manipulation ayant causé cet état.
    \item Proposez la procédure corrective de niveau 1 (sans changer la pâte thermique) pour tenter de stabiliser la température.
    \item Si la température reste $> 85^\circ$C au repos après correction, quelle est la décision (niveau) ?
\end{enumerate}
\tcblower
\textbf{Solution :}
\begin{enumerate}
    \item \textbf{Erreurs identifiées :}
          \begin{enumerate}
              \item \textbf{Usage aspirateur :} Génère forte électricité statique (ESD) $\rightarrow$ risque latence de mort composants (contrôleur ventilateur, capteur température, CPU) + aspiration pâte thermique possible.
              \item \textbf{Ventilateur CPU tourné à l'air comprimé sans blocage :} Rotation à haute vitesse $\rightarrow$ effet générateur (courant inverse injecté sur carte mère $\rightarrow$ grille contrôleur PWM / brûle piste) + usure prématurée roulement à billes. \textbf{Toujours bloquer les pales (doigt ou pince plastique).}
              \item \textbf{Pâte thermique non vérifiée/remise :} Si le radiateur a été décroché (même involontairement par aspiration), la pâte est cassée (bulles d'air) $\rightarrow$ contact CPU/Radiateur médiocre $\rightarrow$ surchauffe immédiate. (Note : l'élève n'a pas dit l'avoir décroché, mais l'aspirateur peut tirer le radiateur si vis fragiles).
          \end{enumerate}
    \item \textbf{Procédure corrective Niveau 1 (Sans changer pâte) :}
          \begin{enumerate}
              \item Sécuriser (Méthode 2 : Secteur coupé, vidangé, ESD).
              \item Ouvrir boîtier. \textbf{Inspecter visuellement} : Radiateur CPU bien plaqué ? Vis serrées en croix (diagonales) ? Pâte visible qui déborde légèrement (bon signe) ou sec/craquelé (mauvais) ?
              \item \textbf{Nettoyer correctement ventilateurs/radiateurs :} Bombe air comprimé verticale, \textbf{pales bloquées}. Souffler dans le sens inverse du flux d'air habituel (arrière $\rightarrow$ avant) pour sortir poussière. Nettoyer grilles boîtier.
              \item \textbf{Vérifier flux d'air :} Câbles rangés (colliers) $\rightarrow$ pas d'obstruction entrée (avant) / sortie (arrière/haut). Ventilateurs boîtier : Entrée avant/bas, Sortie arrière/haut.
              \item \textbf{BIOS/UEFI :} Vérifier courbe ventilateur CPU (Mode « Performance » ou « Turbo » / Courbe personnalisée : 50\% $\rightarrow$ 50°C, 100\% $\rightarrow$ 70°C). Désactiver mode « Silencieux ».
              \item \textbf{Test :} Refermer, démarrer. Surveiller température (HWMonitor / Gestionnaire des tâches Onglet Performance) 15 min charge (bench CPU simple).
          \end{enumerate}
    \item \textbf{Décision si $T > 85^\circ$C au repos :} \textbf{Niveau 2 obligatoire.} La pâte thermique est certainement sèche/mal étalée ou le radiateur mal monté (vis foirées, push-pins cassés Intel). Le remplacement de pâte thermique (nettoyage isopropanol 99\%, application grain de riz/pois, remontage pression croisée) ou le remplacement du ventirad (si ventilateur HS malgré nettoyage) dépasse le niveau 1. \textbf{Action :} Éteindre, étiqueter « Surchauffe CPU - Pâte thermique / Ventirad à vérifier », transmettre au technicien.
\end{enumerate}
\textbf{Conclusion :} Un nettoyage mal fait (aspirateur, ventilateur libre) peut créer la panne. Le niveau 1 corrige la poussière et le flux d'air ; le niveau 2 gère l'interface thermique CPU/Radiateur.
\end{exoresolu}

\courssub{Maintenance logicielle de premier niveau}

\begin{experience}[Observation : Le syndrome du « PC qui rame »]
    \textbf{Contexte :} Poste professeur, Windows 10, SSD 256 Go, 8 Go RAM. Boot : 4 min. Ouverture Outlook : 30 s. Gèle sur gros PDF.
    \textbf{Diagnostic rapide :} Disque C: à 95\% plein (maj Windows accumulées), 42 programmes au démarrage (dont 35 inutiles), Defender pas à jour, \texttt{sfc} signale erreurs.
    \textbf{Approche :} Une méthode structurée, outillée (outils natifs), tracée, sans réinstaller. C'est le cœur du niveau 1 logiciel.
\end{experience}

\begin{methode}[Réaliser la maintenance logicielle de premier niveau (Nettoyage, Optimisation, Sécurité)]
    \textbf{Objectif :} Restaurer les performances et la sécurité d'un poste Windows par des outils natifs, sans réinstallation.
    \begin{enumerate}
        \item \textbf{Sauvegarde préalable (Impératif) :} Copier documents utilisateur (Bureau, Documents, Images) sur clé USB/Disque externe/OneDrive. \textbf{Ne jamais nettoyer sans sauvegarde.}
        \item \textbf{Nettoyage disque (Espace) :}
              \begin{itemize}
                  \item \texttt{cleanmgr} (Win+S $\rightarrow$ Nettoyage de disque) $\rightarrow$ Sélectionner lecteur C: $\rightarrow$ Cocher : Fichiers temporaires, Fichiers temporaires Internet, Corbeille, Miniatures, Rapports d'erreurs, \textbf{Nettoyage mise à jour Windows} (libère plusieurs Go). $\rightarrow$ OK $\rightarrow$ Supprimer.
                  \item \textbf{Storage Sense (Stockage intelligent) :} Paramètres $\rightarrow$ Système $\rightarrow$ Stockage $\rightarrow$ Activer. Configurer : Supprimer fichiers temporaires (quotidien), Corbeille (30 j), Dossier Téléchargements (60 j si non ouverts).
              \end{itemize}
        \item \textbf{Optimisation lecteurs (Défragmentation / TRIM) :}
              \begin{itemize}
                  \item \texttt{dfrgui} (Défragmenter et optimiser les lecteurs).
                  \item \textbf{HDD (Disque dur mécanique) :} \textbf{Défragmenter} (réorganise blocs contigus $\rightarrow$ gain accès séquentiel). Planifier : Hebdomadaire.
                  \item \textbf{SSD (Disque état solide) :} \textbf{NE JAMAIS DÉFRAGMENTER} (usure cellules inutile). L'outil Windows détecte SSD $\rightarrow$ lance \textbf{TRIM} (commande ATA informe contrôleur blocs libérés $\rightarrow$ perf écriture). Planifier : Hebdomadaire (TRIM auto).
              \end{itemize}
        \item \textbf{Gestion démarrage (Rapidité boot) :}
              \begin{itemize}
                  \item Gestionnaire des tâches (Ctrl+Maj+Echap) $\rightarrow$ Onglet \textbf{Démarrage}.
                  \item Désactiver (Clic droit $\rightarrow$ Désactiver) : Logiciels non essentiels au boot (Spotify, Steam, Chrome Update, OneDrive si non utilisé, constructeurs bloatware). \textbf{Laisser :} Antivirus, Pilotes audio/GPU, Sécurité Windows.
              \end{itemize}
        \item \textbf{Sécurité / Antimalware :}
              \begin{itemize}
                  \item Mettre à jour \textbf{Sécurité Windows} (Protection virus \& menaces $\rightarrow$ Mises à jour $\rightarrow$ Rechercher mises à jour).
                  \item Lancer \textbf{Analyse complète} (peut durer 1h+). Quarantainer/Supprimer menaces.
                  \item \textbf{Complément (optionnel) :} \texttt{Malwarebytes Free} (scan unique, pas résident) pour second avis.
              \end{itemize}
        \item \textbf{Mises à jour système :} Paramètres $\rightarrow$ Mise à jour et sécurité $\rightarrow$ Rechercher mises à jour. Installer tout (Qualité, Pilotes, Fonctionnalités). Redémarrer.
        \item \textbf{Vérification intégrité fichiers système (Si instable) :}
              \begin{itemize}
                  \item Invite de commandes \textbf{Admin} (Win+X $\rightarrow$ Terminal Admin).
                  \item \texttt{sfc /scannow} (Répare fichiers système corrompus via cache WinSxS).
                  \item \texttt{DISM /Online /Cleanup-Image /RestoreHealth} (Répare l'image Windows elle-même, source Windows Update). Exécuter \textbf{avant} \texttt{sfc} si gros doute.
              \end{itemize}
    \end{enumerate}
\end{methode}

\begin{attention}[Erreur classique SSD : La défragmentation destructrice]
    « Je défragmente mon SSD pour qu'il soit plus vite. » $\rightarrow$ \textbf{DESTRUCTEUR}. La défragmentation écrit massivement (cycles P/E) sans gain d'accès (temps d'accès constant). Windows le sait : \texttt{dfrgui} lance TRIM sur SSD. Ne forcez jamais \texttt{defrag /c} sur SSD.
\end{attention}

\begin{exoresolu}[Restauration performance poste professeur - Scénario complet]
Le poste fixe du professeur principal (Windows 10 Pro, SSD 256 Go, 8 Go RAM, HDD 1 To données) met 4 min à démarrer, ouvre Outlook en 30 s, gèle sur fichiers lourds. L'élève de maintenance doit intervenir en Niveau 1.
\textbf{Travail demandé :}
\begin{enumerate}
    \item Établir la check-list ordonnée des 6 actions principales à mener (outils natifs uniquement).
    \item Pour l'action « Optimisation lecteurs », justifier techniquement la différence de traitement entre le SSD (C:) et le HDD (D:).
    \item Le \texttt{sfc /scannow} signale « Protection des ressources Windows a trouvé des fichiers corrompus et les a réparés ». Que faire ensuite pour valider la réparation ?
    \item Après intervention, le boot passe à 55 s. L'élève conclut « C'est bon, je rends le poste ». Quelle étape de validation manque-t-il ?
\end{enumerate}
\tcblower
\textbf{Solution :}
\begin{enumerate}
    \item \textbf{Check-list ordonnée (Méthode 4 appliquée) :}
          \begin{enumerate}
              \item \textbf{Sauvegarde} : Copier \texttt{C:\textbackslash Users\textbackslash Prof\textbackslash Desktop}, \texttt{Documents} sur disque D: ou clé USB.
              \item \textbf{Nettoyage disque (C:) :} \texttt{cleanmgr} $\rightarrow$ Fichiers système $\rightarrow$ Cocher « Nettoyage mise à jour Windows », « Fichiers temporaires », « Corbeille ». Exécuter.
              \item \textbf{Gestion démarrage :} \texttt{Ctrl+Maj+Echap} $\rightarrow$ Onglet Démarrage $\rightarrow$ Désactiver tout sauf : Sécurité Windows, Pilote GPU (NVIDIA/AMD/Intel), Realtek Audio, Synaptics Touchpad (si portable).
              \item \textbf{Optimisation lecteurs :} \texttt{dfrgui} $\rightarrow$ Analyser C: (SSD) et D: (HDD) $\rightarrow$ Optimiser (Lance TRIM sur C:, Défrag sur D:).
              \item \textbf{Sécurité :} Mise à jour définitions $\rightarrow$ Analyse complète Windows Defender.
              \item \textbf{Mises à jour Windows + Pilotes :} Paramètres $\rightarrow$ Mise à jour $\rightarrow$ Rechercher/Installer tout $\rightarrow$ Redémarrer. Puis site constructeur (Dell/HP/Lenovo) $\rightarrow$ Assistant maj pilotes (optionnel Niv 1+).
          \end{enumerate}
    \item \textbf{Justification SSD vs HDD :}
          \begin{itemize}
              \item \textbf{HDD (D:) :} Têtes de lecture mécaniques. Fichiers fragmentés $\rightarrow$ sauts de piste multiples $\rightarrow$ latence mécanique forte ($\sim$10 ms). \textbf{Défragmentation} regroupe blocs $\rightarrow$ lecture séquentielle rapide.
              \item \textbf{SSD (C:) :} Pas de pièce mobile. Accès aléatoire instantané ($\sim$0,1 ms). Fragmentation \textbf{sans impact perf}. Défragmentation = écritures inutiles $\rightarrow$ usure mémoire Flash (cycles P/E limités). \textbf{TRIM} (via Optimiser) informe contrôleur des blocs logiquement effacés (fichiers supprimés) pour qu'il les prépare (effacement bloc physique) en arrière-plan $\rightarrow$ maintient vitesse écriture.
          \end{itemize}
    \item \textbf{Post \texttt{sfc /scannow} :} Redémarrer l'ordinateur. Relancer \texttt{sfc /scannow} une seconde fois pour confirmer « Aucune violation d'intégrité trouvée ». Si persistant $\rightarrow$ Lancer \texttt{DISM /Online /Cleanup-Image /RestoreHealth} $\rightarrow$ Redémarrer $\rightarrow$ \texttt{sfc /scannow}.
    \item \textbf{Étape manquante :} \textbf{Validation fonctionnelle utilisateur (Recette).} Ne pas rendre sans : 1. Ouvrir session prof $\rightarrow$ Lancer Outlook, Word, Navigateur, Fichier lourd sur D: $\rightarrow$ Vérifier fluidité. 2. Vérifier imprimante réseau connectée. 3. Faire signer \textbf{Fiche d'intervention} (Date, Poste, Actions menées, Résultat, Signature élève + Prof). \textbf{Traçabilité = Professionnalisme.}
\end{enumerate}
\textbf{Conclusion :} La méthode structurée (Sauvegarde $\rightarrow$ Espace $\rightarrow$ Démarrage $\rightarrow$ Disque $\rightarrow$ Sécurité $\rightarrow$ MAJ $\rightarrow$ Intégrité $\rightarrow$ Recette) garantit un résultat mesurable et traçable.
\end{exoresolu}

\courssub{Installation d'une application simple}

\begin{experience}[Observation : La trappe à adwares]
    \textbf{Contexte :} Un élève installe « 7-Zip » via le premier lien Google « Télécharger gratuit 7-Zip ». Résultat : Page d'accueil Chrome changée (SearchBaron), extension « PDF Converter » installée, lenteur navigateur.
    \textbf{Analyse :} Le lien menait sur un site tiers (wrapper) qui a repackagé l'installeur officiel avec des offres groupées (adware). L'élève a cliqué « Suivant-Suivant » sans lire.
    \textbf{Leçon :} L'installation sécurisée repose sur 3 piliers : Source officielle + Vérification intégrité (Hash/Signature) + Installation personnalisée (lecture écrans).
\end{experience}

\begin{propriete}[Signature numérique Windows : Garantie d'origine et d'intégrité]
    Avant d'exécuter un installeur (.exe, .msi) : Clic droit $\rightarrow$ Propriétés $\rightarrow$ Onglet \textbf{Signatures numériques}.
    \begin{itemize}
        \item Vérifier : Nom signataire = Éditeur connu (ex: \texttt{VideoLAN}, \texttt{Igor Pavlov}, \texttt{The Document Foundation}, \texttt{Mozilla Corporation}).
        \item Vérifier : « La signature numérique est valide. »
    \end{itemize}
    Cela garantit que le fichier n'a pas été modifié depuis sa signature par l'éditeur. Complémentaire à la vérification par Hash SHA-256.
\end{propriete}

\begin{methode}[Installer une application simple de manière sécurisée et vérifiée]
    \textbf{Objectif :} Installer un logiciel (ex: VLC, 7-Zip, LibreOffice, Firefox) en garantissant l'intégrité du fichier et la propreté de l'installation (pas de bloatware).
    \begin{enumerate}
        \item \textbf{Source officielle uniquement :} Site éditeur (videolan.org, 7-zip.org, libreoffice.org, mozilla.org). \textbf{Interdit :} Sites tiers (01net, Clubic, Softonic, FileHippo) $\rightarrow$ Installateurs « wrapper » avec publicitaires, barres outils, malware.
        \item \textbf{Choix version :} Windows $\rightarrow$ \textbf{64-bit (x64)} si OS 64 bits (standard depuis 2010). Vérifier : Paramètres $\rightarrow$ Système $\rightarrow$ Informations $\rightarrow$ Type système : « Système d'exploitation 64 bits, processeur x64 ».
        \item \textbf{Vérification intégrité (Hash SHA-256) — \emph{Recommandé pro} :}
              \begin{enumerate}
                  \item Sur site éditeur, noter le hash SHA-256 affiché (ex: VLC 3.0.20 win64 : \texttt{a1b2c3...}).
                  \item Télécharger l'installeur (.exe ou .msi) dans \texttt{Téléchargements}.
                  \item PowerShell (Admin) : \texttt{Get-FileHash -Algorithm SHA256 .\textbackslash vlc-3.0.20-win64.exe}.
                  \item \textbf{Comparer caractère par caractère.} Si différent $\rightarrow$ \textbf{Supprimer, ne pas exécuter}, retélécharger (réseau corrompu ou MITM).
              \end{enumerate}
        \item \textbf{Exécution installation :}
              \begin{enumerate}
                  \item Clic droit sur .exe $\rightarrow$ « Exécuter en tant qu'administrateur » (nécessaire pour \texttt{Program Files}, registre HKLM, pilotes).
                  \item \textbf{Lire chaque écran} (Suivant $\neq$ Installer). \textbf{Décocher} : « Installer la barre d'outils X », « Faire de Yahoo mon moteur », « Participer au programme d'amélioration », « Installer Avast/Chrome/Autre ».
                  \item Choisir « Installation personnalisée / Avancée » si dispo $\rightarrow$ Dossier par défaut OK $\rightarrow$ Raccourcis : Bureau (optionnel), Menu Démarrer (oui).
                  \item Terminer. \textbf{Ne pas lancer l'appli tout de suite.}
              \end{enumerate}
        \item \textbf{Post-installation :}
              \begin{itemize}
                  \item Lancer l'appli $\rightarrow$ Vérifier fonctionnement (ouvrir fichier test).
                  \item Vérifier association de fichiers (Paramètres $\rightarrow$ Applications $\rightarrow$ Applications par défaut).
                  \item \textbf{Supprimer l'installeur} de \texttt{Téléchargements} (gain espace, sécurité).
              \end{itemize}
    \end{enumerate}
\end{methode}

\begin{exoresolu}[Installation sécurisée de 7-Zip 24.08 sur poste élève]
Un élève télécharge \texttt{7z2408-x64.exe} depuis un lien « Télécharger gratuit » trouvé sur Google (premier résultat non officiel). Le fichier fait 1,2 Mo (au lieu de $\sim$1,5 Mo officiel). Il lance l'installation en cliquant vite sur « Suivant », « J'accepte », « Installer ». Au premier lancement, 7-Zip fonctionne mais la page d'accueil Chrome a changé pour « SearchBaron » et une extension « PDF Converter » est installée.
\textbf{Questions :}
\begin{enumerate}
    \item Listez 4 erreurs de procédure commises par l'élève.
    \item Détaillez la procédure correcte complète (Source, Hash, Installation, Nettoyage).
    \item Comment supprimer « SearchBaron » et l'extension parasite proprement (Niveau 1) ?
\end{enumerate}
\tcblower
\textbf{Solution :}
\begin{enumerate}
    \item \textbf{Erreurs :}
          \begin{enumerate}
              \item \textbf{Source non officielle :} Lien Google « gratuit » $\rightarrow$ Site tiers (wrapper) $\rightarrow$ Fichier modifié (1,2 Mo vs 1,5 Mo = installateur wrapper + payload).
              \item \textbf{Pas de vérification Hash / Signature :} Impossible de détecter la modification.
              \item \textbf{Installation « Express / Suivant-Suivant » :} A accepté les offres groupées (Bloatware/Adware) cachées dans les cases pré-cochées.
              \item \textbf{Pas de droit admin vérifié :} (Mineur ici, mais si admin, infection système possible).
          \end{enumerate}
    \item \textbf{Procédure correcte :}
          \begin{enumerate}
              \item Navigateur $\rightarrow$ \texttt{https://www.7-zip.org/} (HTTPS, certificat valide).
              \item Télécharger \texttt{7z2408-x64.msi} (préférable .msi pour déploiement silencieux/propre) ou .exe. Noter Hash SHA-256 affiché sur page : \texttt{7b1a...} (exemple).
              \item PowerShell Admin : \texttt{Get-FileHash .\textbackslash 7z2408-x64.msi -Algo SHA256} $\rightarrow$ Comparer. OK.
              \item Clic droit .msi $\rightarrow$ Installer (MSI n'a pas d'écrans « offres », installation silencieuse standard). Ou .exe $\rightarrow$ Exécuter admin $\rightarrow$ Personnalisée $\rightarrow$ Décocher tout superflu.
              \item Test : Clic droit fichier $\rightarrow$ 7-Zip $\rightarrow$ Ouvrir archive. OK.
              \item Supprimer \texttt{7z2408-x64.msi} de \texttt{Téléchargements}.
          \end{enumerate}
    \item \textbf{Nettoyage Adware (Niv 1) :}
          \begin{enumerate}
              \item \textbf{Chrome :} Paramètres $\rightarrow$ Moteur de recherche $\rightarrow$ Gérer $\rightarrow$ Supprimer « SearchBaron ». Démarrage $\rightarrow$ Ouvrir page nouvelle onglet. Extensions $\rightarrow$ Supprimer « PDF Converter » (icône puzzle).
              \item \textbf{Windows :} Paramètres $\rightarrow$ Applications $\rightarrow$ Applications installées $\rightarrow$ Trier par date $\rightarrow$ Désinstaller programmes suspects installés aujourd'hui (ex: « SearchBaron Uninstaller », « PDF Converter Pro », « Driver Updater »).
              \item \textbf{Scan :} Analyse complète Windows Defender + Malwarebytes Free (scan unique).
              \item \textbf{Réinitialisation navigateur (si persiste) :} Chrome $\rightarrow$ Paramètres $\rightarrow$ Réinitialiser les paramètres $\rightarrow$ Restaurer les paramètres par défaut.
          \end{enumerate}
\end{enumerate}
\textbf{Conclusion :} L'installation « Suivant-Suivant » sur fichier non vérifié est la porte d'entrée n°1 des adwares en milieu scolaire. La vérification Hash/Signature + Lecture écrans + Source officielle est la parade.
\end{exoresolu}

\courssub{Mise à jour des applications et démarche de diagnostic}

\begin{experience}[Observation : La mise à jour, c'est de la sécurité avant tout]
    \textbf{Contexte :} WannaCry (2017) a paralysé des hôpitaux UK (NHS) et des universités. Cause : Faille Windows SMBv1 (MS17-010). Correctif dispo 2 mois avant. Postes non mis à jour = victimes.
    \textbf{Réalité 3ème :} Un poste non mis à jour est une porte ouverte. La maintenance Niv 1 inclut la vérification et l'application des correctifs (Windows Update, MAJ apps, définitions AV).
\end{experience}

\begin{methode}[Mettre à jour les applications et mener un diagnostic de premier niveau]
    \textbf{Objectif :} Maintenir le parc logiciel à jour (correctifs sécurité, bogues) et suivre une démarche logique face à un dysfonctionnement.
    \begin{enumerate}
        \item \textbf{Mises à jour Windows (Socle) :} Paramètres $\rightarrow$ Mise à jour et sécurité $\rightarrow$ Rechercher les mises à jour. Installer : Qualité (mensuel), Fonctionnalités (annuel - ex 22H2 $\rightarrow$ 23H2), Pilotes (optionnel mais recommandé GPU/Chipset). \textbf{Redémarrage planifié} hors heures cours.
        \item \textbf{Mises à jour Applications (Deux canaux) :}
              \begin{itemize}
                  \item \textbf{Microsoft Store :} Bibliothèque $\rightarrow$ Obtenir les mises à jour (Apps UWP : Photos, Calculatrice, Netflix, apps éducatives).
                  \item \textbf{Logiciels Win32 (Classiques) :} Menu Aide $\rightarrow$ Vérifier les MAJ (VLC, 7-Zip, LibreOffice, Firefox, Chrome). \textbf{Ou} : \texttt{winget upgrade --all} (Windows Package Manager, PowerShell Admin) $\rightarrow$ Met à jour tout (Store + Win32 + Scoop/Choco si présents). \emph{Outil pro Niv 1+}.
              \end{itemize}
        \item \textbf{Démarche de diagnostic (Méthode scientifique adaptée) :}
              \begin{enumerate}
                  \item \textbf{Recueillir les symptômes :} Message d'erreur exact (copier/coller / capture écran), moment (démarrage, action précise), fréquence, récence (depuis MAJ ? install ? chute ?).
                  \item \textbf{Isoler la cause (Diviser pour régner) :}
                        \begin{itemize}
                            \item \textbf{Matériel ?} Test croisé (souris/clavier/écran sur autre poste). Bips démarrage (codes beep). LED carte mère.
                            \item \textbf{Logiciel / OS ?} Mode sans échec (F8 / Maj+Redémarrer $\rightarrow$ Dépannage $\rightarrow$ Options avancées $\rightarrow$ Paramètres $\rightarrow$ Redémarrer $\rightarrow$ 4 ou F4). \textbf{Si OK en Mode Sans Échec $\rightarrow$ Cause : Driver tiers / Service démarrage / Appli démarrage.}
                            \item \textbf{Réseau ?} \texttt{ping 1.1.1.1}, \texttt{ping google.fr} (DNS), \texttt{ipconfig /all} (IP, Passerelle, DNS).
                        \end{itemize}
                  \item \textbf{Formuler hypothèse} : « Pilote GPU NVIDIA 560.XX cause plantage Outlook » / « Mise à jour KB504XXXX casse impression ».
                  \item \textbf{Tester hypothèse (Action réversible) :} Rollback driver (Gestionnaire périphériques $\rightarrow$ Propriétés $\rightarrow$ Pilote $\rightarrow$ Restaurer). Désinstaller MAJ (Paramètres $\rightarrow$ Historique $\rightarrow$ Désinstaller). Restauration système (Point de restauration antérieur).
                  \item \textbf{Documenter :} Fiche intervention : Date, Poste, Symptôme, Hypothèse, Action, Résultat, Statut (Résolu / Escaladé Niv 2).
              \end{enumerate}
    \end{enumerate}
\end{methode}

\begin{methode}[Commandes réseau diagnostic de base (Invite de commandes Admin)]
    \begin{itemize}
        \item \texttt{ipconfig /all} : Config IP complète (IPv4, Masque, Passerelle, DNS, MAC, Bail DHCP).
        \item \texttt{ping <adresse>} : Test connectivité (ICMP). \texttt{ping -t} (continu), \texttt{ping -n 10} (10 paquets).
        \item \texttt{tracert <adresse>} : Chemin routeurs (où ça coupe).
        \item \texttt{nslookup <domaine>} : Résolution DNS (serveur utilisé, IP retournée).
        \item \texttt{netstat -an} : Connexions actives / ports écoutés (PID).
        \item \texttt{ipconfig /flushdns} : Vider cache DNS local (après changement DNS/serveur).
        \item \texttt{ipconfig /release \&\& ipconfig /renew} : Renouveler bail DHCP.
    \end{itemize}
\end{methode}

\begin{exoresolu}[Diagnostiquer une perte d'accès Internet en salle de TP]
Salle TP 20 postes. Soudain, 5 postes (côté fenêtre) n'ont plus Internet (icône globe barré), les 15 autres fonctionnent. L'enseignant appelle l'élève référent maintenance.
\textbf{Travail :} Appliquer la démarche de diagnostic (Méthode 6) pour identifier la cause et proposer la résolution Niveau 1.
\tcblower
\textbf{Démarche diagnostique pas à pas :}

\begin{enumerate}
    \item \textbf{Symptômes :} 5 postes KO (groupe localisé), 15 OK. Perte Internet (pas forcément réseau local). Moment : Soudain (pas au boot).
    \item \textbf{Isolation (Sur un poste KO) :}
          \begin{itemize}
              \item \texttt{ipconfig /all} $\rightarrow$ \textbf{Observation clé :} Adresse IPv4 : \texttt{169.254.x.x} (APIPA) \textbf{OU} \texttt{192.168.x.x} (IP valide) mais Passerelle/DNS vides ou incorrects ?
              \item \texttt{ping 192.168.1.1} (Passerelle box/routeur) ? \texttt{ping 8.8.8.8} ? \texttt{ping google.fr} ?
          \end{itemize}
    \item \textbf{Scénarios probables \& Tests :}
          \begin{tabularx}{\linewidth}{|l|X|X|}
          \hline
          \rowcolor{popblueL}
          \textbf{Hypothèse} & \textbf{Test / Observation} & \textbf{Action Niv 1} \\
          \hline
          \textbf{Port Switch HS (5 ports contigus)} & Link LED éteinte sur switch + PC. \texttt{ipconfig} $\rightarrow$ « Média déconnecté » ou APIPA. & 1. Vérifier câble RJ45 (clips cassé ?) $\rightarrow$ Remplacer câble (stock). 2. Changer port switch (port libre connu bon). 3. Si OK $\rightarrow$ Port switch HS $\rightarrow$ Étiquette switch « Port 12-16 HS », signaler admin réseau (Niv 2). \\
          \hline
          \textbf{Serveur DHCP saturé / Panne} & \texttt{ipconfig} $\rightarrow$ APIPA (169.254.x.x) sur les 5 KO. Les 15 OK ont bail valide. \texttt{ipconfig /renew} $\rightarrow$ « Impossible de contacter serveur DHCP ». & 1. Redémarrer service DHCP (si serveur Windows accessible Niv 2). 2. \textbf{Niv 1 :} Config IP statique temporaire (IP libre, Masque, Passerelle, DNS 1.1.1.1) sur les 5 postes pour dépanner séance. 3. Signaler DHCP. \\
          \hline
          \textbf{VLAN / Port Switch mal configuré} & IP OK (DHCP répond), Passerelle OK, \texttt{ping 8.8.8.8} OK, \texttt{ping google.fr} KO (DNS). Ou \texttt{ping passerelle} KO. & Vérifier config DNS (1.1.1.1 / 8.8.8.8). Si VLAN isolé (pas accès passerelle) $\rightarrow$ Niv 2 (config switch). \\
          \hline
          \textbf{Coupure fibre/Box (Général)} & Les 15 OK ? $\rightarrow$ Non, si 15 OK $\rightarrow$ Exclu. \\
          \hline
          \end{tabularx}
    \item \textbf{Résolution probable (Câble/Port) :} L'élève teste : Remplace câble $\rightarrow$ Link LED allumée $\rightarrow$ \texttt{ipconfig /renew} $\rightarrow$ IP valide $\rightarrow$ \texttt{ping google.fr} OK. \textbf{Résolu Niv 1.}
    \item \textbf{Documentation :} Fiche intervention : « Salle TP2, Postes 12-16 : Câbles RJ45 défectueux (clips cassés) remplacés. Ports switch 12-16 OK après test. Rétablissement 10h15. Signé Élève X. »
\end{enumerate}
\textbf{Conclusion :} La localisation spatiale (5 postes côte à côte) oriente vers infrastructure câblage/switch. La commande \texttt{ipconfig /all} est le stéthoscope : APIPA = Pas de DHCP (câble/port/switch/DHCP). IP Valide + Pas Internet = DNS/Pare-feu/Passerelle. La méthode structurée évite de réinstaller Windows (erreur classique élève) pour un câble à 2€.
\end{exoresolu}

\begin{methode}[Rédiger une fiche d'intervention de maintenance traçable et exploitable]
    \textbf{Objectif :} Produire un compte-rendu écrit professionnel (papier ou numérique) permettant la traçabilité, la continuité de service et l'analyse de parc (MTBF, coûts).
    \begin{enumerate}
        \item \textbf{En-tête (Identification) :}
              \begin{itemize}
                  \item Établissement / Salle / Numéro de poste (Étiquette inventaire : ex: \texttt{LYC-COU-TP2-PC12}).
                  \item Date / Heure début / Heure fin.
                  \item Nom intervenant (Élève / Technicien), Niveau (1 / 2 / 3).
                  \item Demandeur (Professeur, Élève, Admin).
              \end{itemize}
        \item \textbf{Description du symptôme (Factuel) :} « Écran noir, ventilateurs tournent, bip 1 long 2 courts (AMI) » $\neq$ « Ça marche pas ». Citer messages exacts, codes erreur, comportement.
        \item \textbf{Diagnostic / Hypothèse :} Résultat des tests (Mode sans échec OK, \texttt{sfc} erreurs, \texttt{ping} KO, Pâte thermique sèche, Câble coupé).
        \item \textbf{Actions menées (Procédure) :} Liste numérotée, outils utilisés (\texttt{cleanmgr}, \texttt{sfc}, Bombe air, Remplacement câble RJ45, Installation VLC 3.0.20). Préciser si pièces changées (Référence, N° série).
        \item \textbf{Résultat / Validation :} « Poste redémarre normalement, session prof ouverte, Outlook OK, Imprimante OK, Température CPU 42°C au repos. \textbf{Validé par utilisateur (Signature)}. »
        \item \textbf{Statut de clôture :}
              \begin{itemize}
                  \item \cle{Résolu Niv 1} : Intervenant clôture.
                  \item \cle{Escaladé Niv 2} : Transmis à \textit{Nom Technicien} le \textit{Date} à \textit{Heure}. Pièce commandée : \textit{Ref}.
                  \item \cle{En attente} : Pièce / Validation budget / Retour constructeur.
                  \item \cle{Irréparable / Reforme} : Avis reforme n°...
              \end{itemize}
        \item \textbf{Archivage :} Classeur « Maintenance Salle TP2 » + Saisie GMAO (Gestion Maintenance Assistée par Ordinateur) si existant (GLPI, OCS Inventory, ou simple Excel partagé).
    \end{enumerate}
\end{methode}

\begin{aretenir}[Indicateurs clés pour le chef d'établissement]
    Une fiche bien remplie alimente : \textbf{MTBF} (Mean Time Between Failures - Fiabilité), \textbf{MTTR} (Mean Time To Repair - Réactivité), \textbf{Coût maintenance / poste / an}, \textbf{Top 5 pannes récurrentes} (justifie achat matériel qualité / formation utilisateurs).
\end{aretenir}

\begin{exoresolu}[Rédaction fiche intervention : Cas « Poste Bibliothécaire - Impression impossible »]
\textbf{Contexte :} Poste \texttt{BIB-ACC-01} (Windows 10, HP LaserJet Pro M404dn réseau IP fixe 192.168.10.50). Bibliothécaire signale : « Imprime plus depuis ce matin, document reste en file d'attente, icône imprimante avec triangle jaune ». Intervenant : Élève maintenance 3ème, Niv 1. Date : 15/10/2024. Début 09h15.
\textbf{Travail :} Rédiger la fiche d'intervention complète (sections 1 à 6) en simulant la découverte : File d'attente bloquée sur job « Rapport\_Oct.pdf », service Spouleur planté, redémarrage service OK, test page OK.
\tcblower
\textbf{FICHE D'INTERVENTION MAINTENANCE NIVEAU 1}

\begin{center}
\begin{tabularx}{\linewidth}{|l|X|}
\hline
\rowcolor{popblueL}
\textbf{Champ} & \textbf{Contenu} \\
\hline
\textbf{1. Identification} & \textbf{Établissement :} Lycée de la Cité (Exemple). \textbf{Salle :} Bibliothèque (CDI). \textbf{Poste :} BIB-ACC-01 (Inv: 2023-IT-045). \textbf{Date :} 15/10/2024. \textbf{Heure début :} 09h15. \textbf{Heure fin :} 09h35. \textbf{Intervenant :} [Nom Élève], Classe 3ème, Niveau 1. \textbf{Demandeur :} Mme [Nom], Bibliothécaire. \\
\hline
\textbf{2. Symptôme} & « Impossible d'imprimer depuis ce poste depuis ce matin 08h30. Document \texttt{Rapport\_Oct.pdf} bloqué en file d'attente (Statut : \textit{Impression...} puis \textit{Erreur}). Icône imprimante HP M404dn avec triangle jaune \textit{Attention}. Autres postes bibliothèque impriment normalement. » \\
\hline
\textbf{3. Diagnostic} & \begin{enumerate}
    \item Ouverture file d'attente : 3 jobs dont \texttt{Rapport\_Oct.pdf} en erreur « \textit{Erreur - Impression impossible} ».
    \item Suppression jobs $\rightarrow$ Bloqué « \textit{Suppression...} » indéfini.
    \item \texttt{services.msc} $\rightarrow$ Service \textbf{Spouleur d'impression (Spooler)} : Statut « \textit{En cours d'exécution} » mais redémarrage manuel nécessaire.
    \item \texttt{Event Viewer} (Observateur d'événements) $\rightarrow$ Journaux Windows $\rightarrow$ System $\rightarrow$ Source \texttt{PrintService} : Event ID 372 « Le pilote d'imprimante a planté » (Driver HP Universal Printing PCL 6).
\end{enumerate}
\textbf{Hypothèse :} Plantage du service Spouleur suite à job corrompu / incompatibilité pilote PCL6 sur PDF complexe. \\
\hline
\textbf{4. Actions menées} & \begin{enumerate}
    \item Arrêt service Spouleur (\texttt{net stop spooler} ou services.msc).
    \item Suppression fichiers \texttt{C:\textbackslash Windows\textbackslash System32\textbackslash spool\textbackslash PRINTERS\textbackslash *.spl} et \texttt{*.shd} (Nettoyage file d'attente orpheline).
    \item Redémarrage service Spouleur (\texttt{net start spooler}).
    \item Test page de test Windows (Propriétés imprimante $\rightarrow$ Imprimer page de test) $\rightarrow$ \textbf{Succès}.
    \item Impression \texttt{Rapport\_Oct.pdf} (réouvert depuis source) $\rightarrow$ \textbf{Succès}.
    \item Conseil utilisateur : Privilégier impression via « Microsoft Print to PDF » $\rightarrow$ puis impression PDF si plantage récurrent (contournement pilote).
\end{enumerate} \\
\hline
\textbf{5. Résultat / Validation} & Impression fonctionnelle. File d'attente vide. Page de test OK. Document métier imprimé. \textbf{Validé par Mme [Nom] (Signature) à 09h34.} \\
\hline
\textbf{6. Statut} & \textbf{\cle{RÉSOLU NIVEAU 1}}. Aucune pièce remplacée. Aucune escalade. Fichier archivé classeur CDI + Saisie registre maintenance Excel partagé (Ligne 142). \\
\hline
\end{tabularx}
\end{center}

\textbf{Conclusion :} Cette fiche est exploitable immédiatement par un technicien Niv 2 si le problème revient (historique pilote PCL6), par l'administration pour le parc (stats pannes imprimantes), et prouve le professionnalisme de l'élève (traçabilité, vocabulaire technique, validation utilisateur).
\end{exoresolu}

\begin{attention}[Les 5 erreurs qui coûtent des points au BEPC (Maintenance 3ème)]
    \begin{enumerate}
        \item \textbf{Confondre « Défragmenter » et « Optimiser » sur SSD :} Écrire « J'ai défragmenté le SSD pour le rendre plus rapide » $\rightarrow$ \textbf{0 point} sur la question optimisation. \textbf{Attendu :} « L'outil Optimiser lance la commande TRIM sur SSD, la défragmentation est interdite (usure). »
        \item \textbf{Oublier la coupure secteur + vidange (Bouton Power 5s) avant ouverture boîtier :} Dans une procédure de nettoyage interne ou vérification RAM, ne pas mentionner « Débrancher prise murale + Appuyer 5s sur Power » $\rightarrow$ \textbf{Faute grave sécurité (Électrique + ESD)}. Point(s) retirés.
        \item \textbf{Installer un logiciel depuis un site tiers (01net, Clubic) sans vérifier la signature/hash :} Réponse « Je télécharge sur Google et j'installe » $\rightarrow$ \textbf{Non-conformité sécurité / intégrité}. Attendu : « Site officiel, vérification signature numérique ou hash SHA-256, installation personnalisée (décocher offres). »
        \item \textbf{Ne pas sauvegarder les données utilisateur avant maintenance logicielle :} Lancer \texttt{cleanmgr}, \texttt{sfc}, restauration système, ou réinstallation sans avoir copié Bureau/Documents sur support externe $\rightarrow$ \textbf{Faute professionnelle majeure}. Risque perte données irréversible. Point(s) retirés.
        \item \textbf{Rédiger une fiche d'intervention vague / sans validation :} « Réparé l'ordinateur » ou « Nettoyé » sans : Symptôme précis, Actions détaillées (outils, commandes), Résultat mesuré, Signature utilisateur, Statut clôture $\rightarrow$ \textbf{Compétence « Rédiger et justifier une démarche » non validée}. La traçabilité \textbf{est} l'évaluation.
    \end{enumerate}
\end{attention}

\newpage

\coursec{Exercices}

\coursec{Maintenance de premier niveau}

\courssub{Notions de base sur la maintenance informatique}

\begin{definition}[Maintenance informatique]
Ensemble des opérations techniques (matérielles et logicielles) visant à maintenir un système informatique en état de fonctionnement, à prévenir les pannes, à corriger les dysfonctionnements et à faire évoluer les performances.
\end{definition}

\begin{propriete}[Les trois types de maintenance]
\begin{itemize}
    \item \textbf{Maintenance préventive :} Actions planifiées \emph{avant} l'apparition d'une panne pour en réduire la probabilité (nettoyage physique, mises à jour, sauvegardes, vérification des disques).
    \item \textbf{Maintenance corrective :} Interventions \emph{après} la survenue d'une panne pour rétablir le service (remplacement pièce défectueuse, réinstallation OS, suppression virus).
    \item \textbf{Maintenance évolutive :} Modifications pour améliorer les performances ou ajouter des fonctionnalités (ajout RAM, changement SSD, migration OS, installation nouveaux logiciels).
\end{itemize}
\end{propriete}

\begin{propriete}[Niveaux de maintenance (MINESEC 3ème)]
\begin{itemize}
    \item \textbf{Niveau 1 (Utilisateur / Élève formé) :} Opérations sans ouverture du boîtier (ou ouverture simple pour nettoyage à l'air sec / vérification câbles), installation logiciels standards, mises à jour, nettoyage disque, création points de restauration. \textbf{Pas de soudure, pas d'alimentation ouverte, pas de composants nus (CPU, GPU) manipulés.}
    \item \textbf{Niveau 2 (Technicien) :} Remplacement composants (alimentation, carte mère, écran portable), diagnostic approfondi, configuration BIOS/UEFI, récupération données disque défaillant.
    \item \textbf{Niveau 3 (Spécialiste/Constructeur) :} Réparation niveau composant (soudure CMS), micro-programmation (firmware), expertise judiciaire, architecture serveurs complexes.
\end{itemize}
\end{propriete}

\courssub{Règles de sécurité avant toute intervention}

\begin{attention}[Règles d'or de sécurité -- À respecter impérativement]
\begin{enumerate}
    \item \textbf{Coupure secteur :} Éteindre l'OS \emph{correctement} (Arrêter), puis débrancher le cordon d'alimentation \textbf{avant} d'ouvrir le boîtier. Attendre 30~s (décharge condensateurs).
    \item \textbf{Antistatique :} Se mettre à la masse (bracelet antistatique ou toucher une partie métallique nue du boîtier \emph{tant qu'il est branché à la terre mais secteur coupé}) avant de toucher tout composant interne (RAM, carte mère, disque). Travailler sur surface non conductrice (tapis, carton).
    \item \textbf{Bloc d'alimentation (PSU) :} \textbf{Jamais} l'ouvrir. Condensateurs haute tension (jusqu'à 400~V) conservent une charge mortelle des semaines après débranchement. Aucune maintenance niveau 1 à l'intérieur.
    \item \textbf{Protection données :} Avant toute maintenance logicielle majeure (MAJ OS, nettoyage registre, suppression masse) : \textbf{Sauvegarder} les données utilisateur (Documents, Bureau, Favoris) sur support externe ou serveur.
    \item \textbf{Outils adaptés :} Tournevis cruciformes (Phillips \#2), bombe air sec (tenue verticale, pas de secousse), pinceau antistatique, chiffon microfibre, alcool isopropylique (écrans/contacts). Pas d'aspirateur (électricité statique), pas de chiffon pelucheux.
\end{enumerate}
\end{attention}

\courssub{Maintenance matérielle de premier niveau}

\begin{methode}[Nettoyage physique de l'unité centrale (Mensuel / Semestriel)]
\begin{enumerate}
    \item Éteindre, débrancher, appuyer 3~s sur bouton Power (vidange capacité résiduelle).
    \item Ouvrir le panneau latéral (généralement côté vitre/fenêtre ou côté carte mère).
    \item Souffler la poussière avec bombe air sec : ventilateurs (CPU, GPU, boîtier, alimentation \emph{par la grille arrière seulement}), radiateurs, filtres à poussière. \textbf{Bloquer les pales des ventilateurs avec un doigt} pour ne pas les faire tourner trop vite (génération courant inverse / usure roulements).
    \item Brosser délicatement les zones inaccessibles (pinceau antistatique).
    \item Vérifier visuellement : câbles bien enfichés (ATX 24 broches, CPU 4/8 broches, PCIe GPU, SATA), pas de condensateurs bombés/fuyants, pas de câbles abîmés.
    \item Refermer, rebrancher, démarrer, vérifier bruits anormaux et températures (BIOS ou logiciel type HWMonitor).
\end{enumerate}
\end{methode}

\begin{aretenir}[Périphériques externes]
Nettoyage souris/clavoir (air sec, coton-tige alcool), écran (chiffon microfibre sec ou légèrement humide \emph{produit écran seulement}, jamais direct sur dalle), imprimante (nettoyage têtes via logiciel, bac papier).
\end{aretenir}

\courssub{Maintenance logicielle de premier niveau}

\begin{propriete}[Opérations courantes niveau 1 -- Windows]
\begin{itemize}
    \item \textbf{Mises à jour Windows Update :} Correctifs sécurité (critiques), pilotes, améliorations. Configurer \emph{Heures d'activité} pour éviter redémarrage en cours de travail.
    \item \textbf{Antivirus / Windows Defender :} Mise à jour définitions (quotidienne auto), analyse rapide (hebdo), analyse complète (mensuel).
    \item \textbf{Nettoyage de disque (\texttt{cleanmgr}) :} Fichiers temporaires, corbeille, fichiers journaux, \emph{Nettoyage des fichiers système} (anciennes versions Windows, rapports d'erreur). \textbf{Ne pas cocher} « Fichiers d'installation ESD » si on veut pouvoir « Réinitialiser ce PC ».
    \item \textbf{Défragmentation (UNIQUEMENT HDD mécanique) :} \texttt{dfrgui} ou \texttt{Optimiser les lecteurs}. \textbf{Interdit sur SSD} (usure prématurée, pas de gain performance). Windows gère le \texttt{TRIM} auto sur SSD.
    \item \textbf{Vérification erreurs système de fichiers :} \texttt{chkdsk C: /f /r} (nécessite redémarrage) ou Propriétés $\to$ Outils $\to$ Vérifier.
    \item \textbf{Point de restauration :} Création manuelle avant installation logiciel/driver risqué. Activation protection système sur C: (par défaut souvent désactivé sur petits SSD).
    \item \textbf{Gestion démarrage (Gestionnaire des tâches $\to$ Démarrage) :} Désactiver programmes inutiles au boot (gain temps, RAM).
\end{itemize}
\end{propriete}

\begin{attention}[Erreurs fréquentes élèves]
\begin{itemize}
    \item Défragmenter un SSD $\to$ \textbf{Non}, cela use les cellules mémoire pour rien.
    \item Nettoyer la base de registre avec des « nettoyeurs » tiers (CCleaner, Advanced SystemCare...) $\to$ \textbf{Risqué}, instabilité possible. Windows gère seul.
    \item Supprimer \texttt{Windows.old} manuellement dans l'explorateur $\to$ \textbf{Impossible/Erreurs}. Utiliser \texttt{cleanmgr} $\to$ Nettoyer fichiers système $\to$ « Installations Windows précédentes ».
    \item Oublier de créer un point de restauration avant d'installer un pilote graphique $\to$ Écran noir au redémarrage, pas de retour arrière facile.
\end{itemize}
\end{attention}

\courssub{Installation d'une application simple}

\begin{methode}[Procédure standard d'installation (Windows)]
\begin{enumerate}
    \item \textbf{Source fiable :} Site officiel éditeur, Microsoft Store, dépôt établissement. \textbf{Pas} de sites « telecharger-gratuit.com », « 01net », « clubic » (installateurs « wrapper » avec adwares).
    \item \textbf{Vérification :} Architecture (x64 / x86 / ARM64), version OS (Win 10/11), prérequis (VC++ Redist, .NET, Java).
    \item \textbf{Analyse antivirus :} Clic droit sur le fichier $\to$ « Analyser avec Windows Defender » (ou AV tiers).
    \item \textbf{Exécution :} Double-clic. Si \texttt{.msi} ou \texttt{.exe} non signé $\to$ « Exécuter en tant qu'administrateur » (UAC). Lire \textbf{chaque écran} (refuser offres partenaires, barres d'outils, modification page d'accueil).
    \item \textbf{Choix installation :} « Personnalisée » (Advanced) pour choisir dossier, raccourcis, associations de fichiers. Éviter « Express/Recommandé » qui installe tout par défaut.
    \item \textbf{Validation :} Lancer l'application, vérifier mise à jour interne (Aide $\to$ Vérifier MAJ), tester fonctionnalité de base (ouvrir/enregistrer).
    \item \textbf{Nettoyage :} Supprimer le fichier d'installation téléchargé (gain espace), vider corbeille.
\end{enumerate}
\end{methode}

\begin{aretenir}[Droits administrateur en établissement]
En session élève (utilisateur standard), l'UAC demande identifiants admin. \textbf{L'élève ne les connaît pas.} Il doit appeler le professeur ou l'intendant informatique. Ne jamais tenter de contourner (cmd admin, outils tiers) $\to$ sanction disciplinaire + risque sécurité.
\end{aretenir}

\courssub{Mise à jour des applications et démarche de diagnostic}

\begin{methode}[Démarche de diagnostic de premier niveau (Panen matériel -- « Ne démarre plus »)]
\begin{enumerate}
    \item \textbf{Alimentation secteur :} Tester prise (lampe/téléphone), multiprise (interrupteur, disjoncteur), câble d'alimentation (échange avec câble connu bon).
    \item \textbf{Alimentation PC (PSU) :} Interrupteur arrière sur « I » (marche), tension secteur sélectionnée (230~V Cameroun). LED carte mère allumée ? (si visible).
    \item \textbf{Démarrage :} Appuyer Power. Observer : Ventilateurs tournent ? LED façade allumée/clignotante ? Bips (codes POST) ? Affichage écran (logo BIOS, curseur) ?
    \item \textbf{Écran / Câble vidéo :} Écran allumé (LED) ? Source correcte (HDMI/DP/VGA) ? Câble vidéo bien vissé des deux côtés ? Test avec autre écran / autre câble.
    \item \textbf{Périphériques :} Tout débrancher (USB, imprimante, réseau) sauf clavier/souris/écran/alimentation. Redémarrer.
    \item \textbf{Interne (si autorisé/formation) :} Débrancher secteur, ouvrir, vérifier connecteurs ATX 24 broches, CPU 4/8 broches, RAM (encliquetée), carte graphique (verrou PCIe), disques SATA/alimentation. Refermer, tester.
    \item \textbf{Conclusion :} Identifier cause probable $\to$ Action niveau 1 (rebrancher, nettoyer contact RAM gomme) $\to$ Ou escalade niveau 2 (PSU, Carte mère, CPU, Écran).
\end{enumerate}
\end{methode}

\begin{savaistu}[Contexte camerounais : ENEO et onduleurs]
Au Cameroun, le délestage et les variations de tension (surtension/sous-tension) sont quotidiens. Un \textbf{onduleur (UPS)} n'est pas un luxe mais une nécessité pour tout poste de travail critique (secrétariat, comptabilité, serveur, salle multimédia).
\begin{itemize}
    \item \textbf{Onduleur « Line-Interactive » (AVR) :} Régule la tension sans passer sur batterie (idéal pour micro-coupures et variations). Autonomie 10--20~min = le temps de sauver le travail et d'éteindre proprement.
    \item \textbf{Logiciel de gestion (ex: PowerChute, ViewPower) :} Configure l'extinction automatique de l'OS si courant ne revient pas après $X$ minutes.
    \item \textbf{Prise « Surge only » (para-foudre) :} Pour imprimantes laser, scanners (gros appel courant au démarrage) -- ne pas brancher sur prises « Battery Backup » de l'onduleur (risque surcharge).
\end{itemize}
\end{savaistu}

\begin{exemplebox}[Commandes utiles diagnostic / maintenance (Invite de commandes \texttt{cmd} en Admin)]
\begin{lstlisting}[language=bash]
:: Vérification intégrité fichiers système (répare Windows)
sfc /scannow

:: Réparation image Windows (si sfc échoue)
DISM /Online /Cleanup-Image /RestoreHealth

:: Vérification disque (planifie au redémarrage si C:)
chkdsk C: /f /r

:: Gestion services (ex: Serveur de fichiers)
net stop LanmanServer
net start LanmanServer
:: Ou PowerShell :
Restart-Service LanmanServer -Force

:: Informations système / Batterie (portable)
powercfg /batteryreport
systeminfo

:: Nettoyage disque en ligne de commande (avancé)
cleanmgr /sagerun:1  (après config avec cleanmgr /sageset:1)
\end{lstlisting}
\legende{Commandes d'administration système courantes en maintenance niveau 1/2.}
\end{exemplebox}

% ============================================================
% EXERCICES
% ============================================================

\courssub{Je vérifie mes connaissances}

\begin{exercice}[QCM : Types de maintenance]
Parmi les propositions suivantes, une seule correspond à une opération de \textbf{maintenance préventive}. Laquelle ?
\begin{enumerate}
    \item Remplacer un disque dur défectueux par un neuf.
    \item Nettoyer la poussière à l'intérieur de l'unité centrale avec une bombe à air sec.
    \item Ajouter une barrette de mémoire vive (RAM) pour accélérer l'ordinateur.
    \item Réinstaller Windows après une infection virale.
\end{enumerate}
\end{exercice}

\begin{exercice}[Vrai ou Faux justifié]
Indiquez si les affirmations suivantes sont \textbf{Vraies} ou \textbf{Fausses} et justifiez brièvement votre réponse.
\begin{enumerate}
    \item La maintenance corrective consiste à agir \emph{avant} qu'une panne ne survienne.
    \item Ouvrir le bloc d'alimentation pour le nettoyer à l'intérieur est une opération de maintenance de premier niveau.
    \item Mettre à jour la base de signatures d'un antivirus relève de la maintenance logicielle préventive.
    \item La défragmentation d'un disque SSD (Solid State Drive) est fortement recommandée chaque semaine.
\end{enumerate}
\end{exercice}

\begin{exercice}[Définitions et vocabulaire]
Donnez une définition précise et concise pour chacun des termes suivants :
\begin{enumerate}
    \item Maintenance évolutive.
    \item Pilote (ou \emph{driver}).
    \item Point de restauration système.
    \item Logiciel malveillant (\emph{malware}).
\end{enumerate}
\end{exercice}

\begin{exercice}[Ordre des opérations : Diagnostic de démarrage]
Un ordinateur de la salle multimédia ne démarre plus (écran noir, aucun bip). Rangez les vérifications suivantes dans l'ordre logique d'un diagnostic de premier niveau (de 1 à 5) :
\begin{enumerate}
    \item Vérifier si le ventilateur du processeur tourne à la mise sous tension.
    \item Tester la prise électrique avec une autre lampe ou un téléphone.
    \item Remplacer la carte mère.
    \item Vérifier que le câble d'alimentation est bien enfiché derrière l'unité centrale et sur la multiprise.
    \item Écouter s'il y a des bips (codes sonores) au démarrage.
\end{enumerate}
\end{exercice}

\courssub{Je m'entraîne}

\begin{exercice}[Classement des opérations : Premier niveau ou pas ?]
Voici une liste d'opérations. Pour chacune, cochez si elle relève de la \textbf{maintenance de premier niveau} (réalisable par un élève/utilisateur formé) ou si elle \textbf{dépasse ce niveau} (nécessite un technicien spécialisé). Justifiez en une phrase.
\begin{center}
\begin{tabularx}{\linewidth}{|l|X|X|}
\hline
\textbf{Opération} & \textbf{Premier niveau ? (Oui/Non)} & \textbf{Justification} \\
\hline
Vider la corbeille de Windows & & \\
\hline
Ouvrir le bloc d'alimentation pour souffler la poussière & & \\
\hline
Mettre à jour l'antivirus (signatures et moteur) & & \\
\hline
Défragmenter le disque dur (HDD) & & \\
\hline
Changer la carte mère & & \\
\hline
Installer une imprimante réseau via l'assistant Windows & & \\
\hline
Remplacer la pâte thermique du processeur & & \\
\hline
Créer un point de restauration avant une installation & & \\
\hline
\end{tabularx}
\end{center}
\end{exercice}

\begin{exercice}[Installation d'une application depuis une clé USB]
Un élève possède le fichier d'installation \texttt{LibreOffice\_24.2\_Win\_x64.msi} sur sa clé USB. Il souhaite l'installer sur le poste de travail de la bibliothèque (Windows 10, session utilisateur standard sans droits administrateur).
\begin{enumerate}
    \item Énumérez les \textbf{5 étapes successives} pour installer proprement cette application.
    \item Citez \textbf{deux précautions indispensables} à prendre \emph{avant} de lancer le fichier \texttt{.msi}.
    \item L'installation demande les droits administrateur. Que doit faire l'élève dans le contexte de l'établissement ?
\end{enumerate}
\end{exercice}

\begin{exercice}[Libérer de l'espace disque : Cas du secrétariat]
Le disque dur principal (C:) du secrétariat du collège a une capacité de \SI{250}{Go}. Il ne reste plus que \SI{10}{Go} libres (96\% plein). Le secrétariat ne peut plus enregistrer les bulletins.
Proposez \textbf{trois opérations de maintenance logicielle de premier niveau}, classées par efficacité probable, pour libérer de l'espace rapidement. Pour chacune, expliquez \emph{quoi} faire et \emph{pourquoi} cela libère de l'espace.
\end{exercice}

\begin{exercice}[Mise à jour Windows et contexte camerounais]
Lors d'une mise à jour majeure de Windows (ex: passage 21H2 $\to$ 22H2), l'ordinateur affiche « \emph{Ne pas éteindre votre ordinateur} ». Une coupure de courant ENEO survient au milieu de l'installation.
\begin{enumerate}
    \item Expliquez techniquement \textbf{pourquoi} l'interruption brutale (coupure secteur) pendant cette phase est critique pour le système d'exploitation.
    \item Proposez \textbf{deux mesures de protection matérielle ou organisationnelle} adaptées au contexte camerounais (coupures fréquentes, délestage) pour éviter ce scénario.
\end{enumerate}
\end{exercice}

\courssub{Je me prépare au Probatoire}

\begin{exercice}[Calendrier de maintenance préventive : Salle multimédia]
Vous êtes désigné responsable de la maintenance des \textbf{10 ordinateurs} de la salle multimédia du lycée. L'objectif est d'assurer leur bon fonctionnement sur l'année scolaire.
\begin{enumerate}
    \item Établissez un \textbf{calendrier de maintenance préventive} sous forme de tableau avec trois colonnes : \textbf{Périodicité} (Hebdomadaire / Mensuelle / Semestrielle), \textbf{Opérations matérielles}, \textbf{Opérations logicielles}.
    \item Pour \textbf{chaque ligne} du tableau, justifiez en une phrase l'utilité de l'opération pour la pérennité du parc.
    \item Le proviseur demande un \textbf{rapport écrit} après chaque intervention semestrielle. Listez \textbf{quatre informations obligatoires} qui doivent figurer dans ce rapport.
\end{enumerate}
\end{exercice}

\begin{exercice}[Diagnostic complet : « Mon ordinateur ne démarre plus »]
Un camarade vous appelle au secours : « \emph{Mon PC ne veut plus s'allumer, j'ai un exposé à rendre dans 2 heures !} » Vous vous rendez sur place. L'écran est noir, la LED d'alimentation de l'unité centrale ne s'allume pas.
\begin{enumerate}
    \item Rédigez la \textbf{démarche de diagnostic complète en 6 étapes numérotées}, depuis l'arrivée devant la machine jusqu'à la conclusion. Pour chaque étape, précisez l'action effectuée, l'observation attendue et la décision qui en découle (ex: « Si OK $\to$ étape suivante ; Si KO $\to$ action corrective »).
    \item Au terme de votre diagnostic, vous identifiez que le \textbf{câble d'alimentation interne (connecteur ATX 24 broches) est débranché} sur la carte mère. Cette réparation relève-t-elle du premier niveau ? Justifiez.
    \item Si le problème était une \textbf{carte mère grillée} (odeur de brûlé, condensateurs bombés), quelle serait votre conclusion et votre recommandation auprès de l'administration ?
\end{enumerate}
\end{exercice}

\begin{exercice}[Gestion d'une mise à jour critique et sécurité des données]
Le lycée dispose d'un serveur de fichiers (partage Samba) où les professeurs déposent les sujets d'examen. Une mise à jour de sécurité critique de Windows Server est disponible. L'administrateur réseau (vous) doit la déployer.
\begin{enumerate}
    \item Avant de lancer la mise à jour, citez \textbf{trois actions préparatoires obligatoires} pour garantir l'intégrité des données et la continuité de service.
    \item La mise à jour nécessite un redémarrage. Proposez une \textbf{stratégie de planification} (jour, heure, communication) adaptée à la vie de l'établissement (cours, examens, administration).
    \item Après redémarrage, le service de partage de fichiers ne démarre pas automatiquement. Décrivez la \textbf{procédure de vérification et de remise en service} en ligne de commande (invite de commandes \texttt{cmd} ou PowerShell) pour redémarrer le service \texttt{LanmanServer}.
    \item En cas d'échec définitif de la mise à jour (boucle de redémarrage, erreur \texttt{0x800F0922}), expliquez comment utiliser le \textbf{Point de restauration} ou le \textbf{Mode sans échec} pour revenir à l'état antérieur.
\end{enumerate}
\end{exercice}

\newpage

\coursec{Corrigés des exercices}

\coursec{Corrigés des exercices d'application — Maintenance de premier niveau}

\begin{corrige}[Exercice 1 — QCM : Types de maintenance]
\textbf{Bonne réponse : 2.} Nettoyer la poussière à l'intérieur de l'unité centrale avec une bombe à air sec.

\textbf{Justification :} la maintenance \cle{préventive} consiste à agir \emph{avant} la panne pour en réduire la probabilité. Le nettoyage régulier de la poussière évite la surchauffe des composants (ventilateurs encrassés, radiateurs obstrués), donc prévient les pannes.

\textbf{Analyse des autres propositions :}
\begin{itemize}
    \item \textbf{1.} Remplacer un disque dur défectueux : la panne a déjà eu lieu $\to$ maintenance \textbf{corrective}.
    \item \textbf{3.} Ajouter de la RAM pour accélérer l'ordinateur : amélioration des performances $\to$ maintenance \textbf{évolutive}.
    \item \textbf{4.} Réinstaller Windows après une infection virale : rétablissement du service après un dysfonctionnement $\to$ maintenance \textbf{corrective}.
\end{itemize}

\textbf{Point de vigilance :} le critère de classement est le \emph{moment} et le \emph{but} de l'intervention : avant la panne (préventive), après la panne (corrective), ou pour améliorer (évolutive).
\end{corrige}

\begin{corrige}[Exercice 2 — Vrai ou Faux justifié]
\begin{enumerate}
    \item \textbf{FAUX.} La maintenance corrective intervient \emph{après} la survenue d'une panne, pour rétablir le service. Agir avant la panne est la définition de la maintenance \emph{préventive}.
    \item \textbf{FAUX.} Le bloc d'alimentation contient des condensateurs haute tension (jusqu'à 400~V) qui conservent une charge mortelle longtemps après le débranchement. Il ne faut \textbf{jamais} l'ouvrir ; on ne souffle la poussière que par la grille arrière. C'est une opération interdite au premier niveau.
    \item \textbf{VRAI.} Mettre à jour les signatures de l'antivirus est une opération logicielle effectuée \emph{avant} toute infection pour prévenir les attaques : c'est bien de la maintenance logicielle préventive.
    \item \textbf{FAUX.} La défragmentation est \textbf{interdite sur un SSD} : elle use prématurément les cellules mémoire (nombre d'écritures limité) sans aucun gain de performance, car un SSD n'a pas de tête de lecture mécanique à déplacer. Elle ne concerne que les disques durs mécaniques (HDD).
\end{enumerate}
\end{corrige}

\begin{corrige}[Exercice 3 — Définitions et vocabulaire]
\begin{enumerate}
    \item \textbf{Maintenance évolutive :} ensemble des modifications apportées à un système informatique pour améliorer ses performances ou lui ajouter des fonctionnalités (ajout de RAM, remplacement d'un disque par un SSD, installation d'un nouveau logiciel).
    \item \textbf{Pilote (driver) :} petit logiciel qui permet au système d'exploitation de communiquer avec un matériel (imprimante, carte graphique, clavier) et de le faire fonctionner correctement.
    \item \textbf{Point de restauration système :} photographie (copie de sauvegarde) des fichiers système et des paramètres de Windows à une date donnée, permettant de ramener le système à un état antérieur stable en cas de problème, sans effacer les documents personnels.
    \item \textbf{Logiciel malveillant (malware) :} programme conçu pour nuire à un système informatique ou à son utilisateur : virus, vers, chevaux de Troie, rançongiciels, logiciels espions.
\end{enumerate}
\end{corrige}

\begin{corrige}[Exercice 4 — Ordre des opérations : Diagnostic de démarrage]
Ordre logique (du plus simple et externe vers le plus complexe et interne) :
\begin{enumerate}
    \item[\textbf{1.}] \textbf{Tester la prise électrique} avec une autre lampe ou un téléphone (proposition b) : on commence par vérifier que le courant arrive.
    \item[\textbf{2.}] \textbf{Vérifier que le câble d'alimentation est bien enfiché} derrière l'unité centrale et sur la multiprise (proposition d).
    \item[\textbf{3.}] \textbf{Écouter s'il y a des bips} (codes sonores) au démarrage (proposition e) : les bips renseignent sur l'état du matériel.
    \item[\textbf{4.}] \textbf{Vérifier si le ventilateur du processeur tourne} à la mise sous tension (proposition a) : signe que la carte mère est alimentée.
    \item[\textbf{5.}] \textbf{Remplacer la carte mère} (proposition c) : opération de dernier recours, qui relève du \textbf{niveau 2 (technicien)}, jamais du premier niveau.
\end{enumerate}

\textbf{Point de vigilance :} un bon diagnostic va toujours des causes les plus simples et les plus fréquentes (absence de courant, câble débranché) vers les causes les plus graves et coûteuses (carte mère). On ne remplace jamais un composant avant d'avoir éliminé les causes simples.
\end{corrige}

\begin{corrige}[Exercice 5 — Classement des opérations : Premier niveau ou pas ?]
\begin{center}
\begin{tabularx}{\linewidth}{|X|c|X|}
\hline
\rowcolor{popblueL}
\textbf{Opération} & \textbf{1\textsuperscript{er} niveau ?} & \textbf{Justification} \\
\hline
Vider la corbeille de Windows & \textbf{Oui} & Opération logicielle simple, sans risque, accessible à tout utilisateur. \\
\hline
Ouvrir le bloc d'alimentation pour souffler la poussière & \textbf{Non} & Condensateurs haute tension (400~V) : danger mortel même débranché. Interdit à tout niveau sauf spécialiste. \\
\hline
Mettre à jour l'antivirus & \textbf{Oui} & Maintenance logicielle préventive courante, souvent automatique. \\
\hline
Défragmenter le disque dur (HDD) & \textbf{Oui} & Outil intégré à Windows (\texttt{dfrgui}), sans ouverture du boîtier ; réservé aux HDD. \\
\hline
Changer la carte mère & \textbf{Non} & Manipulation de composants nus, démontage complet : niveau 2 (technicien). \\
\hline
Installer une imprimante réseau via l'assistant Windows & \textbf{Oui} & Installation logicielle guidée, sans ouverture du boîtier. \\
\hline
Remplacer la pâte thermique du processeur & \textbf{Non} & Démontage du ventirad et manipulation du CPU (composant nu) : niveau 2. \\
\hline
Créer un point de restauration & \textbf{Oui} & Opération logicielle de prévention, réalisable par l'utilisateur. \\
\hline
\end{tabularx}
\end{center}

\textbf{Critère général à retenir :} le premier niveau regroupe les opérations \emph{sans ouverture du boîtier} (ou ouverture simple pour nettoyage à l'air sec), les installations logicielles standards et les mises à jour. Toute manipulation de composants nus ou à risque électrique dépasse ce niveau.
\end{corrige}

\begin{corrige}[Exercice 6 — Installation d'une application depuis une clé USB]
\textbf{1. Les 5 étapes successives :}
\begin{enumerate}
    \item \textbf{Vérifier la compatibilité} du fichier : \texttt{Win\_x64} convient à un Windows 10 64 bits (vérifier dans Paramètres $\to$ Système $\to$ Informations système).
    \item \textbf{Analyser le fichier} avec l'antivirus : clic droit sur \texttt{LibreOffice\_24.2\_Win\_x64.msi} $\to$ « Analyser avec Windows Defender ».
    \item \textbf{Lancer l'installation} : double-clic sur le fichier \texttt{.msi}, puis accepter la demande de contrôle de compte utilisateur (UAC).
    \item \textbf{Suivre l'assistant} en lisant chaque écran : choisir l'installation « Personnalisée » si proposée, accepter la licence, choisir le dossier d'installation, puis cliquer sur « Installer ».
    \item \textbf{Valider l'installation} : lancer LibreOffice depuis le menu Démarrer, ouvrir et enregistrer un document test, puis vérifier les mises à jour internes (Aide $\to$ Vérifier les mises à jour).
\end{enumerate}

\textbf{2. Deux précautions avant de lancer le fichier \texttt{.msi} :}
\begin{itemize}
    \item S'assurer que le fichier provient d'une \textbf{source fiable} (site officiel \texttt{libreoffice.org}) et non d'un site de téléchargement douteux qui ajoute des logiciels parasites.
    \item \textbf{Analyser le fichier avec l'antivirus} pour détecter une éventuelle infection (une clé USB circule de machine en machine et transmet facilement les virus).
\end{itemize}

\textbf{3. Droits administrateur :} en session utilisateur standard, l'élève \textbf{ne doit pas tenter de contourner} la demande de mot de passe administrateur. Il doit appeler le professeur d'informatique ou le responsable du parc, qui saisira lui-même ses identifiants. Contourner la sécurité expose l'établissement à des risques et est passible de sanction disciplinaire.
\end{corrige}

\begin{corrige}[Exercice 7 — Libérer de l'espace disque : Cas du secrétariat]
Trois opérations classées par efficacité probable :

\textbf{1. Nettoyage de disque avec \texttt{cleanmgr} (le plus efficace).}
\emph{Quoi faire :} lancer « Nettoyage de disque », sélectionner C:, cliquer sur « Nettoyer les fichiers système », cocher : fichiers temporaires, corbeille, « Installations Windows précédentes » (\texttt{Windows.old}), rapports d'erreur.
\emph{Pourquoi :} les anciennes installations de Windows et les fichiers temporaires occupent souvent plusieurs dizaines de Go. C'est le gain le plus important et le plus rapide, sans risque pour les documents.

\textbf{2. Vider la corbeille et trier les dossiers personnels (Documents, Bureau, Téléchargements).}
\emph{Quoi faire :} vider la corbeille ; déplacer les archives (anciens bulletins, scans, vidéos) vers un disque externe ou le serveur de fichiers du collège ; supprimer les doublons.
\emph{Pourquoi :} les fichiers supprimés restent dans la corbeille et occupent encore de l'espace tant qu'elle n'est pas vidée ; les dossiers personnels accumulent des fichiers lourds inutilisés au quotidien.

\textbf{3. Désinstaller les applications inutilisées.}
\emph{Quoi faire :} Paramètres $\to$ Applications $\to$ trier par taille, désinstaller les logiciels jamais utilisés (jeux, anciens logiciels, versions multiples d'un même programme).
\emph{Pourquoi :} certaines applications occupent plusieurs Go inutilement ; les désinstaller proprement libère l'espace et allège le démarrage.

\textbf{Point de vigilance :} il ne faut \emph{jamais} supprimer \texttt{Windows.old} manuellement dans l'explorateur (erreurs, fichiers protégés) : on passe obligatoirement par \texttt{cleanmgr}. De même, on ne touche pas aux dossiers \texttt{Windows} et \texttt{Program Files}.
\end{corrige}

\begin{corrige}[Exercice 8 — Mise à jour Windows et contexte camerounais]
\textbf{1. Pourquoi la coupure est critique :}
Pendant une mise à jour majeure, Windows \emph{remplace ses propres fichiers système} (noyau, pilotes, registre) alors qu'ils sont en cours d'utilisation. Les fichiers sont écrits de façon séquentielle : une coupure brutale laisse le système dans un \textbf{état incohérent} — certains fichiers sont nouveaux, d'autres anciens, certains à moitié écrits. Conséquences possibles : impossibilité de démarrer (écran noir, boucle de redémarrage), corruption du registre, perte de données personnelles en cours d'écriture sur le disque. Le système de fichiers lui-même (NTFS) peut être endommagé si l'écriture est interrompue.

\textbf{2. Deux mesures de protection adaptées au contexte camerounais :}
\begin{itemize}
    \item \textbf{Équiper chaque poste critique d'un onduleur (UPS) de type Line-Interactive :} en cas de délestage ENEO, la batterie prend le relais instantanément et offre 10 à 20 minutes d'autonomie, suffisantes pour terminer la mise à jour ou éteindre proprement. L'onduleur régule aussi les surtensions et sous-tensions fréquentes.
    \item \textbf{Planifier les mises à jour majeures hors période de délestage connu et hors heures de cours} (par exemple le samedi matin ou en soirée), en configurant les « Heures d'activité » de Windows Update pour interdire tout redémarrage automatique pendant les heures de travail, et en vérifiant la météo/les annonces de délestage avant de lancer l'opération.
\end{itemize}

\textbf{Point de vigilance :} une imprimante laser ne se branche jamais sur les prises « Battery Backup » de l'onduleur (fort appel de courant au démarrage $\to$ surcharge), mais sur les prises « Surge only » (protection contre la foudre seulement).
\end{corrige}

\begin{corrige}[Exercice 9 — Calendrier de maintenance préventive : Salle multimédia]
\textbf{1. Calendrier de maintenance préventive :}

\begin{center}
\begin{tabularx}{\linewidth}{|l|X|X|}
\hline
\rowcolor{popblueL}
\textbf{Périodicité} & \textbf{Opérations matérielles} & \textbf{Opérations logicielles} \\
\hline
\textbf{Hebdomadaire} & Dépoussiérage extérieur (écrans, claviers, souris) ; vérification des câbles et des onduleurs. & Analyse antivirus rapide ; vérification de l'espace disque libre ; vidage des corbeilles. \\
\hline
\textbf{Mensuelle} & Nettoyage intérieur à la bombe à air sec (poste éteint et débranché) ; contrôle visuel des ventilateurs. & Mises à jour Windows et antivirus ; nettoyage de disque (\texttt{cleanmgr}) ; défragmentation des HDD ; création d'un point de restauration. \\
\hline
\textbf{Semestrielle} & Nettoyage complet et approfondi des 10 postes ; vérification des connecteurs internes (RAM, ATX) ; test des onduleurs et des batteries. & Vérification des disques (\texttt{chkdsk}) ; désinstallation des logiciels inutiles ; sauvegarde complète des données ; bilan et rapport écrit au proviseur. \\
\hline
\end{tabularx}
\end{center}

\textbf{2. Justifications :}
\begin{itemize}
    \item \textbf{Hebdomadaire :} des petites actions fréquentes (poussière de surface, antivirus à jour, espace disque surveillé) évitent l'accumulation des problèmes et détectent tôt les anomalies.
    \item \textbf{Mensuelle :} la poussière intérieure provoque la surchauffe (pannes, durée de vie réduite) ; les mises à jour corrigent les failles de sécurité exploitées par les virus.
    \item \textbf{Semestrielle :} un contrôle approfondi et une sauvegarde complète garantissent la pérennité du parc et la protection des données sur le long terme (année scolaire).
\end{itemize}

\textbf{3. Quatre informations obligatoires du rapport semestriel :}
\begin{enumerate}
    \item La \textbf{date} de l'intervention et le \textbf{nom} de l'intervenant.
    \item La \textbf{liste des postes concernés} (numéros d'inventaire des 10 ordinateurs).
    \item Les \textbf{opérations réalisées} et les \textbf{anomalies constatées} (pannes, pièces à remplacer).
    \item Les \textbf{recommandations} (achats nécessaires : bombe à air, barrettes RAM, onduleurs) et la date de la prochaine intervention.
\end{enumerate}

\textbf{Barème indicatif (sur 10) :} tableau complet et cohérent 4 pts ; justifications pertinentes 3 pts ; rapport (4 informations) 2 pts ; présentation 1 pt.
\end{corrige}

\begin{corrige}[Exercice 10 — Diagnostic complet : « Mon ordinateur ne démarre plus »]
\textbf{1. Démarche de diagnostic en 6 étapes :}

\begin{enumerate}
    \item \textbf{Tester la prise électrique.} \emph{Action :} brancher une lampe ou un téléphone sur la prise / la multiprise. \emph{Observation attendue :} la lampe s'allume. \emph{Décision :} Si OK $\to$ étape 2 ; Si KO $\to$ le problème vient de l'installation électrique (disjoncteur, délestage ENEO), changer de prise.
    \item \textbf{Vérifier la chaîne d'alimentation externe.} \emph{Action :} contrôler l'interrupteur de la multiprise, le câble d'alimentation (bien enfiché des deux côtés), l'interrupteur arrière du bloc d'alimentation sur « I ». \emph{Observation :} tout est enfoncé et en position marche. \emph{Décision :} Si OK $\to$ étape 3 ; Si KO $\to$ rebrancher et retester.
    \item \textbf{Observer les signes de vie à la mise sous tension.} \emph{Action :} appuyer sur le bouton Power. \emph{Observation :} LED de façade ? ventilateurs ? bips ? \emph{Décision :} Si des bips ou des LED $\to$ noter le code et orienter le diagnostic ; Si rien du tout (notre cas) $\to$ étape 4.
    \item \textbf{Éliminer les périphériques.} \emph{Action :} débrancher tous les périphériques USB, imprimante, réseau, ne garder que écran, clavier, alimentation. \emph{Observation :} le PC démarre-t-il ? \emph{Décision :} Si OK $\to$ un périphérique était en cause, les rebrancher un par un ; Si KO $\to$ étape 5.
    \item \textbf{Contrôle interne (poste éteint et débranché du secteur).} \emph{Action :} ouvrir le panneau latéral, se mettre à la masse, vérifier visuellement : connecteur ATX 24 broches, connecteur CPU 4/8 broches, RAM bien encliquetée. \emph{Observation :} ici, le connecteur ATX 24 broches est débranché. \emph{Décision :} Si un connecteur est desserré $\to$ le rebrancher fermement, refermer, tester ; Si tout est bien branché $\to$ étape 6.
    \item \textbf{Conclusion et escalade.} \emph{Action :} résumer la cause identifiée. Si la panne persiste après toutes les vérifications, la cause est probablement matérielle profonde (bloc d'alimentation ou carte mère hors service) $\to$ \textbf{escalade vers un technicien de niveau 2}, avec un compte rendu précis des tests déjà effectués.
\end{enumerate}

\textbf{2. Rebrancher le connecteur ATX 24 broches : premier niveau ?}
\textbf{Oui}, cette réparation relève du premier niveau, à condition de respecter les règles de sécurité : poste éteint, cordon débranché, attente de 30 secondes, mise à la masse avant de toucher l'intérieur. Il s'agit d'une simple vérification/enfichage de câble, sans manipulation de composant nu ni soudure : c'est exactement le type d'intervention prévu au niveau 1.

\textbf{3. Si la carte mère est grillée :}
\emph{Conclusion :} la panne dépasse le premier niveau : le remplacement d'une carte mère est une opération de niveau 2 (technicien), et une carte mère avec condensateurs bombés et odeur de brûlé n'est pas réparable à notre niveau.
\emph{Recommandation à l'administration :} rédiger un rapport écrit décrivant les symptômes et les tests effectués, mettre le poste hors service (étiquette « EN PANNE »), proposer soit l'intervention d'un technicien agréé pour remplacement, soit — si le poste est ancien — son remplacement, en récupérant si possible le disque dur pour sauvegarder les données (ici, l'exposé du camarade !).

\textbf{Barème indicatif (sur 10) :} démarche ordonnée en 6 étapes avec observations et décisions 5 pts ; justification niveau 1 avec règles de sécurité 2 pts ; conclusion et recommandation 3 pts.
\textbf{Points de vigilance :} exiger l'ordre « externe $\to$ interne », la mention explicite des règles de sécurité (débrancher avant d'ouvrir), et la distinction nette niveau 1 / niveau 2.
\end{corrige}

\begin{corrige}[Exercice 11 — Gestion d'une mise à jour critique et sécurité des données]
\textbf{1. Trois actions préparatoires obligatoires :}
\begin{enumerate}
    \item \textbf{Sauvegarder les données} du serveur (sujets d'examen, documents des professeurs) sur un support externe ou un autre serveur : en cas d'échec de la mise à jour, aucune donnée n'est perdue.
    \item \textbf{Créer un point de restauration système} (ou une image système) pour pouvoir revenir à l'état antérieur stable en cas de problème.
    \item \textbf{Vérifier l'alimentation électrique} : s'assurer que le serveur est branché sur un onduleur chargé et fonctionnel, afin qu'une coupure ENEO n'interrompe pas la mise à jour ; vérifier aussi que l'espace disque libre est suffisant.
\end{enumerate}

\textbf{2. Stratégie de planification :}
Lancer la mise à jour \textbf{un samedi après-midi ou en soirée (par exemple à 18~h)}, hors heures de cours, hors période d'examens et hors saisie des notes par l'administration. \textbf{Communication :} informer par écrit (note de service, affichage, message au groupe des enseignants) au moins 48~h à l'avance que le serveur sera indisponible, en précisant la date, l'heure et la durée estimée, pour que chacun sauvegarde son travail et ne soit pas surpris.

\textbf{3. Remise en service du partage de fichiers :}
Ouvrir une invite de commandes \textbf{en tant qu'administrateur}, puis :
\begin{enumerate}
    \item Vérifier l'état du service :
\begin{lstlisting}[language=bash]
sc query LanmanServer
\end{lstlisting}
    \item Si le service est arrêté, le démarrer :
\begin{lstlisting}[language=bash]
net start LanmanServer
\end{lstlisting}
    \item Vérifier à nouveau son état (\texttt{sc query LanmanServer} doit afficher \texttt{RUNNING}), puis tester concrètement : depuis un poste client, ouvrir \texttt{\textbackslash\textbackslash SERVEUR\textbackslash Examens} et lire un fichier.
    \item Si le service ne démarre pas, consulter l'observateur d'événements (\texttt{eventvwr}) pour identifier l'erreur.
\end{enumerate}

\textbf{4. Retour à l'état antérieur en cas d'échec :}
\begin{itemize}
    \item \textbf{Point de restauration :} démarrer sur les options de récupération (touche F8 ou clé d'installation Windows $\to$ « Réparer l'ordinateur » $\to$ Dépannage $\to$ Options avancées $\to$ « Restauration du système »), choisir le point créé \emph{avant} la mise à jour, laisser la restauration s'effectuer : le système revient à son état stable antérieur, sans toucher aux fichiers personnels.
    \item \textbf{Mode sans échec :} si Windows redémarre en boucle, forcer le démarrage en mode sans échec (Windows ne charge alors que les pilotes essentiels). Depuis ce mode, on peut désinstaller la mise à jour défectueuse (Paramètres $\to$ Windows Update $\to$ Historique $\to$ Désinstaller les mises à jour) ou lancer la restauration du système, puis redémarrer normalement.
\end{itemize}

\textbf{Barème indicatif (sur 10) :} préparation (3 actions justifiées) 3 pts ; planification et communication 2 pts ; procédure en ligne de commande correcte 3 pts ; restauration / mode sans échec 2 pts.
\textbf{Points de vigilance :} la sauvegarde doit toujours être citée \emph{en premier} ; la commande doit être exacte (\texttt{net start LanmanServer}) ; rappeler que la restauration du système ne supprime pas les documents personnels.
\end{corrige}

\newpage

\coursec{Fiche bilan}

\begin{popfigure}
\centering
\begin{tikzpicture}[
    mindmap,
    grow cyclic,
    every node/.style={concept, execute at begin node=\hskip0pt},
    concept color=popblue,
    level 1/.append style={level distance=4.5cm, sibling angle=360/6},
    level 2/.append style={level distance=3cm, sibling angle=45},
    text=white,
    font=\small\bfseries
]
\node[concept, scale=1.2, align=center]{Maintenance\\ de premier niveau}
    child[concept color=poporange] { node[concept, align=center]{Notions de base\\ (curative, préventive,\\ évolutive)} }
    child[concept color=popgreen] { node[concept, align=center]{Règles de sécurité\\ (électrique, ESD,\\ données)} }
    child[concept color=poppurple] { node[concept, align=center]{Matérielle\\ (nettoyage, câbles,\\ périphériques)} }
    child[concept color=poppink] { node[concept, align=center]{Logicielle\\ (antivirus, nettoyage,\\ désinstallation)} }
    child[concept color=popgold] { node[concept, align=center]{Installation\\ (sources sûres,\\ assistant, droits)} }
    child[concept color=popdark] { node[concept, align=center]{Mise à jour\\ \& Diagnostic\\ (auto, manuel,\\ observer/tester)} };
\end{tikzpicture}
\legende{Carte mentale : Maintenance de premier niveau — 3ème}
\end{popfigure}

\begin{bilan}

\courssub{Définitions clés}
\begin{itemize}
    \item \cle{Maintenance informatique} : Ensemble des actions visant à maintenir ou rétablir un bien (matériel/logiciel) dans un état spécifié pour qu'il puisse assurer un service déterminé.
    \item \cle{Maintenance préventive} : Actions planifiées pour réduire la probabilité de panne (nettoyage, mises à jour, sauvegardes).
    \item \cle{Maintenance curative} : Actions de réparation après une panne (remplacement pièce, réinstallation).
    \item \cle{Maintenance évolutive} : Amélioration des performances ou adaptation aux nouveaux besoins (ajout RAM, maj OS).
    \item \cle{ESD (ElectroStatic Discharge)} : Décharge électrostatique ; risque majeur pour les composants sensibles (RAM, CPU, CM).
\end{itemize}

\courssub{Règles de sécurité (Ordre impératif)}
\begin{enumerate}
    \item \textbf{Électrique} : Débrancher l'alimentation secteur (prise murale) \emph{avant} d'ouvrir le boîtier. Attendre 30 s (condensateurs).
    \item \textbf{Antistatique (ESD)} : Porter un bracelet antistatique relié à la masse (boîtier métallique) ou se décharger en touchant le boîtier nu avant toute manipulation.
    \item \textbf{Données} : Sauvegarder les fichiers importants \emph{avant} toute intervention logicielle majeure (formatage, maj OS, nettoyage registre).
    \item \textbf{Environnement} : Plan de travail propre, éclairé, sec, hors de portée des liquides. Outils adaptés (tournevis cruciforme/plat, bombe à air sec, chiffon microfibre, pinceau antistatique).
\end{enumerate}

\courssub{Maintenance matérielle de premier niveau (Check-list rapide)}
\begin{center}
\begin{tabularx}{\linewidth}{|l|X|}
\hline
\rowcolor{popblueL}
\textbf{Composant} & \textbf{Action de premier niveau} \\
\hline
Boîtier / Ventilateurs & Souffler la poussière (bombe air sec, brosse souple). Vérifier rotation ventilos (CPU, alim, boîtier). \\
\hline
Écran & Nettoyer avec chiffon microfibre sec ou légèrement humide (produit écran). Pas de spray direct. \\
\hline
Clavier / Souris & Retourner + tapoter + souffler. Nettoyer touches (chiffon + alcool isopropylique \textless 70\%). Capteur souris propre. \\
\hline
Câblage & Vérifier branchement ferme (alim, vidéo, USB, RJ45). Remplacer câbles abîmés (coupures, broches tordues). \\
\hline
Périphériques externes & Imprimante : nettoyage têtes, bac papier. Clé USB/Disque dur : éjection logicielle avant retrait physique. \\
\hline
\end{tabularx}
\end{center}

\courssub{Maintenance logicielle de premier niveau}
\begin{itemize}
    \item \textbf{Antivirus / Antimalware} : Mise à jour des signatures \rightarrow Analyse complète (quotidienne/hebdo). Quarantaine/Suppression menaces.
    \item \textbf{Nettoyage disque} : Supprimer fichiers temporaires (\texttt{\%temp\%}, \texttt{cleanmgr}), vider corbeille, désinstaller logiciels inutiles (Panneau de configuration \textgreater Programmes).
    \item \textbf{Mises à jour système} : Windows Update / \texttt{apt update \&\& apt upgrade} (Linux). Redémarrer si requis.
    \item \textbf{Outils intégrés} : \texttt{chkdsk /f} (erreur disque), \texttt{sfc /scannow} (fichiers système Windows), \texttt{msconfig} / Gestionnaire des tâches (démarrage).
\end{itemize}

\courssub{Méthode express : Installation d'une application simple}
\begin{methode}[Installation standard (Windows \texttt{.exe} / \texttt{.msi})]
\begin{enumerate}
    \item Télécharger \textbf{uniquement} depuis le site officiel de l'éditeur ou magasin de confiance (Microsoft Store).
    \item Vérifier l'extension (\texttt{.exe}, \texttt{.msi}) et la signature numérique (Propriétés \textgreater Signatures numériques).
    \item Exécuter en \textbf{Administrateur} (clic droit \textgreater Exécuter en tant qu'administrateur) si demande UAC.
    \item Suivre l'assistant : Accepter licence (EULA), choisir dossier (défaut \texttt{Program Files}), refuser logiciels parasites (barres d'outils, antivirus d'essai) \rightarrow \textbf{Installation personnalisée}.
    \item Lancer l'application, vérifier fonctionnement, créer raccourci si besoin.
\end{enumerate}
\end{methode}

\courssub{Méthode express : Mise à jour \& Diagnostic de base}
\begin{methode}[Démarche de diagnostic (Approche systématique)]
\begin{enumerate}
    \item \textbf{Observer} : Message d'erreur exact, code, comportement (bip, écran bleu, lenteur, absence réseau). Noter.
    \item \textbf{Identifier le périmètre} : Matériel ? Logiciel ? Réseau ? Un poste ? Tout le réseau ? Depuis quand ? Changement récent ?
    \item \textbf{Tester hypothèses simples} : Redémarrer (premier remède). Vérifier câbles/alimentation. Test avec un autre périphérique/câble/prise.
    \item \textbf{Consulter journaux} : Observateur d'événements (\texttt{eventvwr.msc}) \textgreater Journaux Windows \textgreater Système/Application. Filtrer \textbf{Erreur}/\textbf{Critique}.
    \item \textbf{Appliquer correctif} : Maj pilote, réinstallation logiciel, restauration système (point de restauration), nettoyage malware.
    \item \textbf{Valider \& Documenter} : Confirmer retour à la normale. Noter cause/solution pour base de connaissances.
\end{enumerate}
\end{methode}

\end{bilan}

\begin{aretenir}[Checklist avant le BEPC]
    $\square$ Je sais définir maintenance préventive, curative, évolutive et donner un exemple de chaque.
    $\square$ Je cite les 4 règles de sécurité impératives avant d'ouvrir un boîtier (électrique, ESD, données, environnement).
    $\square$ Je sais nettoyer correctement : ventilateurs (air sec), écran (microfibre), clavier (retourner/souffler), contacts RAM (gommé).
    $\square$ Je sais vérifier l'intégrité des câbles (alimentation, vidéo, RJ45) et l'éjection logicielle des médias amovibles.
    $\square$ Je sais effectuer une analyse antivirus complète et mettre à jour les signatures.
    $\square$ Je sais libérer de l'espace disque (fichiers temporaires, corbeille, désinstallation logiciels inutiles).
    $\square$ Je sais lancer les mises à jour système (Windows Update / \texttt{apt upgrade}) et pourquoi redémarrer est parfois nécessaire.
    $\square$ Je sais installer un logiciel \texttt{.exe}/\texttt{.msi} depuis une source officielle en refusant les logiciels parasites (installation personnalisée).
    $\square$ J'applique la démarche de diagnostic en 6 étapes : Observer $\rightarrow$ Identifier $\rightarrow$ Tester $\rightarrow$ Consulter journaux $\rightarrow$ Corriger $\rightarrow$ Valider/Documenter.
    $\square$ Je connais les outils de base Windows : \texttt{cleanmgr}, \texttt{chkdsk}, \texttt{sfc /scannow}, \texttt{eventvwr.msc}, Gestionnaire des tâches (démarrage).
    $\square$ Je sais expliquer pourquoi on ne souffle pas dans les connecteurs (humidité) et pourquoi on ne vaporise pas directement sur l'écran.
\end{aretenir}

\findecours

\end{document}
